\documentclass[10pt]{article}
\usepackage[margin=0.85in]{geometry}
\usepackage[T1]{fontenc}
\usepackage{amsmath,amssymb,amsthm}
\usepackage{multicol}
\usepackage{hyperref}
\newtheorem{theorem}{Theorem}
\newtheorem{lemma}{Lemma}
\newtheorem{proposition}{Proposition}
\newtheorem{corollary}{Corollary}
\newcommand{\good}{\mathcal{G}}
\newcommand{\cubes}{\mathcal{C}}
\title{Cube-Sum Permutations: Global Involutions, Reflection Entropy, and Algebraic Families}
\author{Qichao Wang\\{\small Hebei University of Technology, Tianjin, China}\\[0.4em]Daoyu Dong\\{\small University of Electronic Science and Technology of China, Chengdu, China}\\[0.4em]{\small Independent mathematical and computational review requested}}
\date{October 18, 2026}
\begin{document}
\maketitle
\begin{abstract}
Let $P(n)$ count permutations $p$ of $[1,n]$
for which every $i+p(i)$ is an integer cube.
We prove that every odd $n\ge70001$ admits
a legal involution with exactly one fixed
point, while every even $n\ge230000$
admits one with no fixed points.
These parity-minimal involutions have
uniform cube-root entropy:
\[
\liminf_{n\to\infty}
\frac{\log J_{\min}(n)}{(2n)^{1/3}}
\ge\frac{\log2}{4}.
\]
The proof uses two short explicit
seed permutations, finite parity-specific
odd-cube reflection closures, and a
uniform gap induction, with no
bipartite matching search.
We also establish exponential numbers
of parity-minimal involutions along
explicit infinite subsequences
using taxicab equal-cube identities.
We develop further global, threshold-free structure:
the cube-sum graph admits a partial matching
of size $n-O(n^{2/3})$ for every $n$,
and all Hall obstructions can be restricted
to independent sets spanning the two
boundary regions. The undirected cycle-cover
expansion gives a parity determinant over
$\mathbb F_2$, while a single valid matching
factors exact counting through the strongly
connected components of its alternation graph.
An explicit four-reflection permutation
on $171$ elements yields the affine
renormalization
\[
P(343n+171)\ge P(n)^{343},
\]
and an entropy monotonicity theorem along
its multiplicative orbits.
Separately, an elementary weighted identity
of cubes gives a fully algebraic
exponential-subsequence construction
\[
P(681707u^3)\ge2^{148877u^3-1}
\qquad(u\ge1).
\]
For the finite A006063 existence theorem, a
unimodular block-order matrix gives a
single symbolic determinant formula for
all $57$ residual cubic-word/output-order
shapes. Six additional explicit algebraic
families supply $25$ further base lengths.
Every one of the $57$ shapes also admits an
infinite cubic dilation family.
The earlier endpoint-multiset and $191$
augmenting-word certificates are retained
as independent redundant checks;
Hall and
alternating-cycle certificates classify
the related multiplicity sequence A096680.
This last classification remains
computer-assisted, unlike the global
structural and algebraic constructions.
All certificates and independent exact
checkers accompany the manuscript;
literature priority and external
mathematical review remain open.
\end{abstract}

\section{The reflection principle}
Write $\cubes=\{k^3:k=2,3,\dots\}$ and let $P(n)$ be the
number of cube-sum permutations of $\{1,\dots,n\}$.
By convention the empty permutation is permitted only in the
general extension lemma; our classification concerns $n\ge1$.
Let $\good=\{n\ge1:P(n)>0\}$.

\begin{lemma}[Reflection and counting]\label{lem:reflect}
If $0\le m<n$ and $c=m+n+1$ is a cube, every valid
permutation $p$ on $\{1,\dots,m\}$ extends to a valid permutation
$q$ on $\{1,\dots,n\}$ via
\[
q_i=\begin{cases}
p_i,&i\le m,\\ c-i,&m<i\le n.
\end{cases}
\]
For different admissible cubes $c$, the images of these
extensions are disjoint. In particular,
\begin{equation}\label{eq:count}
P(n)\ge\!\!\!\!\sum_{\substack{c\in\cubes\\n+1\le c\le 2n}}
P(c-n-1),
\end{equation}
with the convention $P(0)=1$.
\end{lemma}
\begin{proof}
Because $c=m+n+1$, reflection $i\mapsto c-i$
permutes precisely the new interval $\{m+1,\dots,n\}$,
and $i+q_i=c$ there. Old indices retain their cube sums.
The embedding is injective for each $c$ and two different
$c$ yield different last entries $q_n=c-n$. The claimed
sum counts disjoint families.
\end{proof}

\begin{lemma}[The cube-gap bound]\label{lem:gap}
For each $n\ge365$ some cube $c$ satisfies
$n+174\le c\le2n$.
For each $n\ge556$ some cube $c$ satisfies
$n+176\le c\le2n$.
\end{lemma}
\begin{proof}
For $365\le n\le555$, take $c=729$.
For $n\ge556$, let $k$ be the least integer with
$k^3\ge n+176$; then $k\ge10$ and
$n\ge(k-1)^3-175$. Thus
\[
2n-k^3\ge k^3-6k^2+6k-352\ge0.
\]
The final polynomial equals $108$ at $k=10$ and
has strictly positive successive differences
$3k^2-9k+1$ for $k\ge10$.
Hence the second assertion holds and implies the first
for $n\ge556$.
\end{proof}

\section{Global graph structure, factorization, and parity}
\label{sec:graph}

The cube-sum condition is symmetric in its two indices.
This observation gives a conceptual interpretation of
existence, multiplicity, exact counting, and parity that
does not refer to any numerical threshold.

Let $G_n$ be the undirected graph on $[1,n]$,
with an edge $\{i,j\}$ whenever $i\ne j$
and $i+j$ is a perfect cube. At a vertex $i$,
allow a loop exactly when $2i$ is a cube.
A \emph{spanning cycle-cover configuration}
is a decomposition of its vertices into
disjoint allowed loops, single edges
(dimer components), and unoriented simple
cycles of length at least three.

\begin{lemma}[Arithmetic location of fixed points]
\label{lem:fixedpoints}
In any cube-sum permutation, a fixed point
$p(i)=i$ can occur only at
\[
\boxed{i=4r^3\qquad(r\ge1).}
\]
Consequently a valid involution of odd
length must have at least one fixed point
at such an index. In particular, a
cube-sum permutation of odd length
$n<4$ cannot be an involution.
\end{lemma}
\begin{proof}
At a fixed point $2i=k^3$.
The cube on the right is even, so
$k=2r$, whence $i=4r^3$.
An involution partitions its vertices
into fixed points and transpositions;
when $n$ is odd at least one fixed
point is necessary.
\end{proof}

\begin{theorem}[Independent-set Hall criterion]
\label{thm:independent-hall}
A cube-sum permutation exists on $[1,n]$
if and only if every vertex set $I$ of
$G_n$ containing no internal edge
\emph{or loop} satisfies
\begin{equation}\label{eq:independent-hall}
\boxed{|N_{G_n}(I)|\ge |I|,}
\end{equation}
where $N_{G_n}(I)$ includes all vertices
adjacent to some vertex of $I$.
Thus the full family of Hall inequalities
can be reduced to \emph{independent}
vertex sets in the symmetric cube-sum graph.
\end{theorem}
\begin{proof}
The bipartite double cover of $G_n$
has a perfect matching precisely when
Hall's inequality $|N(S)|\ge|S|$
holds for every $S\subseteq[1,n]$.
Its perfect matchings are exactly
the cube-sum permutations.
Hence the proposed condition is necessary.

Conversely, suppose Hall's inequality
fails for some $S$. Put
$I=S\setminus N(S)$.
No two vertices of $I$ can be adjacent,
nor can $I$ contain a loop: otherwise
one of those vertices would lie in
$N(S)$.
Moreover $N(I)\cap S=\varnothing$,
because any edge from $I$ to $S$
would again contradict $I\cap N(S)=
\varnothing$. Thus
\[
|N(I)|\le |N(S)\setminus S|
=|N(S)|-|N(S)\cap S|
<|S|-|N(S)\cap S|=|I|.
\]
So an independent Hall obstruction
exists whenever any Hall obstruction
exists. This proves sufficiency.
\end{proof}

In particular, every negative
certificate supplied for $n<173$ in
the associated JSON is already an
independent vertex set: none of its
members are adjacent to each other
or carry a permitted loop.
The criterion is global and does not
depend on any conjectural cutoff.

\begin{theorem}[Uniform near-perfect partial matchings]
\label{thm:nearperfect}
For every integer $n\ge7$, write
\[
k=\lfloor(n+1)^{1/3}\rfloor,\qquad
g(k)=(k+1)^3-k^3=3k^2+3k+1.
\]
There is an explicitly constructible
cube-sum \emph{partial permutation}
matching at least
\begin{equation}\label{eq:nearperfect}
\boxed{\displaystyle
n-\left\lfloor\frac{g(k)}2\right\rfloor}
\end{equation}
indices to distinct values within $[1,n]$.
Consequently the maximum matching of
the cube-sum bipartite graph has
deficiency $O(n^{2/3})$ for \emph{every}
$n$, including lengths where no
complete cube-sum permutation exists.
\end{theorem}
\begin{proof}
Choose whichever of the consecutive cubes
$c_-=k^3$ and $c_+=(k+1)^3$
is closest to $n+1$, and denote it by $c$.
Their gap is $g(k)$, hence
$\delta=|c-(n+1)|\le\lfloor g(k)/2\rfloor$.
Since $n\ge7$, $c_-\ge8$.
If $c_+>2n$, then
$c_+-(n+1)\ge n$ while
$(n+1)-c_-\le n-7$, so the nearer
cube must be $c_-$. Thus the chosen
$c$ lies in $[2,2n]$.

The interval
\[
J_c=[\max(1,c-n),\ \min(n,c-1)]
\cap\mathbb Z
\]
has precisely $n-\delta$ elements.
It is invariant under $i\mapsto c-i$.
On this interval set $p(i)=c-i$,
giving a bijection $J_c\to J_c$
with the constant cube sum $i+p(i)=c$.
Therefore it defines a partial perfect
matching of size $n-\delta$.
The claimed bound follows.
\end{proof}

\begin{theorem}[Localization of Hall obstructions]
\label{thm:obstruction-localization}
For $n\ge7$ let
\[
K=\lfloor(2n)^{1/3}\rfloor,\qquad
G=3K^2+3K+1.
\]
Every nonempty subset $S\subseteq[1,n]$
with diameter
$\operatorname{diam}(S)=\max S-\min S$
satisfying
\[
\boxed{\operatorname{diam}(S)\le n-G}
\]
automatically obeys Hall's condition
$|N(S)|\ge|S|$ in the symmetric
cube-sum bipartite graph.
Consequently \emph{every} independent
Hall obstruction must satisfy
\[
\boxed{\operatorname{diam}(S)>n-G,}
\]
and must contain vertices from both
the lower and upper boundary regions
of width at most $G$.
\end{theorem}
\begin{proof}
Write $a=\min S$, $b=\max S$.
The candidate integer interval
\[
[b+1,a+n]
\]
has $n-(b-a)\ge G$ consecutive
integers, all between $2$ and $2n$.
Within this range there cannot be
$G$ consecutive noncubes: the
difference between consecutive cubes
is $3j^2+3j+1\le G$ for
$j\le K$, while the initial integers
$2,\dots,7$ also give fewer than
$G$ noncubes.
Consequently some perfect cube $c$
lies in $[b+1,a+n]$.

For every $i\in S$ we then have
$1\le c-i\le n$, and all
$c-i$ are distinct.
Thus $i\mapsto c-i$ is an injection
of $S$ into its neighborhood $N(S)$,
proving Hall's condition.
The contrapositive gives the claimed
diameter restriction. The latter implies
$\min S<G$ and $\max S>n-G+1$,
so a Hall obstruction has to meet
both boundary regions.
\end{proof}

These theorems supply a completely
noncomputational approach to the
existence question. Instead of
checking all subsets, one need only
understand \emph{long-spanning
independent sets} whose vertices
reach both ends of $[1,n]$.
Eliminating those obstructions by a
uniform argument would remove the
remaining finite matching certificates;
the near-perfect bound alone does
not yet yield such an elimination.

\begin{theorem}[Arithmetic codegrees and expansion]
\label{thm:arith-expansion}
Let $H_n$ be the cube-sum bipartite graph
with left and right vertices $[1,n]$.
For distinct left vertices $i\ne j$,
their number of common right neighbors
satisfies
\begin{equation}\label{eq:codegree}
\boxed{\displaystyle
|N(i)\cap N(j)|\le\tau(|i-j|),}
\end{equation}
where $\tau(d)$ is the number of
positive divisors of $d$.
Write
\[
\delta_n=\min_i|N(i)|,\quad
\Delta_n=\max_i|N(i)|,\quad
D_n=\max_{1\le d<n}\tau(d).
\]
For every nonempty $S\subseteq[1,n]$,
with $s=|S|$,
\begin{equation}\label{eq:second-moment-expansion}
\boxed{\displaystyle
|N(S)|\ge
\frac{s\,\delta_n^2}
{\Delta_n+(s-1)D_n}.}
\end{equation}
In particular, for each fixed
$\varepsilon>0$ there exist constants
$c_\varepsilon>0$ and $N_\varepsilon$
such that any Hall-deficient set
$S$ in $H_n$, for $n\ge N_\varepsilon$,
must satisfy
\begin{equation}\label{eq:hall-size}
\boxed{\displaystyle
|S|\ge c_\varepsilon
n^{2/3-\varepsilon}.}
\end{equation}
\end{theorem}
\begin{proof}
If $v$ is adjacent to distinct left
vertices $i>j$, there are cube roots
$a>b$ with
$i+v=a^3$, $j+v=b^3$.
Therefore
\[
i-j=a^3-b^3
=(a-b)(a^2+ab+b^2).
\]
The positive integer $h=a-b$
must divide $i-j$.
For a fixed $h>0$, the function
$b\mapsto(b+h)^3-b^3$ is strictly
increasing, hence at most one
$b$ solves the equation.
This proves \eqref{eq:codegree}.

For a set $S$ of $s$ left vertices,
put $r_v=|\{i\in S:i\sim v\}|$.
Then
\[
\sum_v r_v
=\sum_{i\in S}|N(i)|\ge s\delta_n,
\]
whereas the codegree estimate gives
\[
\sum_v r_v^2
=\sum_{i\in S}|N(i)|
+2\!\!\sum_{\substack{i<j\\i,j\in S}}
|N(i)\cap N(j)|
\le s\Delta_n+s(s-1)D_n.
\]
Cauchy--Schwarz on the nonzero $r_v$
now yields \eqref{eq:second-moment-expansion}.

For completeness, because the possible
right neighbors of $i$ correspond
precisely to cubes in $[i+1,i+n]$,
\[
|N(i)|
=\lfloor(i+n)^{1/3}\rfloor
 -\lfloor i^{1/3}\rfloor.
\]
Concavity of $x^{1/3}$ implies
\[
\delta_n\ge
(2^{1/3}-1)n^{1/3}-1,
\qquad
\Delta_n\le\lfloor(2n)^{1/3}\rfloor.
\]
The classical elementary divisor bound
$\tau(d)=O_\varepsilon(d^\varepsilon)$
gives $D_n=O_\varepsilon(n^\varepsilon)$
for every fixed $\varepsilon>0$.
If $S$ violates Hall, so
$|N(S)|<s$, formula
\eqref{eq:second-moment-expansion}
forces
\[
s>1+\frac{\delta_n^2-\Delta_n}{D_n}.
\]
Since $\delta_n^2$ grows like
$n^{2/3}$ and $\Delta_n=O(n^{1/3})$,
this implies \eqref{eq:hall-size}.
\end{proof}

Combining Theorem~\ref{thm:arith-expansion}
with independent-set reduction and
Theorem~\ref{thm:obstruction-localization}
gives a two-dimensional localization
of any possible Hall obstruction:
it must \emph{both} contain
$n^{2/3-o(1)}$ vertices and
span a range of length
$n-O(n^{2/3})$. These constraints
hold whether or not an all-length
existence threshold is already known.
They do not, by themselves, rule out
the remaining large independent sets.

\begin{theorem}[Undirected cycle-cover expansion]
\label{thm:cyclecover}
If $\mathfrak F_n$ denotes the set of
spanning cycle-cover configurations of $G_n$,
and $c(F)$ is the number of its cycles
of length at least three, then
\begin{equation}\label{eq:cyclecover}
\boxed{P(n)=\sum_{F\in\mathfrak F_n}2^{c(F)}.}
\end{equation}
Thus $P(n)>0$ if and only if $G_n$
has a spanning union of loops, dimers,
and simple cycles. In particular,
\[
\boxed{P(n)\equiv I(n)\pmod 2,}
\]
where $I(n)$ counts valid involutions.
\end{theorem}
\begin{proof}
Every valid permutation decomposes uniquely
into fixed points, transpositions, and
directed cycles of length at least three.
The first correspond to permitted loops,
the second to undirected dimer edges,
and each longer undirected cycle admits
exactly two orientations. Choices on
different components are independent,
giving \eqref{eq:cyclecover}.
All terms with $c(F)>0$ are even,
whereas configurations with $c(F)=0$
are exactly the involutions.
\end{proof}

Let $H_n$ be the symmetric $n$-by-$n$
zero-one matrix with
$(H_n)_{ij}=1$ precisely when $i+j$
is a perfect cube.

\begin{corollary}[Efficient parity via binary algebra]
\label{cor:parity}
\[
\boxed{P(n)=\operatorname{per}(H_n),\qquad
P(n)\bmod2=\det(H_n)\bmod2.}
\]
In particular $P(n)$ is odd exactly
when $H_n$ is nonsingular over
$\mathbb F_2$, and its parity can be
computed by polynomial-time Gaussian
elimination without enumerating
any permutations. Moreover,
\[
\det(H_n)=1\pmod2
\quad\Longrightarrow\quad
\text{there exists a legal involution}
\]
since the number of such involutions
is then odd. The converse need not hold.
\end{corollary}
\begin{proof}
The permanent expansion enumerates
the admissible permutations. In
characteristic two every permutation
sign in the determinant is $+1$,
so determinant and permanent coincide.
\end{proof}

Even more structure appears once any
single legal permutation has been found.
For such a permutation $p$, define the
directed \emph{alternation graph} $D_p$
with vertex set $[1,n]$ and a
non-loop arc $i\to j$ when
$i+p(j)$ is a cube. Let
$\mathcal S(D_p)$ be its strongly
connected components, and let
$B_C$ be the square zero-one matrix
indexed by $i,j\in C$, with
$(B_C)_{ij}=1$ if $i+p(j)$
is a cube (including all diagonal entries).

\begin{theorem}[Exact strong-component factorization]
\label{thm:sccfactor}
For \emph{every} legal reference permutation
$p$ of length $n$,
\begin{equation}\label{eq:sccfactor}
\boxed{\displaystyle
P(n)=\prod_{C\in\mathcal S(D_p)}
\operatorname{per}(B_C).}
\end{equation}
In particular $P(n)=1$ precisely when
every strong component is a singleton;
and every nontrivial component contributes
at least two possible permutations.
\end{theorem}
\begin{proof}
Any second legal permutation $q$
determines a permutation
$f=p^{-1}\circ q$ of $[1,n]$,
and its arcs $i\to f(i)$ are
allowed by the extended alternation
graph with the identity loops included.
Every cycle of $f$ lies entirely
inside one strongly connected
component of $D_p$.
Hence $f$ restricts to independent
permutations within the components,
and their numbers multiply.
Conversely, any such collection of
internal permutation choices assembles
into a legal $q$.
The factorization follows.
A nontrivial strong component contains
a directed cycle, which can be toggled
against the identity on that component,
proving the final assertion.
\end{proof}

This theorem explains the previously
computed $P(236)=512$ from nine
independent directed five-cycles and
$191$ singleton components.
It also unifies the alternating-cycle
criterion of Lemma~\ref{lem:cycle}
with exact counting.
The companion verifier
\texttt{global\_graph\_audit.py} checks
the parity identity and component
factorization in concrete, bounded-size
instances using exact binary arithmetic.

\section{Threshold-free renormalization and counting entropy}
\label{sec:renormalization}

The specific all-$n$ threshold is not needed for
a powerful global closure property of the
cube-sum permutation family.

\begin{lemma}[Independent-fiber renormalization]
\label{lem:renorm-general}
Fix an integer exponent $d\ge2$ and an odd
positive integer $t$; set
$L=t^d$, $A=(L-1)/2$, and
$T_t(m)=Lm+A$.
Let $P_d(m)$ denote the number of permutations
of $[1,m]$ whose index-image sums are
positive perfect $d$th powers.
Then for every nonnegative integer $m$,
\begin{equation}\label{eq:renorm-general}
\boxed{P_d(T_t(m))\ge P_d(A)P_d(m)^L,}
\end{equation}
with $P_d(0)=1$. Moreover,
\[
2T_t(m)+1=L(2m+1),\qquad
T_s(T_t(m))=T_{st}(m).
\]
\end{lemma}
\begin{proof}
Partition $[1,T_t(m)]$ into a prefix
$[1,A]$ and the remaining $m$ consecutive
blocks of length $L$.
Choose an arbitrary valid permutation on
the prefix, and choose independently
$L$ valid permutations
$p_1,\ldots,p_L$ of $[1,m]$.
At $x=A+(i-1)L+r$, $1\le r\le L$, define
\[
q(x)=A+(p_r(i)-1)L+(L+1-r).
\]
This permutes the remaining interval
and satisfies
\[
x+q(x)=L(i+p_r(i)).
\]
Each right-hand side is a perfect $d$th
power, since $L=t^d$.
Different prefix permutations and
different tuples of $p_r$ produce
different full permutations, yielding
\eqref{eq:renorm-general}.
The identities for $T_t$ follow directly
from its affine definition.
\end{proof}

\begin{lemma}[A four-interval $171$-permutation]
\label{lem:171}
The following involution is a cube-sum
permutation of $[1,171]$:
\[
p(i)=
\begin{cases}
8-i,&1\le i\le7,\\
27-i,&8\le i\le19,\\
64-i,&20\le i\le44,\\
216-i,&45\le i\le171.
\end{cases}
\]
\end{lemma}
\begin{proof}
The four displayed intervals are
pairwise disjoint and partition $[1,171]$.
Each is invariant under its indicated
reflection, whose index-image sum is,
respectively, $2^3,3^3,4^3$, or $6^3$.
\end{proof}

\begin{theorem}[Cubic affine closure and entropy monotonicity]
\label{thm:renorm}
For every integer $m\ge1$,
\[
\boxed{P(343m+171)\ge P(m)^{343}.}
\]
Set $T(m)=343m+171$. Then for each
positive integer $k$,
\[
T^k(m)=343^k\Bigl(m+\frac12\Bigr)-\frac12,
\qquad
\boxed{P(T^k(m))\ge P(m)^{343^k}.}
\]
For $P(m)>0$, define the normalized entropy
\[
h(m)=\frac{\log P(m)}{m+\frac12}.
\]
It is nondecreasing on each orbit:
\[
\boxed{h(T^k(m))\ge h(m).}
\]
Consequently, without requiring any
eventual-existence theorem,
\begin{equation}\label{eq:globalentropy}
\boxed{\displaystyle
\limsup_{n\to\infty}\frac{\log P(n)}n
\ge\sup_{\substack{m\ge1\\P(m)>0}}
\frac{\log P(m)}{m+\frac12}.}
\end{equation}
\end{theorem}
\begin{proof}
Apply Lemma~\ref{lem:renorm-general}
with $d=3$, $t=7$, $L=343$, and
$A=(343-1)/2=171$.
Lemma~\ref{lem:171} gives $P(171)\ge1$.
Iterating the inequality and the affine
composition law yields the second display.
Since $T^k(m)+1/2=343^k(m+1/2)$,
the entropy inequality follows on
taking logarithms.
For each fixed $m$ with $P(m)>0$,
the orbit $T^k(m)$ tends to infinity,
so its entropy gives a lower bound
for the global limsup. Taking the
supremum over all such $m$ proves
\eqref{eq:globalentropy}, interpreted
in the extended-real sense if necessary.
\end{proof}

The transformation $2m+1\mapsto343(2m+1)$
is a multiplicative renormalization of
the odd coordinate $2m+1$.
In particular the set of admissible
lengths is closed under
$m\mapsto343m+171$, independently of
the numerical threshold in
Theorem~\ref{thm:exist}.
The explicit arithmetic verifier
\texttt{renormalization.py} checks both
the four-block seed and independent
fiber choices without relying on the
large finite-certificate archive.

\section{A single-seed eventual existence theorem}
\label{sec:one-seed}
There is a substantially more conceptual route
to \emph{eventual existence} than proving the
sharp cutoff at $173$. It uses only one
seven-element permutation, seven rounds of
integer reflection closure, and a uniform
cubic gap inequality. In particular, this
proof never invokes a bipartite matching
algorithm or any of the earlier $173$--$364$
permutation certificates.

Let $\mathcal R_0=\{7\}$ and define a
sequence of finite sets recursively by
\begin{equation}\label{eq:sevenrounds}
\mathcal R_{j+1}=\mathcal R_j\cup
\{k^3-m-1:\ m\in\mathcal R_j,\
m<k^3-m-1\le10975,\ k\ge2\}.
\end{equation}

\begin{lemma}[Seven-round arithmetic certificate]
\label{lem:sevenrounds}
For $j=0,\ldots,7$ the exact cardinalities
$|\mathcal R_j|$ are
\[
1,\ 21,\ 189,\ 995,\ 3280,\ 6588,\ 8533,\ 8924.
\]
Moreover,
\[
\boxed{[8915,10975]\cap\mathbb Z
\subseteq\mathcal R_7.}
\]
\end{lemma}
\begin{proof}
This is a bounded integer-only calculation,
reproduced by \texttt{one\_seed\_eventual.py}.
For each member $m$ and each integer
$k\ge2$ with $k^3\le21950$,
form $n=k^3-m-1$ and insert it
when $m<n\le10975$. After seven
simultaneous rounds the cardinalities
are exactly those displayed and all
$2061$ required lengths are present.
The computation tests only elementary
integer identities and finite set
membership, not the existence of
matchings. Independently, the verifier
\texttt{one\_seed\_interval\_proof.py}
operates entirely on unions of integer
intervals. For every current interval
$[a,b]$ and cube $c$, it restricts to
\[
\ell=\max(a,c-10976),\qquad
h=\min(b,\lfloor(c-2)/2\rfloor)
\]
and, if $\ell\le h$, adds the reflected
interval $[c-h-1,c-\ell-1]$.
The exact numbers of disjoint intervals
after rounds $0$ through $7$ are
\[
1,\ 21,\ 186,\ 870,\ 2077,\ 2037,\ 1031,\ 770.
\]
The last interval of the final union
is exactly $[8915,10975]$, which proves
the finite coverage using \emph{interval
arithmetic alone}. Both verifiers
reproduce the same seven-stage cardinalities.
This finite check is separate from, and
much weaker than, a search for perfect
matchings.
\end{proof}

\begin{theorem}[Eventual existence from one permutation]
\label{thm:single-seed}
The explicit permutation
\[
\boxed{p=(7,6,5,4,3,2,1)}
\]
at length $7$, whose cube sums all equal
$8$, together with the reflection-extension
lemma, proves that
\[
\boxed{P(n)>0\quad\text{for every }n\ge8915.}
\]
Thus eventual existence requires only one
elementary starting permutation and
the seven-round arithmetic certificate,
without locating the optimal cutoff.
\end{theorem}
\begin{proof}
Each element of $\mathcal R_j$ is a
length possessing a valid cube-sum
permutation, by induction on $j$
using Lemma~\ref{lem:reflect}.
Therefore the finite interval
$8915\le n\le10975$ is covered.

For $n\ge10976$, let $k$ be the
least integer such that
$k^3\ge n+8916$ and put
$m=k^3-n-1$. Then $m\ge8915$,
and $k\ge28$.
We prove $k^3\le2n$.
If $k=28$, this follows from
$n\ge10976$, since $28^3=21952$.
If $k\ge29$, minimality gives
$n\ge(k-1)^3-8915$, whence
\[
\begin{aligned}
2n-k^3
&\ge 2((k-1)^3-8915)-k^3\\
&=k^3-6k^2+6k-17832\\
&\ge1685>0.
\end{aligned}
\]
The last polynomial equals $1685$
at $k=29$ and is strictly increasing
for all larger integer $k$.
Thus $8915\le m<n$.
By strong induction $P(m)>0$,
so reflection through the cube $k^3$
produces a valid permutation of length
$n$. This proves the theorem.
\end{proof}

\begin{theorem}[Eventual existence of involutions]
\label{thm:involution-tail}
Let $I_4(n)$ denote the number of
cube-sum \emph{involutions} of $[1,n]$
satisfying $p(4)=4$, where an
involution obeys $p^2=\mathrm{id}$.
Then
\[
\boxed{I_4(n)\ge F(n)>0
\qquad\text{for every }n\ge8915.}
\]
More strongly,
\begin{equation}\label{eq:involution-entropy}
\boxed{\displaystyle
\liminf_{n\to\infty}
\frac{\log I_4(n)}{(2n)^{1/3}}
\ge\frac{\log2}{2}.}
\end{equation}
In particular, the cube-sum graph
admits a spanning configuration
consisting entirely of loops and
disjoint dimers at every $n\ge8915$.
\end{theorem}
\begin{proof}
The sole starting permutation
$p(i)=8-i$ on $[1,7]$ is an
involution and fixes the index $4$.
Every suffix reflection $i\mapsto
c-i$ in Lemma~\ref{lem:reflect}
is itself an involution of the new
interval, which is disjoint from
all previously defined positions.
Thus every constructed permutation
in the one-seed reflection tree
is an involution fixing $4$.
Theorem~\ref{thm:single-seed}
proves $F(n)>0$ for $n\ge8915$,
and the uniqueness of reflection
histories proved in
Lemma~\ref{lem:reflection-count}
gives $I_4(n)\ge F(n)$.
Finally, Theorem~\ref{thm:treeentropy}
supplies \eqref{eq:involution-entropy}.
The undirected loop-and-dimer
formulation follows by decomposing
an involution into fixed points and
transpositions.
\end{proof}

This separates two logically different
facts. The $236$-element cube-sum graph
admits $512$ ordinary permutations but
no involutions, by
Theorem~\ref{thm:236}. Nevertheless,
the one-seed method proves that
noninvolutive obstructions cannot
persist for all large lengths.
It does so without any matching
solver or sharp-cutoff data.

\section{Parity-optimal fixed-point structure}
\label{sec:minimal-fixed-points}

The single-seed eventual-involution theorem
can be strengthened from the existence of
an involution to the existence of an
involution with the \emph{minimum possible
number of fixed points}. This does not
follow from the $173$-threshold theorem
or from general perfect $2$-matchings:
it uses parity-preserving odd-cube
reflections and two tiny explicit seeds.

A reflection step from $m$ to $n$
through an \emph{odd} cube $c=k^3$
satisfies
\[
n=c-m-1\equiv m\pmod2.
\]
Its new interval reflection
$i\mapsto c-i$ has no fixed point,
since $c/2\notin\mathbb Z$.
Thus the parity of the length and
the exact set of fixed points
are preserved under every odd-cube
extension.

\begin{lemma}[Two parity-specific finite closures]
\label{lem:parity-closures}
Starting from the involution
$p_7(i)=8-i$ on $[1,7]$ and allowing
strict extensions through \emph{odd}
cubes only, the reachable lengths
include \emph{every odd integer}
in
\[
\boxed{[70001,100001].}
\]
Starting from the fixed-point-free
involution $p_{26}(i)=27-i$
on $[1,26]$, the same odd-cube
reflection rule reaches \emph{every
even integer} in
\[
\boxed{[230000,302000].}
\]
\end{lemma}
\begin{proof}
This is an exact finite integer
calculation, reproduced by
\texttt{parity\_minimal\_involutions.py}.
For each seed $a$ and even or odd
candidate $n$ of matching parity,
iterate the rule
\[
R_a(a)=1,\quad
R_a(n)=
\bigvee_{\substack{k\ge3\ \mathrm{odd}\\
a\le k^3-n-1<n}}
R_a(k^3-n-1).
\]
The two consecutive-parity intervals
contain respectively $15001$
and $36001$ certified lengths.
The last smaller parity-compatible
lengths which are \emph{not} reachable
below the respective interval
endpoints are $69601$ (odd)
and $228026$ (even).
A second, algorithmically
independent checker
\texttt{parity\_interval\_proof.py}
works with \emph{unions of integer
intervals} in the coordinate
$t=(n-p)/2$ where $p\in\{0,1\}$
is the length parity. An odd-cube
reflection becomes $t\mapsto d-t$
with $d=(k^3-1-2p)/2$.
For every current interval $[a,b]$,
the strict-extension constraints
give the exactly reflected interval
\[
\left[d-\min\left(b,
\left\lfloor\tfrac{d-1}{2}\right\rfloor\right),
d-\max(a,d-T_{\max})\right]
\]
when the lower and upper source
bounds are consistent, with
$T_{\max}=(H-p)/2$.
After just twelve simultaneous
rounds, the merged unions contain
both claimed parity-specific
coverage intervals. The resulting
numbers of reachable parity lengths
are respectively $41057$ and
$121707$, in $2997$ and $7047$
disjoint integer intervals.
The first checker also reconstructs
full involutions at the endpoints
and verifies the exact fixed
point sets. Neither checker uses
a matching solver or the
sharp-threshold certificates.
\end{proof}

\begin{theorem}[Eventual parity-minimal involutions]
\label{thm:minimal-fixed}
For every odd integer $n\ge70001$
there exists a legal cube-sum involution
on $[1,n]$ whose \emph{unique} fixed
point is $4$. For every even integer
$n\ge230000$ there exists a legal
cube-sum involution of $[1,n]$
with \emph{no fixed point}.
Equivalently,
\[
\boxed{\displaystyle
\begin{array}{ll}
n\ge70001\ \text{odd}:&
\min_{p^2=\mathrm{id}}\#\mathrm{Fix}(p)=1,\\
n\ge230000\ \text{even}:&
\min_{p^2=\mathrm{id}}\#\mathrm{Fix}(p)=0,
\end{array}}
\]
where the minima are over legal
cube-sum involutions.
\end{theorem}
\begin{proof}
The preceding lemma supplies
the finite intervals and the
corresponding minimal-fixed-point
involutions, because odd-cube
reflections preserve the set of
fixed points.

For either parity put
$(M,H)=(70001,100001)$ or
$(230000,302000)$, respectively.
If $n>H$ has the required parity,
take the \emph{least odd} $k$
with
\[
k^3\ge n+M+1,\qquad
m=k^3-n-1.
\]
Then $m\ge M$, and $m$ has the
same parity as $n$. We show
$m<n$. Since $k$ is the least
such odd root,
$n\ge(k-2)^3-M$.
Therefore
\[
2n-k^3\ge
2(k-2)^3-k^3-2M
=k^3-12k^2+24k-16-2M.
\]
For the odd case $n\ge100003$,
hence $k\ge57$, and the final
polynomial is increasing for
$k\ge57$ and is $7555>0$
at $k=57$.
For the even case $n\ge302002$,
we have $k\ge83$; its value
at $k=83$ is $31095>0$ and
the polynomial is again increasing.
Consequently $k^3\le2n$,
and $M\le m<n$.

Strong induction within the
fixed parity now produces a
minimal-fixed-point involution
at $m$, followed by the
odd-cube suffix reflection
from $m$ to $n$. This adds
no fixed points. The parity
lower bounds (one fixed point
for odd $n$, none for even $n$)
are tautological for involutions,
so the achieved counts are minimal.
\end{proof}

\begin{corollary}[Odd-cube perfect matchings for every large even length]
\label{cor:odd-cube-matching}
For every even $n\ge230000$, there
is a fixed-point-free involution
$p$ such that \emph{every}
$i+p(i)$ is an \emph{odd}
perfect cube. Equivalently,
the odd and even integers of
$[1,n]$ can be paired bijectively
using edges whose endpoint sums
are odd cubes.

For every odd $n\ge70001$
there is a cube-sum involution
with exactly one fixed point
at $4$, whose cube sums are
odd at \emph{all} indices
$i\ge8$; the first seven sums
are all $8$.
\end{corollary}
\begin{proof}
The even seed $p_{26}(i)=27-i$
has only the odd cube sum $27$,
and every subsequent extension
uses an odd cube. The final
involution is fixed-point-free,
so each transposition joins
an odd integer to an even
integer.

The odd seed $p_7(i)=8-i$
has precisely seven positions
with the even cube sum $8$,
and later odd-cube reflections
introduce no further even
cube sums. Theorem~\ref{thm:minimal-fixed}
supplies the required constructions
at all sufficiently large lengths.
\end{proof}

Define $J_{\min}(n)$ to count
legal cube-sum involutions
with no fixed point when $n$
is even, or with exactly one
fixed point at $4$ when $n$
is odd.

\begin{theorem}[Uniform entropy of minimally fixed involutions]
\label{thm:minimal-entropy}
For all sufficiently large $n$,
\[
\boxed{J_{\min}(n)\ge
\exp\bigl(\Omega(n^{1/3})\bigr).}
\]
More precisely,
\begin{equation}\label{eq:minimal-entropy}
\boxed{\displaystyle
\liminf_{n\to\infty}
\frac{\log J_{\min}(n)}{(2n)^{1/3}}
\ge\frac{\log2}{4}.}
\end{equation}
\end{theorem}
\begin{proof}
For $a\in\{7,26\}$, let
$F_a^{\rm odd}(n)$ count the
distinct strict-reflection
histories from seed $a$ which
use \emph{odd} cube roots only.
All such histories have the
seed's parity and fixed-point
structure; their permutations
are distinct by
Lemma~\ref{lem:reflection-count}.
Theorem~\ref{thm:minimal-fixed}
implies $F_a^{\rm odd}(n)\ge1$
on its corresponding parity
tail $n\ge M_a$.

Fix $s\ge2$, let $\rho_s$
solve
$\sum_{h=1}^{s}\rho_s^{-h}=1$,
and set
\[
B_s=
\max\{8(s-1)^2+24(s-1)+15,\ 4s+80\}.
\]
For $n\ge M_a$, put
$k=\lfloor(2n)^{1/3}\rfloor$.
We prove by strong induction
\[
F_a^{\rm odd}(n)
\ge\rho_s^{(k-B_s)/4}.
\]
This is immediate if $k<B_s$
by positivity. Otherwise,
let $q$ be the largest odd
integer with $q\le k$ and
use the $s$ distinct odd cubes
\[
c_j=(q-2j)^3,\qquad 0\le j<s.
\]
Their predecessors
$m_j=c_j-n-1$ have the same
parity as $n$. Since
$2n<(k+1)^3$ and
$q\ge k-1$,
\[
\begin{aligned}
2m_j-(k-4j-4)^3
&\ge 2(k-2j-1)^3
 -(k+1)^3-1
 -(k-4j-4)^3\\
&=3k^2-(24j^2+72j+45)k\\
&\qquad+
48j^3+168j^2+180j+60
\ge0.
\end{aligned}
\]
The last inequality holds by
$k\ge B_s\ge8j^2+24j+15$.
As $k-4s\ge80$, it follows
that $m_j\ge80^3/2
=256000\ge M_a$.
Also $c_j\le k^3\le2n$,
so $m_j<n$.
Writing $K(m)=\lfloor(2m)^{1/3}\rfloor$,
we have
$K(m_j)\ge k-4(j+1)$.
Consequently, strong induction
and the exact reflection
counting recurrence give
\[
\begin{aligned}
F_a^{\rm odd}(n)
&\ge\sum_{j=0}^{s-1}F_a^{\rm odd}(m_j)\\
&\ge\rho_s^{(k-B_s)/4}
\sum_{j=0}^{s-1}\rho_s^{-(j+1)}
=\rho_s^{(k-B_s)/4}.
\end{aligned}
\]
For fixed $s$, the constant
$B_s$ is independent of $n$.
Letting $n\to\infty$ in
each parity, followed by
$s\to\infty$ and
$\rho_s\uparrow2$, proves
\eqref{eq:minimal-entropy},
because $J_{\min}(n)$ dominates
the corresponding $F_a^{\rm odd}(n)$.
\end{proof}

This shows that the minimal
fixed-point restriction costs
at most a change in the proved
cube-root entropy coefficient:
it does not destroy uniform
superpolynomial combinatorial
multiplicity. The direct
taxicab family in
Section~\ref{sec:taxicab-involutions}
separately establishes stronger
\emph{linear} entropy on
an infinite subsequence.

\begin{corollary}[Efficient one-seed construction]
\label{cor:single-seed-alg}
For every $n\ge8915$, a cube-sum
permutation of length $n$ is
constructible using $O(n)$
unit-cost arithmetic operations
and $O(\log\log n)$ reflection
layers beyond a fixed finite lookup.
\end{corollary}
\begin{proof}
For $n\ge10976$, take the least
cube $k^3\ge n+8916$ as in the
theorem. Its predecessor satisfies
\[
8915\le m=k^3-n-1
=O(n^{2/3}),
\]
since consecutive cubes differ by
$O(k^2)=O(n^{2/3})$.
Hence $\log m\le(2/3)\log n+O(1)$,
which yields $O(\log\log n)$
recursive stages before reaching
the finite interval. Every reflection
writes its newly added positions
exactly once, so the output work
is $O(n)$; exact cube roots add
only polylogarithmic overhead.
The finite interval is traced through
at most seven predecessor links,
ending at the sole seed of length $7$.
\end{proof}

Theorem~\ref{thm:single-seed} is weaker
numerically than the fully classified
cutoff in Theorem~\ref{thm:exist},
but stronger conceptually in a different
direction: it establishes eventual
existence without any large explicit
matching witness. The residual finite
verification concerns \emph{reflection
reachability only}, not bipartite
matching, and is entirely independent
of the sharp-threshold classification.
A completely computation-free proof
of eventual existence would still
require eliminating this last
seven-round finite arithmetic check.

\section{Monotone signed Waring representations}
\label{sec:signed-waring}
The reflection orbit has a direct reformulation
as an additive-number-theoretic problem.
The important feature is not just the
presence of signed cubes, but the
\emph{monotonicity of every intermediate
partial reflection}.

\begin{proposition}[Reflection paths and alternating cubes]
\label{prop:signed-waring}
Fix any positive starting length $a$ and
set $x_0=2a+1$.
A sequence of $r$ strict reflection
extensions from $a$ to $n$ exists if
and only if there are integers
$k_1,\ldots,k_r\ge2$ for which
\[
x_j=2k_j^3-x_{j-1}>x_{j-1}
\qquad(1\le j\le r),
\]
and $x_r=2n+1$.
Equivalently,
\[
\boxed{\displaystyle
2n+1=(-1)^r(2a+1)
+2\sum_{j=1}^{r}(-1)^{r-j}k_j^3,}
\]
where the intermediate quantities
specified by the recursion must be
strictly increasing positive odd integers.
\end{proposition}
\begin{proof}
Writing $x_j=2n_j+1$,
the reflection condition
$n_{j-1}+n_j+1=k_j^3$
is equivalent to
$x_{j-1}+x_j=2k_j^3$.
The strict extension $n_j>n_{j-1}$
is exactly $x_j>x_{j-1}$.
Iterating this affine identity
gives the alternating sum formula.
Conversely every such increasing
sequence recovers the necessary
reflection predecessors and cubes.
\end{proof}

For the canonical seed $a=7$, a path
of exactly seven reflections reaches
$n$ precisely when it admits
\[
\boxed{\displaystyle
n+8=
k_7^3-k_6^3+k_5^3-k_4^3+
k_3^3-k_2^3+k_1^3}
\]
with the corresponding strictly
increasing intermediate odd coordinates.
This converts the global question
about \emph{bounded reflection depth}
into an explicit restricted
signed-Waring representation problem.
It is stronger than mere eventual
existence: Theorem~\ref{thm:single-seed}
proves a growing but
$O(\log\log n)$-bounded depth,
not a universal depth of seven.

Experiments with exact interval arithmetic
suggest that the seven-step orbit already
covers every length $n\ge8915$,
but this stronger unbounded statement
is \emph{not proved here}. It would
require a uniform theorem about
alternating cubic representations
with their intermediate-order
constraints, rather than an
unbounded numerical extrapolation.

\section{The reflection tree as a standalone counting model}
\label{sec:reflection-tree}

The one-seed construction contains not merely one
permutation for every sufficiently large length,
but a recursively countable subfamily
whose size itself grows uniformly.

Let $F(7)=1$, and for $n>7$ define
\begin{equation}\label{eq:reflection-tree}
\boxed{\displaystyle
F(n)=
\sum_{\substack{7\le m<n\\m+n+1=k^3\ \text{for some }k\ge2}}
F(m).}
\end{equation}
We regard $F(n)$ as zero at lengths not
reachable from the seven-element seed.

\begin{lemma}[Exact reflection-history enumeration]
\label{lem:reflection-count}
For every $n\ge7$, $F(n)$ equals exactly
the number of distinct cube-sum permutations
obtainable from the seed
$(7,6,5,4,3,2,1)$ by a finite sequence
of strict reflection extensions.
In particular,
\[
\boxed{0\le F(n)\le P(n),\qquad
F(n)>0\ \ (n\ge8915).}
\]
\end{lemma}
\begin{proof}
Every strictly increasing history has
a unique final predecessor $m<n$,
and its last extension uses the cube
$c=m+n+1$. The final entry of the
permutation is then
\[
p_n=c-n=m+1.
\]
Thus the last predecessor $m$ can be
recovered uniquely from $p_n$, while
the previous permutation is recovered
by restricting to its first $m$
positions. Hence two distinct
reflection histories cannot produce
the same final permutation.
The recursive count is exactly
\eqref{eq:reflection-tree}.
The eventual positivity is
Theorem~\ref{thm:single-seed}.
\end{proof}

\begin{theorem}[Uniform entropy inside one reflection tree]
\label{thm:treeentropy}
Fix $s\ge2$ and let $\rho_s\in(1,2)$
be determined by
$\sum_{h=1}^{s}\rho_s^{-h}=1$.
Put
\[
B_s^\ast=\max\{2s^2+4s-1,\ 2s+27\},
\qquad K(n)=\lfloor\sqrt[3]{2n}\rfloor.
\]
For every $n\ge8915$,
\[
\boxed{\displaystyle
F(n)\ge\rho_s^{(K(n)-B_s^\ast)/2}.}
\]
Consequently,
\begin{equation}\label{eq:tree-entropy}
\boxed{\displaystyle
\liminf_{n\to\infty}
\frac{\log F(n)}{(2n)^{1/3}}
\ge\frac{\log2}{2}.}
\end{equation}
The same bound holds for $P(n)$,
but this proof uses only the
seven-element seed, its seven-round
arithmetic reachability certificate,
and elementary cubic inequalities.
\end{theorem}
\begin{proof}
Write $k=K(n)$, $\rho=\rho_s$,
and $B=B_s^\ast$. If $k<B$,
the right-hand side of the claimed
inequality is at most $1$,
and $F(n)\ge1$ by
Lemma~\ref{lem:reflection-count}.

If $k\ge B$, for $j=0,\ldots,s-1$
take $c_j=(k-j)^3$ and
$m_j=c_j-n-1$. The polynomial
inequality in the proof of
Theorem~\ref{thm:entropy}
(applied with the same $j,k$)
gives
\[
K(m_j)\ge k-2j-2.
\]
Since $k-2s\ge27$, this
guarantees $m_j\ge
\lceil27^3/2\rceil>8915$.
Also $c_j\le k^3\le2n$,
so $m_j<n$.
By strong induction, the exact
recurrence \eqref{eq:reflection-tree}
therefore yields
\[
\begin{aligned}
F(n)&\ge\sum_{j=0}^{s-1}F(m_j)\\
&\ge \rho^{(k-B)/2}
\sum_{j=0}^{s-1}\rho^{-(j+1)}
=\rho^{(k-B)/2}.
\end{aligned}
\]
Finally, for every fixed $s$,
$B_s^\ast$ does not depend on $n$,
and $\rho_s\uparrow2$ as $s\to\infty$.
The limit inferior is at least
$(\log\rho_s)/2$ for each $s$,
which proves \eqref{eq:tree-entropy}.
\end{proof}

The reflection-tree recurrence
is a useful independent combinatorial
object: its members are explicit
permutations, not merely abstract
paths or potential matchings.
Unlike a claim that every large
permutation can be generated by
reflection (which is false, for
example at $n=173$ from this seed),
this theorem counts a rigorously
defined \emph{subfamily} only.

\section{The entropy barrier of pure reflection}
\label{sec:reflection-barrier}

The reflection-tree estimates admit a sharp
structural complement: its combinatorial entropy
has exactly the cube-root growth \emph{order}.
Consequently, the exponentially many
permutations in the independent-cycle constructions
cannot be produced from finitely many seeds by
suffix reflections alone.

\begin{theorem}[Strict ordering of reflection cubes]
\label{thm:cube-order}
Consider any chain of strict reflection extensions
\[
n_0<n_1<\cdots<n_r=n
\]
whose step from $n_{j-1}$ to $n_j$
uses the cube $c_j=n_{j-1}+n_j+1$.
Then
\[
\boxed{c_1<c_2<\cdots<c_r\le2n.}
\]
If the initial seed permutation is fixed,
the set $\{c_1,\ldots,c_r\}$ uniquely
determines the entire history and its
final permutation.
\end{theorem}
\begin{proof}
The step from $n_{j-1}$ to $n_j$
is strict, hence
$c_j=n_{j-1}+n_j+1\le2n_j\le2n$.
For consecutive steps,
\[
c_{j+1}-c_j
=n_{j+1}-n_{j-1}>0.
\]
Once the starting length $n_0$ and
the strictly increasing list of cubes
are specified, every subsequent
length is uniquely determined by
$n_j=c_j-n_{j-1}-1$.
The reflected permutation is then
uniquely determined by the starting
permutation and these successive
lengths. A set of cubes therefore
determines at most one history.
\end{proof}

\begin{theorem}[Optimal cube-root order for the one-seed tree]
\label{thm:reflection-barrier}
Let $F(n)$ count cube-sum permutations
obtained by strict reflection extensions
from the unique seed
$(7,6,5,4,3,2,1)$.
For all $n\ge7$,
\begin{equation}\label{eq:reflection-upper}
\boxed{\displaystyle
F(n)\le 2^{K(n)-1},
\qquad K(n)=\lfloor(2n)^{1/3}\rfloor.}
\end{equation}
Together with
Theorem~\ref{thm:treeentropy}, this yields
\[
\boxed{\displaystyle
\log F(n)=\Theta(n^{1/3})\qquad(n\to\infty).}
\]
The same upper estimate holds with a
factor $S$ for the union of strict-reflection
descendants of any fixed finite collection
of $S$ initial cube-sum permutations:
\begin{equation}\label{eq:finite-seeds-upper}
\boxed{\displaystyle
|\mathcal R_{\mathcal S}(n)|
\le S\,2^{K(n)-1}.}
\end{equation}
\end{theorem}
\begin{proof}
All possible strict reflection cubes up to
length $n$ are among
$2^3,3^3,\ldots,K(n)^3$,
a list of $K(n)-1$ cubes.
By Theorem~\ref{thm:cube-order},
every reflection history from a fixed
seed corresponds injectively to a subset
of these cube roots (including the empty
subset for the initial length).
There are at most $2^{K(n)-1}$ such
subsets. Summing over $S$ seeds gives
\eqref{eq:finite-seeds-upper}.
Theorem~\ref{thm:treeentropy} provides a
strictly positive lower liminf for
$\log F(n)/(2n)^{1/3}$ as $n\to\infty$,
and the upper bound gives a finite
limsup. Hence the logarithm is
$\Theta(n^{1/3})$.
\end{proof}

\begin{corollary}[A strict entropy separation]
\label{cor:entropy-separation}
For any fixed finite family of
cube-sum seed permutations,
\[
\lim_{n\to\infty}
\frac1n\log\bigl(
1+|\mathcal R_{\mathcal S}(n)|\bigr)=0.
\]
By contrast, the entire legal family
has positive upper exponential entropy:
\[
\boxed{\displaystyle
\limsup_{n\to\infty}
\frac{\log P(n)}n
\ge\frac{148877\log2}{681707}>0.}
\]
Consequently, no fixed finite-seed
strict-reflection scheme can capture a
positive fraction of the exponential-scale
freedom exhibited by the independent
three-cycle family.
\end{corollary}
\begin{proof}
The upper estimate
\eqref{eq:finite-seeds-upper} gives
$\log(1+|\mathcal R_{\mathcal S}(n)|)
=O(n^{1/3})=o(n)$.
The lower bound for the full family
is Theorem~\ref{thm:three-cycle}.
In particular, on its explicit
subsequence $n=681707u^3$,
the logarithm of the full permutation
count is at least a positive constant
times $n$, while the reflected
finite-seed subfamily has logarithm
$o(n)$.
\end{proof}

More generally, with perfect $d$th
powers in place of cubes, the strictly
increasing-root proof gives at most
\[
S\,2^{\lfloor(2n)^{1/d}\rfloor-1}
\]
permutations from $S$ fixed seeds.
The exponent $n^{1/d}$ is a
\emph{methodological upper bound}
for reflection-only generation, not
an upper bound for all admissible
power-sum permutations.

\begin{theorem}[Fibonacci majorant for the reflection tree]
\label{thm:fibonacci-majorant}
Let $F(n)$ be the number of
strict-reflection descendants of the
seven-point seed, as in
\eqref{eq:reflection-tree}, and
$K(n)=\lfloor(2n)^{1/3}\rfloor$.
With $\mathrm{Fib}_1=\mathrm{Fib}_2=1$
and $\mathrm{Fib}_{j+2}=
\mathrm{Fib}_{j+1}+\mathrm{Fib}_{j}$,
\begin{equation}\label{eq:fibonacci-majorant}
\boxed{\displaystyle
F(n)\le\mathrm{Fib}_{K(n)-2}
\qquad(K(n)\ge3),}
\end{equation}
and $F(n)\le1$ when $K(n)=2$.
In particular, writing
$\varphi=(1+\sqrt5)/2$,
\begin{equation}\label{eq:fibonacci-entropy}
\boxed{\displaystyle
\limsup_{n\to\infty}
\frac{\log F(n)}{(2n)^{1/3}}
\le\log\varphi.}
\end{equation}
Combined with
Theorem~\ref{thm:treeentropy},
\[
\boxed{\displaystyle
\frac{\log2}{2}\le
\liminf_{n\to\infty}\frac{\log F(n)}{(2n)^{1/3}}
\le
\limsup_{n\to\infty}\frac{\log F(n)}{(2n)^{1/3}}
\le\log\varphi.}
\]
\end{theorem}
\begin{proof}
Write $k=K(n)$, and let
$m=(k-j)^3-n-1\ge7$
be any valid strict predecessor
appearing in \eqref{eq:reflection-tree}.
Necessarily $j\ge0$ and
$(k-j)^3>n\ge k^3/2$,
so $k-2j>0$. Moreover,
\[
\begin{aligned}
2m
&=2(k-j)^3-2n-2\\
&\le2(k-j)^3-k^3-2\\
&=(k-2j)^3-6j^2(k-j)-2\\
&<(k-2j)^3.
\end{aligned}
\]
Thus the integer cube-root height of
every predecessor satisfies
\[
\boxed{K(m)\le k-(2j+1).}
\]

Define $B_\ell=0$ for $\ell\le1$,
$B_2=1$, and for $\ell\ge3$ set
\[
B_\ell=\sum_{j\ge0} B_{\ell-2j-1}.
\]
All sums here are finite. The recurrence
gives $B_3=B_4=1$ and
$B_\ell=B_{\ell-1}+B_{\ell-2}$
for every $\ell\ge5$;
hence $B_\ell=\mathrm{Fib}_{\ell-2}$
for every $\ell\ge3$.
In particular $B_\ell$ is
nondecreasing for $\ell\ge2$.

Strong induction on $n$ gives
$F(n)\le B_{K(n)}$: the seed
$n=7$ has height $K(7)=2$ and
$F(7)=B_2=1$; other lengths
with height two have $F(n)=0$;
and for $K(n)=k\ge3$,
\[
F(n)\le
\sum_{j\ge0}B_{K(m_j)}
\le\sum_{j\ge0}B_{k-(2j+1)}
=B_k.
\]
This proves the finite inequality.
Finally Binet's formula implies
$\log\mathrm{Fib}_k=k\log\varphi+O(1)$,
so \eqref{eq:fibonacci-entropy}
follows. The lower inequality is
Theorem~\ref{thm:treeentropy}.
\end{proof}

This replaces the coarse
$2^{K(n)-1}$ encoding ceiling
with a strict Fibonacci bound.
It captures the important fact that
the distinct predecessor branches
do not just consume one increasing
cube root: they reduce the available
height by the odd increments
$1,3,5,\ldots$.
It remains an open question whether
the exact upper or lower limiting
coefficients can be attained.

\section{A dimension barrier for bounded reflection depth}
\label{sec:depth-barrier}

Even though the seven-element seed reaches
every sufficiently large length, a genuinely
\emph{uniformly bounded} number of reflections
is a separate Diophantine question.
There is an elementary counting obstruction
to any bound of at most three reflections.

\begin{theorem}[Reflection-depth sparsity]
\label{thm:reflection-depth}
Fix one cube-sum permutation as an initial
seed, and let $\mathcal D_r(N)$ be the
set of output lengths at most $N$
obtainable from it by at most $r$
strict reflection extensions.
For $K=\lfloor(2N)^{1/3}\rfloor$,
\begin{equation}\label{eq:depth-sum}
\boxed{\displaystyle
|\mathcal D_r(N)|
\le \sum_{j=0}^{r}\binom{K-1}{j}.}
\end{equation}
For each fixed $r\le2$,
\[
\boxed{\displaystyle
|\mathcal D_r(N)|=O(N^{r/3})=o(N),}
\]
and for three or fewer reflections,
\begin{equation}\label{eq:three-depth}
\boxed{\displaystyle
|\mathcal D_3(N)|\le
\frac N3+O(N^{2/3}).}
\end{equation}
In particular, no one-seed reflection
construction of depth at most three
can cover every sufficiently large length.
At least four strict reflections are
necessary for any one-seed, bounded-depth
cofinite construction.
\end{theorem}
\begin{proof}
By Theorem~\ref{thm:cube-order}, a
history producing a length $n\le N$
uses a strictly increasing sequence
of cubes chosen from the $K-1$
possible cubes $2^3,\dots,K^3$.
A history with exactly $j$ steps
is thus encoded by a $j$-element
subset, and there are at most
$\binom{K-1}{j}$ choices.
Summing yields \eqref{eq:depth-sum}.
For fixed $r\le2$ this is
$O(K^r)=O(N^{r/3})$.
For $r=3$, the sum is
\[
\binom{K-1}{3}+O(K^2)
=\frac{K^3}{6}+O(K^2)
\le\frac N3+O(N^{2/3}),
\]
because $K^3\le2N$.
Thus at least $(2/3+o(1))N$
of the first $N$ lengths cannot be
reached in at most three steps.
\end{proof}

More generally, if a fixed finite
family of $S$ initial permutations
is allowed, the right side of
\eqref{eq:depth-sum} is multiplied
by at most $S$. This makes clear
why the counting barrier does not
extend unchanged to arbitrarily
many input seeds.

The numerical seven-round closure
in Lemma~\ref{lem:sevenrounds}
raises the sharper question whether
\emph{a fixed depth of seven} already
suffices for all sufficiently large
lengths. The present theorem does
not prove this; it places a rigorous
lower bound of four on any such
single-seed uniform depth.

\section{A hand-checkable finite interval proof}
\label{sec:finite-interval-proof}

We give a standalone \emph{finite algebraic}
certificate for every base length $173\le n\le364$.
It is independent of the augmenting-path
search, Hall tests, and permutation matching
algorithms. Unlike the global theory, this
part still uses an explicit finite case table:
the numerical expressions below can each
be verified solely by integer arithmetic.

\begin{lemma}[Boundary cancellation criterion]
\label{lem:endpoint-cancellation}
Let $n\ge1$ and choose a positive integer $t$,
endpoints $0=b_0<b_1<\cdots<b_t=n$,
and perfect cubes $c_1,\dots,c_t$.
Define
\[
L_j=c_j-b_j,\qquad R_j=c_j-b_{j-1}.
\]
Suppose that $1\le L_j<R_j\le n+1$
for all $j$. Then the piecewise formula
\[
\boxed{p(i)=c_j-i\quad
\text{whenever }b_{j-1}<i\le b_j}
\]
is a cube-sum permutation of $[1,n]$
\emph{if and only if}
\begin{equation}\label{eq:endpoint-cancellation}
\boxed{
\{1,R_1,\ldots,R_t\}_{\rm multi}
=
\{n+1,L_1,\ldots,L_t\}_{\rm multi}.}
\end{equation}
\end{lemma}
\begin{proof}
The $j$th input segment is mapped
bijectively onto the consecutive integers
$[L_j,R_j-1]$, with $i+p(i)=c_j$.
Let $f(x)$ count the number of
image segments containing the integer
$x$, for $1\le x\le n$.
Set $f(0)=f(n+1)=0$.
The difference $f(x)-f(x-1)$ is
the number of left boundaries at $x$
minus the number of right boundaries
at $x$, where the image intervals
are half-open $[L_j,R_j)$.
The multiset identity
\eqref{eq:endpoint-cancellation}
is precisely the assertion that
this difference is $1$ at $x=1$,
$-1$ at $x=n+1$, and $0$ at all
other integers. Hence $f(x)=1$
at every $1\le x\le n$:
the image intervals partition $[1,n]$.
Conversely, a partition has precisely
those boundary differences, and thus
satisfies the stated multiset identity.
\end{proof}

\begin{lemma}[Reflected-block rearrangement]\label{lem:main-block-rearrange}
Choose positive input-block lengths
$\ell_1,\ldots,\ell_t$ of sum $n$ and an arbitrary
output order $\pi$ of these $t$ blocks.
Let $A_j$ be the sum of input lengths before
block $j$, and let $B_j$ be the sum of output
lengths before block $j$ in order $\pi$.
If all numbers
\[
  C_j=1+A_j+B_j+\ell_j
\]
are perfect cubes, then on input block $j$
the reflection $p(i)=C_j-i$ gives a cube-sum
permutation of $[1,n]$.
\end{lemma}
\begin{proof}
The input block is $[A_j+1,A_j+\ell_j]$,
which maps exactly to the prescribed output
block $[B_j+1,B_j+\ell_j]$.
Those output blocks partition $[1,n]$
by construction. Each index sum is $C_j$.
\end{proof}

\paragraph{Family I (four blocks).}
Let $U<V$ be integer cubes and put
\[
a=2U-V-1,\qquad b=2V-3U+1.
\]
If $a,b>0$, use input lengths $(a,b,a,b)$
and output order $(3,4,1,2)$.
Lemma~\ref{lem:main-block-rearrange} gives the cubic
sums $(U,V,U,V)$ and length $2(V-U)$.
The pairs $(U,V)=(64,125),(125,216),(216,343)$
yield lengths $122,182,254$ respectively.

\paragraph{Family II (five blocks).}
For cubes $X<Y<Z$, define
\[
a=2Y-Z-1,\quad
b=X+2Z-4Y+1,\quad c=Y-X.
\]
If $a,b,c>0$, the lengths $(a,b,c,a,c)$
in output order $(4,2,5,1,3)$ produce
sums $(Y,X,Z,Y,Z)$ and total length
$2Y-X-1$. With $(Y,Z)=(64,125)$ and
$X=27,8$, the results are $n=100,119$.

\paragraph{Family III (six blocks).}
For integer cubes $X,W,Y,Z$, put
\[
\begin{aligned}
a&=2Y-Z-1,& b&=X-2a-1,\\
c&=W+2a-2X+1,& d&=Y+X-W-2a-1.
\end{aligned}
\]
If these are positive, choose lengths
$(a,b,c,d,a,d)$ and output order
$(5,2,3,6,1,4)$. The cubic sums are
$(Y,X,W,Z,Y,Z)$ and the length is
$2Z-2Y+X-W$. Taking $(Y,Z)=(64,125)$,
the choices $(X,W)=(8,64),(27,64),(8,27)$
give $n=66,85,103$, respectively.

\paragraph{Family IV (twelve blocks and three binary switches).}
Let $X,T,Q$ be integers and put
\[
\begin{aligned}
u&=X-7,&v&=T-2X+7,&w&=X+27-T,\\
z&=Q+2T-2X-121,&h&=X+152-T-Q.
\end{aligned}
\]
Assume $u,v,z\ge0$ and $w,h>0$.
Use input lengths
\[
(3,u,v,w,3,w,z,h,3,27,3,h)
\]
and output block order
\[
(11,2,3,6,9,4,7,12,5,10,1,8).
\]
The reflected-block lemma computes the twelve
nominal index-image sums as
\[
(125,X,T,64,125,64,Q,216,125,216,125,216).
\]
When the second, third, or seventh block is nonempty,
its nominal sum $X$, $T$, or $Q$, respectively,
must be a cube. Any zero-length block is omitted
from both orders and imposes no cube condition.
The total length is
\[
n=276+X-T-Q.
\]
Take $X\in\{7,8\}$, choose
$T\in\{2X-7,27\}$, and independently choose
$Q\in\{121-2T+2X,125\}$. All eight resulting
length choices satisfy the positivity and cube
conditions. The $2^3$ distinct lengths are
\[
\boxed{131,132,150,151,155,156,174,175}.
\]
This is one symbolic rearrangement construction
with three optional empty blocks, not eight
separate finite endpoint certificates.

\paragraph{Family V (15 blocks, two constrained switches).}
Take $\alpha,\beta\in\{0,1\}$ with
$\alpha\le\beta$ and
$t\in\{0,7+4\beta+2\alpha\}$.
Define the ordered input block lengths
\[
\begin{gathered}
\ell_1=2\beta,\quad
\ell_2=7-\alpha-4\beta-\alpha\beta,\quad
\ell_3=\alpha,\quad\ell_4=\alpha\beta,\\
\ell_5=13+4\beta,\quad\ell_6=\alpha,\quad
\ell_7=23-4\beta-\alpha,\quad\ell_8=13+4\beta,\\
\ell_9=48-\alpha-45\beta+\alpha\beta,\quad
\ell_{10}=2\beta,\quad\ell_{11}=\beta(39-\alpha),\\
\ell_{12}=t,\quad\ell_{13}=30+\alpha-t,\quad
\ell_{14}=13+4\beta,\quad
\ell_{15}=91+\alpha-t.
\end{gathered}
\]
The output block order is
\[
(10,4,6,2,8,3,11,1,9,14,7,12,15,5,13).
\]
All lengths are nonnegative, with empty blocks omitted.
Because $\alpha\beta=\alpha$, the nominal reflected
sums are
\[
\begin{gathered}
21+43\beta,\ 8,\ 28-\alpha,\ 8,\ 216,\\
21+2\alpha+4\beta,\ 125,\ 64,\ 125,\ 105-41\beta,\\
125,\ 209-2\alpha-4\beta+t,\ 343,\ 216,\ 343.
\end{gathered}
\]
Every \emph{nonempty} block has a cubic sum:
blocks $1,10,11$ are used only if $\beta=1$;
blocks $3,4,6$ are used only if
$\alpha=\beta=1$; and block $12$ is empty
if $t=0$, otherwise its sum is $216$.
The total length is
\[
\boxed{n=238+\alpha+2\beta-t},
\]
so the three allowed pairs
$(\alpha,\beta)=(0,0),(0,1),(1,1)$
with their two possible $t$ values give
\[
\boxed{228,229,231,238,240,241}.
\]
This is a symbolic cubic-sum construction;
it does not require verifying six separate endpoint lists.

\paragraph{Family VI (ten blocks, nested binary degeneration).}
Let $\alpha,\beta\in\{0,1\}$ with $\beta\le\alpha$.
Use input block lengths
\[
\begin{gathered}
(7,\ 12\alpha,\ 159-24\alpha-99\beta,\
14\beta,\ 28\beta,\ 63\beta,\\
10+24\alpha-6\beta,\ 8\beta,\
159-24\alpha-107\beta,\
10+24\alpha+93\beta)
\end{gathered}
\]
and output order $(1,2,8,5,6,9,4,10,3,7)$.
Direct substitution into the reflected-block lemma
gives the ten nominal sums
\[
\begin{gathered}
8,\ 8+19\alpha,\ 343,\
167-12\alpha+61\beta,\
167-12\alpha-30\beta,\\
167-12\alpha+61\beta,\ 512,\
177+12\alpha+27\beta,\ 343,\ 512.
\end{gathered}
\]
Block $2$ is nonempty only when $\alpha=1$
and has cubic sum $27$; blocks $4,5,6,8$
are nonempty only when $\beta=\alpha=1$
and have sums $216,125,216,216$. Other
nonempty blocks have constant cube sums;
all empty blocks are omitted. Since
\[
n=345+12\alpha-6\beta,
\]
the three admissible choices give the
lengths $345,357,351$, using one
parametric construction rather than
three independent endpoint lists.

\subsection{Universal unimodular determinant construction}
\label{subsec:unimodular-factory}
The remaining $57$ integer witnesses admit a common
algebraic formula, without separately testing
their endpoint cancellation identities.
For an output-block order $\pi$ on $t$ labels,
write $\rho(i)=\pi^{-1}(i)$ and set
\[
M_{ji}=[i<j]+[\rho(i)<\rho(j)]+[i=j],
\quad 1\le i,j\le t.
\]
Here brackets denote indicator functions.

\begin{lemma}[Unimodular block-order matrices]
\label{lem:unimodular-block}
For every finite output-block order $\pi$,
$\det M\in\{0,1\}$.
\end{lemma}
\begin{proof}
With $\boldsymbol1$ denoting the all-ones
vector, write $M=\boldsymbol1\boldsymbol1^T+K$.
The matrix $K$ is skew-symmetric, and for
$i<j$ its $(i,j)$ entry is $-1$ exactly
when $\rho(i)<\rho(j)$, and $0$ otherwise.
It is therefore an oriented chord-intersection
matrix: place first chord endpoints in
input order around one arc of a circle,
and second endpoints in output order around
the complementary arc.

Recall that an even-order chord-intersection
matrix has Pfaffian $0$ or $\pm1$.
Indeed, thickening the chord diagram into an
oriented ribbon surface, the chord loops form an
integral homology basis. Their intersection form
is integrally congruent to a direct sum of
unimodular symplectic pairs
$\left(\begin{smallmatrix}0&1\\-1&0\end{smallmatrix}\right)$
and zero blocks from the boundary radical.
Its full Pfaffian is thus $0$ or $\pm1$.

If $t$ is even, $\operatorname{adj}(K)$ is
skew-symmetric; the rank-one determinant
identity gives
\[
\det M=\det K+
\boldsymbol1^T\operatorname{adj}(K)\boldsymbol1
=\operatorname{Pf}(K)^2\in\{0,1\}.
\]
If $t$ is odd, $\det K=0$. Adjoin a new
chord crossing all original chords, whose
skew matrix is
\[
H=\begin{pmatrix}
0&\boldsymbol1^T\\
-\boldsymbol1&K
\end{pmatrix}.
\]
The same rank-one and bordered determinant
identities give
\[
\det M=\boldsymbol1^T\operatorname{adj}(K)
\boldsymbol1=\det H=\operatorname{Pf}(H)^2
\in\{0,1\}.
\]
\end{proof}

\begin{theorem}[Symbolic determinant families]
\label{thm:unimodular-block-factory}
Fix $\pi$ with $\det M=1$.
For arbitrary integer cube roots
$k_1,\ldots,k_t\ge2$, let
$d=(k_j^3-1)_{j=1}^t$. Define
\[
\ell_j=\det M^{(j)}(d),
\]
where $M^{(j)}(d)$ is obtained from $M$
by replacing column $j$ with $d$.
If every $\ell_j>0$, then with
$N=\sum_j\ell_j$ there is a permutation
of $[1,N]$ with cube sums $k_j^3$ on
successive input blocks of these lengths.
\end{theorem}
\begin{proof}
By the preceding lemma and Cramer's rule,
$M\ell=d$ holds identically over the integers.
If $A_j,B_j$ are the input and output block
prefix lengths, this is precisely
\[
k_j^3=1+A_j+B_j+\ell_j.
\]
Reflection $p(i)=k_j^3-i$ maps the
$j$th input block bijectively to the
$j$th block in output order $\pi$.
The output blocks partition $[1,N]$.
\end{proof}

The $57$ cubic-root words and block orders
in Appendix~\ref{app:algebraic-shapes}
are therefore a collection of \emph{symbolic}
families, rather than $57$ independently
postulated endpoint identities.
For each printed root word,
integer determinant evaluation gives
$\det M=1$, $\ell_j\ge1$, and the
specified length $N$ (the smallest
$\ell_j$ across the full table equals $1$).
The separate exact
\texttt{verify\_unimodular\_factory.py}
checks these finite positivity conditions.
The $57$ parameter choices remain finite data:
this is not a case-free proof of their existence.

\begin{theorem}[Uniform cubic dilation]
\label{thm:unimodular-dilation}
Under Theorem~\ref{thm:unimodular-block-factory},
put $v=M^{-1}\boldsymbol1$ and
$\varepsilon=t\bmod2$. For every integer $h\ge1$,
\[
\ell_j(h)=\ell_j+
(h^3-1)(\ell_j+v_j)>0
\]
defines another cube-sum permutation, with
root word $(hk_1,\ldots,hk_t)$, of length
\[
\boxed{N(h)=(N+\varepsilon)h^3-\varepsilon.}
\]
\end{theorem}
\begin{proof}
Pfaffian minor expansions in the preceding
chord-diagram proof imply
$v_j\in\{-1,0,1\}$: the Cramer minors
are signed Pfaffians of chord diagrams
with a marked chord deleted, divided
by a nonzero Pfaffian of absolute value $1$.
Moreover $\sum_j v_j=\varepsilon$.
Indeed, $Mv=\boldsymbol1$ and skew symmetry
give $S^2=S$ for $S=\boldsymbol1^Tv$.
If $t$ is even, $K$ is invertible and
$v=K^{-1}\boldsymbol1$ yields $S=0$.
If $t$ is odd, invertibility of $M$ and
singularity of $K$ exclude $S=0$:
a kernel vector of $K$ has nonzero pairing
with $\boldsymbol1$. Thus $S=1$.

Since $\ell_j\ge1$ and $v_j\ge-1$,
all displayed scaled lengths are positive.
Finally
\[
M\ell(h)=h^3M\ell+(h^3-1)\boldsymbol1
=h^3(k_j^3)_{j=1}^t-\boldsymbol1.
\]
The determinant construction applies again.
Summing the lengths gives the formula
for $N(h)$.
\end{proof}

\begin{lemma}[Affine-word rigidity]\label{lem:rigid}
Fix a finite list of integer cubes $C_1,\ldots,C_t$.
Suppose all breakpoints $b_j(n)$ are affine
functions of $n$ with $b_0(n)=0$ and $b_t(n)=n$.
The resulting fixed-cube reflection word can
be a valid permutation for at most two distinct
values of $n$.
\end{lemma}
\begin{proof}
Every legal permutation satisfies the sum identity
\[
 \sum_{j=1}^{t} C_j\bigl(b_j(n)-b_{j-1}(n)\bigr)
 =\sum_{i=1}^n(i+p(i))=n(n+1).
\]
Its left side is affine in $n$; its right
side is a monic quadratic. The equality
cannot hold for three distinct values.
\end{proof}
The obstruction concerns \emph{fixed} cube sums
and affine endpoints only, not varying cube
sums or composition of reflection templates.

In total these \emph{six symbolic identities}
give twenty-five base lengths
\[
66,85,100,103,119,122,131,132,150,151,155,156,174,175,182,228,229,231,238,240,241,254,345,351,357
\]
without twenty-five separate endpoint-table checks.

\begin{theorem}[Reduced finite arithmetic certificate]
\label{thm:finite-interval-atlas}
The $57$ cubic-root words and output
block orders in Appendix~\ref{app:algebraic-shapes},
together with six additional symbolic block
families, prove that every $n\ge173$ admits
a cube-sum permutation.
For each of the $57$ shapes, all reflected
block lengths are defined by the determinant
formula of Theorem~\ref{thm:unimodular-block-factory}.
They are strictly positive and total $795$
reflected blocks. The older $57$-row endpoint
atlas is independent redundant evidence.
\end{theorem}
\begin{proof}
The algebraic shape appendix contains
$10$ lower and $47$ middle parameter sets.
All its order matrices have determinant $1$,
and all Cramer-determinant block lengths
are positive integers. Hence
Theorem~\ref{thm:unimodular-block-factory}
constructs each corresponding permutation
directly, without using the redundant
endpoint-multiset atlas.

The six preceding formulas establish
another twelve seeds below $173$, namely
$66,85,100,103,119,122,131,132,150,151,155,156$,
and thirteen seeds in the target interval,
$174,175,182,228,229,231,238,240,241,254,345,351,357$.
Their permutations are guaranteed by
the symbolic block-rearrangement lemma.
Thus $22$ lower and $60$ middle seeds
have been established without matching.

Starting with these $82$ seeds and
the empty permutation, we use the
reflection lemma in strictly increasing
order of the target length.
The following predecessor sets give all
other cases before the uniform tail.

\paragraph{First: 26 reflected lengths below 173.}
\begin{quote}
\footnotesize
\begin{verbatim}
cube     already constructed predecessors m
  8      0
 27      0,7
 64      0,7,19,26
125      0,7,19,26,37,44,56
216      44,56,63,66,68,80,82,85,87,
         100,103,105
\end{verbatim}
\end{quote}
Together with the $10$ exceptional lower seeds
and the twelve algebraic-template lower seeds,
these supply $48$ distinct lengths below 173.

\paragraph{Second: 49 reflected lengths in
$[173,255]\cup[339,364]$.}
\begin{quote}
\footnotesize
\begin{verbatim}
cube     previously constructed predecessors m
216      0,7,19,26,37
343      0,87,98,100,103,105,110,112,
         115,116,117,118,119,120,121,122,
         123,124,125,126,128,129,130,131,
         132,133,135,147,149,150,151,152,
         155,156,159
512      147,149,150,151,152,155,156,
         159,171
\end{verbatim}
\end{quote}
The $47$ upper exceptional rows, the thirteen
upper algebraic-template lengths and these $49$
reflections are exactly the $109$ integers
in $[173,255]\cup[339,364]$.
There is no circular dependency:
each predecessor $m$ is strictly smaller
than its constructed length $n$.

For example, $211=343-131-1$,
$340=512-171-1$, and $342=343-0-1$.
The first example replaces a previously
necessary-looking $53$-segment independent
row by one suffix reflection.


In particular $48$ positive lower lengths
and the $109$ lengths in
$[173,255]\cup[339,364]$ are now
covered. For $256\le n\le338$, let
$m=511-n\in[173,255]$, so that
$m<n$ and $m+n+1=512=8^3$.
Lemma~\ref{lem:reflect} then covers
all $192$ target base lengths.
For every $n\ge365$, the existing
cube-gap inequality of Lemma~\ref{lem:gap}
and strong induction extend the
permutations to all later lengths.
\end{proof}

The primary proof still contains $57$
selected finite algebraic shapes with a
positive-length check for each:
the algebraic templates do not establish
a completely case-free classification.
The primary audit script
\texttt{verify\_unimodular\_factory.py}
is optional and reads only the cubic-order
shape table, not the old endpoint table or original
$191$ augmenting words or a matching
archive. The older $109$-row table
is preserved as an independent
redundant check.

\subsection{Forced reflection peeling at the critical base}
\label{subsec:forced-reflection-peeling}

\begin{lemma}[Forced neighbor and reflected pair]
\label{lem:forced-neighbor}
Let $V$ be a finite set of positive integers, with
cube-sum adjacency $i\sim j$ whenever $i+j$ is a
cube. If $i\in V$ has exactly one neighbor $j\in V$,
every cube-sum permutation $p:V\to V$ satisfies
$p(i)=j$ and $p(j)=i$. Deleting $i,j$ (or only
$i$ when $i=j$) therefore preserves exactly the
number of cube-sum permutations of the induced
residual vertex set.
\end{lemma}
\begin{proof}
The $i$th row can use only image $j$.
Since the cube-sum graph is symmetric, column $i$
can only be supplied by row $j$. Both assignments
are therefore forced, and restriction or extension
gives mutually inverse maps between permutations.
\end{proof}

\begin{theorem}[Three forced bands]
\label{thm:forced-bands}
For every $173\le n\le181$, the following
three successive reflected bands are forced
in \emph{every} cube-sum permutation:
\[
\begin{array}{c|c|c}
\text{round}&\text{first endpoints }i& i+p(i)\\ \hline
1 &[125,342-n]&216\\
2 &[34,215-n]&64\\
3 &[95,276-n]&216
\end{array}
\]
Here each row is deleted together with its
reflected images before the next round.
Exactly $582-3n$ forced reflected pairs are
identified by these three bands.
\end{theorem}
\begin{proof}
Since $2n<512$, all relevant cube sums
are $8,27,64,125,216,343$. On the first
band, cubes up to $125$ give no positive
neighbor, and $343-i>n$; thus the unique
neighbor is $216-i$. Their image interval
is $[n-126,91]$.
On the second band, the possible $125-i$
is in this deleted image interval,
and $216-i>n$, leaving only $64-i$;
its images form $[n-151,30]$.
On the third band, the possible $125-i$
lies in the deleted second image interval,
and $343-i>n$, leaving $216-i$.
All partners are distinct and belong to
the still remaining vertices.
The band sizes are
$218-n$, $182-n$, and $182-n$, so their
sum is $582-3n$.
\end{proof}

\begin{theorem}[Exact counts at $173,174,175$]
\label{thm:forced-base}
For the cube-sum permutation counting function,
\[
\boxed{P(173)=1,\quad P(174)=1,\quad P(175)=3.}
\]
In particular the existence at $n=173$
is established without a determinant shape
or a precomputed endpoint-multiset identity.
\end{theorem}
\begin{proof}
At $n=173$, after the three forced bands
the additional batches of first endpoints
and their required cube sums are
\[
\begin{array}{c|l}
\text{round}& \text{first endpoints : cube}\\ \hline
4&\{4,5\}:8,\ [8,12]:27\\
5&\{45,46,106,107,108,122\}:216\\
6&\{31\}:64,\ \{172,173\}:216\\
7\!-\!15&(92,216),(1,8),(20,125),(111,125),
(13,125),\\
& (104,125),(6,8),(123,216),(32,64)
\end{array}
\]
Each first endpoint has exactly one neighbor
remaining when its round is reached;
this follows by direct subtraction from the
six cubes at most $343$. All $173$ vertices
are removed. Thus the resulting permutation
is unique, with fixed points $4,32,108$.

At $n=174$, the corresponding rounds are
\[
\begin{array}{c|l}
1\!-\!3&[125,168]:216,\ [34,41]:64,\ [95,102]:216\\
4&\{4\}:8,\ [8,11]:27\\
5&\{45,46,47,106,107,108\}:216\\
6&\{15\}:27,\ \{172,173,174\}:216\\
7\!-\!19&(113,216),(22,27),(3,125),(94,125),
(33,125),\\
&(124,125),(7,27),(105,216),(14,27),
(112,216),\\
&(21,27),(2,125),(93,125).
\end{array}
\]
Again every vertex is removed by forced choices,
so the permutation is unique, with fixed
points $4,108$.

At $n=175$ the first three bands are
$[125,167]$, $[34,40]$, $[95,101]$.
Use rounds $4$--$19$ of the preceding table,
but omit $\{4\}$ from round $4$ and
replace its interval $[8,11]$ by $[8,10]$.
The remaining ten vertices are precisely
\[
\{4,11,16,23,41,48,102,114,168,175\}.
\]
Their induced cube-sum graph is the loop at
$4$, the edge $(4,23)$, and the chordless
nine-cycle
\[
23-41-175-168-48-16-11-114-102-23.
\]
If $p(4)=4$, the odd cycle has exactly
two directed cycle covers, one for each
orientation. Otherwise $p(4)=23$
forces $p(23)=4$; the eight remaining
vertices form a path with a unique
cover by four reciprocal pairs. Thus
there are exactly $2+1=3$ permutations.
\end{proof}

\begin{corollary}[One fewer algebraic shape]
\label{cor:56-plus-forcing}
The proof of the finite interval $[173,364]$
may use the $56$ cubic-root/output-order
shapes in Appendix~\ref{app:reduced-shapes},
the six algebraic block families generating
$25$ further seed lengths, and the unique
forced matching of $173$ from
Theorem~\ref{thm:forced-base}.
Together these establish the same $82$
seed lengths as before. The excluded
$173$ determinant shape is redundant.
\end{corollary}

\subsection{The interval-overlap avalanche and its finite obstruction}
\label{subsec:interval-overlap-avalanche}

\begin{lemma}[Exact cubic reflection of a good interval]
\label{lem:reflection-overlap}
Suppose every integer $m\in[a,b]$, $0\le a\le b$,
has a legal cube-sum permutation. For a cube $c$,
every integer $n$ in the nonempty interval
\[
\bigl[
\max\{1,\lceil c/2\rceil,c-b-1\},\
c-a-1
\bigr]
\]
also has a legal cube-sum permutation.
\end{lemma}
\begin{proof}
Take $m=c-n-1$. By the displayed bounds,
$a\le m\le b$ and $0\le m<n$.
Keep the old permutation on $[1,m]$ and set
$p(i)=c-i$ on $[m+1,n]$. The suffix is mapped
bijectively to itself, with constant sum $c$.
\end{proof}

\begin{lemma}[Interval-overlap avalanche]
\label{lem:interval-avalanche}
Suppose every length in $[a,b]$ is legal,
where $1\le a\le b$. If a cube $c$ satisfies
\[
\boxed{a+b+2\le c\le2b+2},
\]
then every length in the larger interval
$[a,c-a-1]$ is legal.
\end{lemma}
\begin{proof}
By Lemma~\ref{lem:reflection-overlap},
the new interval of lengths includes $b+1$
because $a\le c-b-2\le b$ and
$c\le2(b+1)$. It continues through
$c-a-1\ge b+1$; it therefore meets the
already established interval immediately,
with no gap.
\end{proof}

\begin{theorem}[Infinite avalanche from the finite base]
\label{thm:interval-avalanche-tail}
If all lengths $173\le n\le364$ are legal,
then every length $n\ge173$ is legal,
without any additional finite certificate.
\end{theorem}
\begin{proof}
For $a=173$, $b=364$, the cube $9^3=729$
lies in the avalanche window $[539,730]$.
Consequently $[173,555]$ is legal.
Assume $[173,k^3-174]$ is legal for $k\ge9$.
The cube $(k+1)^3$ obeys
\[
k^3+1\le(k+1)^3\le2k^3-346,
\]
since the last inequality is equivalent to
\[
f(k)=k^3-3k^2-3k-347\ge0,
\]
with $f(9)=112$ and
$f(k+1)-f(k)=3k^2-3k-5>0$ for $k\ge9$.
Lemma~\ref{lem:interval-avalanche}
extends the good interval through
$(k+1)^3-174$. Induction on $k$ proves
the assertion.
\end{proof}

\paragraph{Why overlap does not by itself remove the finite base.}
At $n=65$, the domains $[1,63]$ of the
$64$-reflection and $[60,65]$ of the
$125$-reflection overlap and jointly cover
$[1,65]$. Yet
\[
S=\{3,40,59,64\},\qquad
N(S)=\{5,24,61\}
\]
under the available cube sums $8,27,64,125$,
so Hall's inequality fails.
Likewise no smaller legal length can produce
$n=173$ by a single suffix reflection:
the only possible predecessors are
$m=42$ and $m=169$, via cubes $216,343$.
For $m=42$, the six-row set
\[
[1,2]\cup[20,21]\cup[38,39]
\]
has just the four cube-sum neighbors
$[6,7]\cup[25,26]$.
For $m=169$, the fourteen-row set
\[
\begin{gathered}
[8,9]\cup[38,39]\cup[45,46]\cup[99,100]\\
{}\cup[129,130]\cup[136,137]\cup[160,161]
\end{gathered}
\]
has only the twelve neighbors
\[
\begin{gathered}
[18,19]\cup[25,26]\cup[55,56]\\
{}\cup[79,80]\cup[86,87]\cup[116,117].
\end{gathered}
\]
These are direct Hall obstructions.
There is also no cube in the interval
$[513,678]$, the necessary avalanche
window for extending $[173,338]$ through
$339$. Thus the avalanche lemma completely
removes the infinite tail's finite checks,
but \emph{does not derive the original
$173$--$364$ base interval without the
separately constructed finite witnesses}.
Such a claim would require an additional
local augmentation or gluing theorem.

\begin{proposition}[A finite parity obstruction at $n=236$]
\label{prop:236-no-involution-elementary}
The base length $236$ has cube-sum
permutations (as certified by its
$15$-segment atlas row) but admits
no cube-sum \emph{involution}.
\end{proposition}
\begin{proof}
In the undirected cube-sum graph
on $[1,236]$, remove the following
set $S$ of $86$ vertices:
\[
\begin{aligned}
S={}&[1,3]\cup[9,13]\cup[19,24]
\cup[35,39]\cup[45,49]\\
&\cup[56,60]\cup[71,74]\cup[81,85]
\cup[107,115]\\
&\cup[117,121]\cup[126,130]
\cup[136,140]\\
&\cup[147,151]\cup[162,166]
\cup[172,176]\\
&\cup[198,201]\cup[208,212],
\end{aligned}
\]
where $[a,b]$ denotes the
integer interval.
The only possible edge sums are
$8,27,64,125,216,343$.
Direct substitution into
$j=c-i$ shows that $G_{236}-S$
has $87$ singleton components,
one five-vertex component
$\{14,50,75,141,202\}$, and
six even components of sizes
$4,6,10,12,12,14$.
Of the singleton components,
only $\{4\}$ admits a loop
($2\cdot4=8$); the five-vertex
component has no loop.
Hence there are $86+1=87$
odd components which admit
no internal loop. In any
involution, at least one
vertex from each such component
must be paired with a distinct
vertex of $S$. Since $|S|=86$,
this is impossible.
The atlas row gives ordinary
permutations, so the involution
obstruction does not contradict
their existence.
\end{proof}

\section{A short, local existence certificate}\label{sec:words}
The preceding reflection argument is infinite and elementary.
The finite existence verification can itself be made much more
transparent than an array of $192$ long permutations:
we give \emph{one} explicit involution at size $173$,
then $191$ short alternating-path words.
Only $1844$ single-digit cube roots are needed.

\begin{lemma}[An explicit order-$173$ involution]\label{lem:seed173}
There is a cube-sum permutation on $[1,173]$, given on
each listed interval $I$ by $p(i)=c-i$. Intervals
paired by a double arrow use the same cube $c$:
\[
\begin{array}{c|c}
\text{intervals }I&c\\\hline
{[1,7]}&8\\
{[22,42]}&64\\
{[106,110]}&216\\
{[8,12]}\ \longleftrightarrow\ [15,19]&27\\
{[13,14]}\ \longleftrightarrow\ [111,112]&125\\
{[20,21]}\ \longleftrightarrow\ [104,105]&125\\
{[43,103]}\ \longleftrightarrow\ [113,173]&216
\end{array}
\]
\end{lemma}
\begin{proof}
The eleven intervals partition $[1,173]$.
Reflection $i\mapsto c-i$ preserves each of the
first three intervals and exchanges the two intervals
in each of the last four rows. Thus it defines an
involution, and every sum $i+p(i)$ is a cube.
\end{proof}

\begin{lemma}[An alternating-word extension]\label{lem:words}
Suppose $p$ is a cube-sum permutation of $[1,n-1]$.
Let $w=r_1\cdots r_t$ be a finite word of positive
integer cube roots. Starting from $u_0=n$, define
$v_j=r_j^3-u_{j-1}$. For $j<t$, require $v_j\in[1,n-1]$
and let $u_j=p^{-1}(v_j)$. Suppose all these
$u_0,u_1,\dots,u_{t-1}$ are distinct,
all $v_1,\dots,v_t$ are distinct, and $v_t=n$.
Then replacing the old matching edges by
\[
q(u_{j-1})=v_j\quad (1\le j\le t),
\]
while keeping all other old assignments, yields
a cube-sum permutation of $[1,n]$.
\end{lemma}
\begin{proof}
The old matching edges removed along this path are
$u_j\mapsto v_j$ for $j<t$; these form a bijection
between the old visited left and right vertices.
Their replacements join $u_{j-1}$ to $v_j$ and have
sum $r_j^3$ by construction. The new left vertex $n$
and the new right vertex $n$ are incorporated, without
repeating any endpoint. Every untouched vertex retains
its old matching edge.
\end{proof}

The \texttt{augmenting\_cube\_words.txt} appendix lists
one word for each $n=174,\dots,364$.
Every root is a single decimal digit in $\{2,\dots,8\}$.
The words are checked using Lemma~\ref{lem:words}
and the predecessor matching obtained from the previous
word; this is \emph{local} verification, not a
matching search or a scan to a large cutoff.
The independent, standard-library-only
\texttt{verify\_word\_paths.py} checks
all required edge equalities and endpoint distinctness.
The seven-line interval seed plus $1844$ root digits
certifies all $192$ base cases of
Theorem~\ref{thm:exist}; combined with the
cube-gap induction this proves existence for
\emph{every} $n\ge173$.
The older long witness archives remain as secondary
independent cross-checks, not as prerequisites
of this compact proof.

\section{Finite certificates and exact membership}
Our finite data are provided as UTF-8 JSON arrays of integers.
A valid witness $p$ is checked solely by
\[
\operatorname{sort}(p)=(1,\dots,n),\qquad
\forall i\quad i+p_i\in\cubes.
\]
For a negative instance, let
$N_n(S)=\{j\in[1,n]:\exists i\in S,\ i+j\in\cubes\}$.
Hall's marriage theorem gives the immediate obstruction
\begin{equation}\label{eq:hall}
|N_n(S)|<|S|\quad\Longrightarrow\quad P(n)=0.
\end{equation}
The file \texttt{exact\_classification\_below\_173.json}
provides $48$ explicit valid permutations and $124$
explicit Hall-deficient subsets, exactly partitioning
$1\le n\le172$. Every assertion is directly checked
by \texttt{verify\_extended.py}. For instance, for $n=172$
the $27$ left indices
\[
\begin{gathered}
S=\{1,5,8,13,20,31,35,38,42,43,45,92,96,99,\\
104,111,122,126,129,133,134,136,153,157,160,165,172\}
\end{gathered}
\]
have precisely $26$ distinct possible right neighbors.
Thus $172\notin\good$, making the threshold below sharp.

For the unbounded tail, only the $109$ witness sizes
\begin{equation}\label{eq:core}
173\le n\le255,\qquad 339\le n\le364
\end{equation}
require explicit stored permutations.
Indeed for each $256\le n\le338$, put
$m=511-n\in[173,255]$; then
$m+n+1=512=8^3$, so Lemma~\ref{lem:reflect}
constructs the missing $83$ cases.
This compresses the previously saved $192$ base cases
by $83/192>43\%$. All $109$ stored witnesses in
\texttt{core\_109.json}, and all $192$ reconstructed
witnesses, are checked with exact integer arithmetic.
The earlier full $192$-row file is retained as an
additional independent-format audit source.

There is a further useful compression when the $48$
verified positive sizes below $173$ are reused.
An ascending reflection-closure algorithm starts from
those small solutions and, at each $173\le n\le364$,
uses any already certified $m<n$ for which $m+n+1$
is a cube. Precisely $132$ lengths are obtained by
reflection and only $60$ further lengths need explicit
witnesses. Here the empty matching is allowed:
$n=342$ follows directly from $c=343$,
with $p(i)=343-i$. The $60$ additional
witnesses are archived in
\texttt{core\_60\_with\_empty.json}, and
\texttt{verify\_extended.py} certifies this alternate
compression independently of the $109$-case presentation.


\begin{proposition}[Limit of reflection-only compression]\label{prop:atoms}
Suppose the empty matching and the $48$ known positive
lengths below $173$ are provided as initial data. Among proofs of all lengths
$173,\dots,364$ that use only the reflection extension
of Lemma~\ref{lem:reflect} and independently supplied
new full-length seed permutations, at least $60$
extra seeds are necessary, and $60$ suffice.
\end{proposition}
\begin{proof}
Process $n=173,\dots,364$ in ascending order. If no
cube $c$ yields a previously certified positive
predecessor $m=c-n-1\in[0,n-1]$, then reflection
cannot establish $n$ from earlier data, so a new seed
is forced. Otherwise choose such a predecessor and
perform reflection, without adding a seed.
This deterministic closure rule yields $44$ forced
seeds in $[173,255]$, none in $[256,338]$ (use
$c=512$ and $m=511-n$), and $16$ in $[339,364]$.
Thus $44+16=60$, and the attached short
integer-only checker reproduces the forced
index list. Every forced seed is unavoidable by
the definition of the restricted proof rule.
\end{proof}
This is \emph{not} a lower bound on arbitrary
mathematical proofs. This $60$-seed minimum concerns the
restricted reflection-only closure with all $48$ positive
small lengths and the empty seed already available; it is
distinct from the historical $57$ algebraic shape choices.
The alternating-path
extension in Lemma~\ref{lem:words} is a different
operation and indeed replaces those $57$ large
independent seeds by one explicit involution
and a small word table.

\begin{theorem}[Complete A006063 membership]\label{thm:exist}
For every integer $n\ge173$, $P(n)>0$.
Below $173$, the complete membership list is
\begin{align*}
\good\cap[1,172]=\{&
7,19,26,37,44,56,63,66,68,80,82,85,87,98,100,\\
&103,105,110,112,115,116,117,118,119,120,121,\\
&122,123,124,125,126,127,128,129,130,131,132,\\
&133,135,147,149,150,151,152,155,156,159,171\}.
\end{align*}
In sequence-index notation this gives
\[
\boxed{\operatorname{A006063}(j)=j+124\quad(j\ge49).}
\]
\end{theorem}
\begin{proof}
The exact finite witnesses establish $173\le n\le364$,
and the Hall certificates settle the smaller sizes.
For $n\ge365$, Lemma~\ref{lem:gap} gives a cube
$c\in[n+174,2n]$. The predecessor
$m=c-n-1$ satisfies $173\le m<n$.
By strong induction $P(m)>0$, and
Lemma~\ref{lem:reflect} yields $P(n)>0$.
There are precisely $48$ positive indices below $173$,
so the $49$th and all later sequence terms are consecutive.
\end{proof}

\section{Multiplicity and the second OEIS classification}
A cube-sum permutation $p$ is a perfect matching in the
bipartite graph with left and right vertices $1,\dots,n$,
where $(i,j)$ is an edge exactly when $i+j\in\cubes$.
For any matching $p$, form a directed graph $D_p$ on the
left vertices, putting $i\to j$ whenever $i\ne j$
and $(i,p_j)$ is a cube-sum edge.

\begin{lemma}[Alternating-cycle criterion]\label{lem:cycle}
For a valid permutation $p$, there exists another distinct
valid permutation if and only if $D_p$ has a directed cycle.
\end{lemma}
\begin{proof}
Along a directed cycle $i_1\to\dots\to i_t\to i_1$,
replace $p_{i_s}$ by $p_{i_{s+1}}$ (cyclic indices).
Every moved pair is an allowed edge, and all right
vertices remain distinct. Conversely, if $q\ne p$ is
another valid permutation, the nontrivial cycles of
the permutation $p^{-1}\circ q$ consist of directed
cycles in $D_p$. Thus acyclicity is a rigorous
uniqueness certificate.
\end{proof}

The file \texttt{second\_solutions.json} records directed
cycles, not complete duplicate permutations. Direct
finite checking establishes: among the $48$ positive
lengths below $173$, exactly
\[
E=\{112,115,116,117,119,124,125,126,127,128,129,130,133\}
\]
have a directed alternating cycle; every other positive
length below $173$, as well as $173$ and $174$, has
an acyclic $D_p$ and therefore a unique permutation.

Furthermore, all $190$ base cases $175\le n\le364$
have \emph{two distinct} alternating-cycle certificates,
archived as \texttt{second\_solutions.json} and
\texttt{third\_solutions.json}. They yield three different
permutations, counting the original matching.
To extend the multiplicity assertion:
for $365\le n\le553$, the cube $729$ gives
$m=728-n\in[175,363]$. The two exceptional lengths
$554,555$ have separate pairs of alternating-cycle
certificates obtained by extending the unique solutions
at $174,173$ respectively.
For $n\ge556$, the stronger part of
Lemma~\ref{lem:gap} gives
$m=c-n-1\ge175$. Three distinct solutions for $m$
remain distinct after reflection, finishing an induction.

\begin{theorem}[Complete A096680 membership]\label{thm:multi}
The positive integers with at least two cube-sum permutations
are exactly
\[
\boxed{E\ \cup\ \{175,176,177,\dots\}.}
\]
Moreover $P(n)\ge3$ for every $n\ge175$, while
$P(173)=P(174)=1$. The sequence terms
satisfy
\[
\boxed{\operatorname{A096680}(j)=j+161\quad(j\ge14).}
\]
\end{theorem}


\paragraph{Sharpness of the constant three.}
The bound $P(n)\ge3$ for all $n\ge175$ cannot be
improved to four: $P(175)=3$ exactly.
For the stored matching at $175$, the directed
alternation graph has precisely one nontrivial
strongly connected component,
\[
\{4,11,16,23,41,48,102,114,168,175\},
\]
and the other $165$ vertices are singleton components.
Every cycle of a permutation alternative lies in a
strongly connected component, so the permanent reduces
to the $10$-by-$10$ component's allowed-edge matrix
(including the matching self-edges).
A direct bitmask dynamic program gives its permanent
as $3$. This computation is independently reproduced in
\texttt{verify\_extended.py} using exact integers.

\section{A nine-gadget exceptional case at $n=236$}
\label{sec:236}

An \emph{involution} is a permutation satisfying
$p^2=\mathrm{id}$. At $n=236$ there are exactly $512$
cube-sum permutations, \emph{none} of which is an involution.
This follows from a short alternating-cycle decomposition,
not from a general permanent computation.

Define a reference permutation $p$ of $[1,236]$ by
$p(i)=c-i$ on the following fifteen intervals:
\[
\begin{array}{r|r@{\qquad}r|r}
I&c&I&c\\ \hline
{[1,7]}&8 &[8,9]&216\\
{[10,18]}&64 &[19,43]&125\\
{[44,45]}&216 &[46,70]&125\\
{[71,81]}&216 &[82,106]&125\\
{[107,134]}&343 &[135,136]&216\\
{[137,170]}&343 &[171,172]&216\\
{[173,197]}&343 &[198,208]&216\\
{[209,236]}&343 &&
\end{array}
\]
These interval images partition $[1,236]$, and
every prescribed sum is a cube.

\begin{lemma}[Translation of alternating cycles]
\label{lem:transport}
Suppose $p(i)=c_j-i$ on pairwise disjoint
integer intervals $I_j$ for $1\le j\le r$.
If $i_j\in I_j$ form a directed alternating
cycle $i_1\to i_2\to\cdots\to i_r\to i_1$
for $p$, then every integer shift $t$
with $i_j+t\in I_j$ for all $j$ produces
another alternating cycle on $i_j+t$.
If $t$ ranges through a set $T$ of such
shifts, the translated cycles are disjoint
and yield at least $2^{|T|}$ distinct valid
cube-sum permutations.
\end{lemma}
\begin{proof}
Because $p(i_{j+1}+t)=c_{j+1}-i_{j+1}-t$,
\[
(i_j+t)+p(i_{j+1}+t)
=i_j+c_{j+1}-i_{j+1}
=i_j+p(i_{j+1}),
\]
so every translated alternating edge has
the same cube sum as its original.
Since the $I_j$ are pairwise disjoint,
different shifts cannot share a vertex
within any $I_j$. The cycles are
therefore vertex-disjoint and can be
switched independently.
\end{proof}

In the present $236$-element template,
\[
10\to137\to46\to198\to71\to10
\]
is one alternating $5$-cycle, with
successive cube sums
\[
216,\quad216,\quad64,\quad343,\quad125.
\]
The five vertices belong respectively to
the five different reflection blocks
$[10,18]$, $[137,170]$, $[46,70]$,
$[198,208]$, $[71,81]$.
The common shifts $t=0,\dots,8$
produce nine disjoint alternating cycles
\emph{by pure algebra}, giving immediately
\[
\boxed{P(236)\ge2^9=512}
\]
without any matching enumeration.
The remaining work for the exact count
is only to exclude further cycles.

\begin{lemma}[Alternating digraph at $236$]\label{lem:236}
Let $D_p$ be the digraph with an arrow $i\to j$
whenever $i\ne j$ and $i+p(j)$ is a cube.
Its strongly connected components comprise
$191$ singletons and exactly nine five-vertex sets
\[
S_t=\{10+t,46+t,71+t,137+t,198+t\},
\qquad 0\le t\le8.
\]
Within each $S_t$ the nontrivial arrows form
one simple directed $5$-cycle, without chords.
\end{lemma}
\begin{proof}
This is a finite integer check using the only six
possible cubes $8,27,64,125,216,343$.
The interval formula specifies every edge; the
independent \texttt{count\_236.py} verifier checks
all $236$ rows, the component partition, and
the five-cycle properties with integer arithmetic.
It invokes neither a matching solver nor
a permanent-computation routine.
\end{proof}

\begin{theorem}[Exact count and involution obstruction]\label{thm:236}
\[
\boxed{P(236)=2^9=512.}
\]
There is no cube-sum involution on $[1,236]$.
\end{theorem}
\begin{proof}
Given any other cube-sum permutation $q$,
$f=p^{-1}\circ q$ is a permutation whose arrows
lie in $D_p$ or are identity loops. Every cycle of
$f$ is confined to one strongly connected component.
Each singleton is therefore fixed, while each $S_t$
has exactly two choices: five fixed points or the
unique directed $5$-cycle. The choices are independent,
so there are $2^9$ permutations $q$.

For the central set $S_4=\{14,50,75,141,202\}$,
the reference permutation restricts to
\[
14\mapsto50\mapsto75\mapsto141
\mapsto202\mapsto14.
\]
The only two possible restrictions of $q$ to $S_4$
are this $5$-cycle and its inverse, neither of
which is an involution. Thus no $q$ is an involution.
\end{proof}

The failure of involutivity at $236$ is therefore a
specific five-cycle obstruction, not a failure of
cube-sum permutation existence. The proof uses
only a fifteen-row interval template and an easily
audited finite graph-decomposition lemma.

\section{A taxicab principle for exponentially many involutions}
\label{sec:taxicab-involutions}

Positive exponential entropy is not restricted to
general permutations with odd directed cycles.
Equal sums of two perfect powers generate
large families of actual \emph{involutions}.
The argument is valid for any fixed exponent
$d\ge2$.

\begin{theorem}[Perfect-power taxicab switching]
\label{thm:taxicab-switching}
Fix $d\ge2$, and suppose positive perfect
$d$th powers $A,B,C,D$ satisfy
\[
\boxed{A+B=C+D,\qquad
\frac32A\le C<\frac12B.}
\]
Let $I_d(n)$ count permutations $p$
of $[1,n]$ with $p^2=\mathrm{id}$
and every $i+p(i)$ a perfect
$d$th power. For every positive integer
$u$, with $T=u^d$,
\begin{equation}\label{eq:taxicab-involutions}
\boxed{\displaystyle
I_d((B-A)T)\ge
2^{\left\lfloor(AT-1)/2\right\rfloor}.}
\end{equation}
There is an injective encoding of the
$\lfloor(AT-1)/2\rfloor$ binary choices
into these involutions which multiplies
Hamming distances by exactly four.
Consequently,
\[
\boxed{\displaystyle
\limsup_{n\to\infty}
\frac{\log I_d(n)}n
\ge\frac{A\log2}{2(B-A)}>0.}
\]
\end{theorem}
\begin{proof}
Since $A<C<B/2$, we have $B>2A$.
Put $N=(B-A)T$ and start with
the two-block reflection involution
\[
p(i)=
\begin{cases}
AT-i,&1\le i<AT,\\
BT-i,&AT\le i\le N.
\end{cases}
\]
The two blocks are separately invariant,
and the sums are the $d$th powers
$AT=(uA^{1/d})^d$ and
$BT=(uB^{1/d})^d$.

Write $H=\lfloor(AT-1)/2\rfloor$.
For each integer $s=1,\ldots,H$,
consider the four distinct vertices
\[
a_s=s,\quad b_s=AT-s,\quad
c_s=CT-s,\quad
d_s=(B-C)T+s.
\]
Here $a_s,b_s$ belong to the first
block, while $c_s,d_s$ belong
to the second block, since
$H\le AT/2\le(C-A)T$.
The inequalities $2C<B$ and
$2H<AT$ imply that the four
coordinate ranges obtained by varying
$s$ are pairwise disjoint.
Thus all $H$ quadruples are
vertex-disjoint.

The reference involution has
transpositions
$(a_s\,b_s)$ and $(c_s\,d_s)$.
For any chosen quadruple, replace
these by
$(a_s\,c_s)$ and $(b_s\,d_s)$.
The new sums are
\[
a_s+c_s=CT,\qquad
b_s+d_s=(A+B-C)T=DT,
\]
both perfect $d$th powers.
This is again an involution, and
the independent choices on the $H$
disjoint quadruples give $2^H$
different legal involutions.

Changing one choice changes precisely
the four corresponding permutation
values, yielding the exact Hamming
distance multiplier four.
Finally $H=\frac12AT+O(1)$ and
$N=(B-A)T$, so taking logarithms
along the infinite sequence $u\to\infty$
proves the entropy bound.
\end{proof}

\begin{corollary}[Two explicit cubic identities]
\label{cor:taxicab-cubes}
For every integer $u\ge1$, with $T=u^3$,
\[
\boxed{\displaystyle
I_3(54406T)\ge
2^{\lfloor(4913T-1)/2\rfloor}.}
\]
This follows from the small identity
\[
17^3+39^3=26^3+36^3,
\]
using $(A,B,C,D)=(4913,59319,17576,46656)$.
A second identity,
\[
99^3+172^3=120^3+163^3,
\]
gives the stronger optional numerical estimate
\[
\boxed{\displaystyle
\limsup_{n\to\infty}\frac{\log I_3(n)}n
\ge\frac{970299\log2}{8236298}>0.}
\]
Neither identity is asserted to optimize
the constant; the structural
taxicab-switching theorem is independent
of choosing a special example.
\end{corollary}
\begin{proof}
Each identity and its inequalities
$\frac32A\le C<B/2$
are verified by integer multiplication,
so Theorem~\ref{thm:taxicab-switching}
applies. In the second case
$A=970299$, $B=5088448$ and
$B-A=4118149$, producing
$A\log2/[2(B-A)]$ as displayed.
\end{proof}

The construction distinguishes the
two entropy sources in the cube-sum graph.
Reflection-only descendants of finitely
many fixed seeds have zero entropy
per site by
Corollary~\ref{cor:entropy-separation}.
Yet globally there are both ordinary
permutations \emph{and involutions}
with positive upper exponential entropy,
by Theorems~\ref{thm:three-cycle}
and~\ref{thm:taxicab-switching}.
The latter needs only a four-term
equal-power identity; it does not
assume any eventual-existence threshold.

\begin{corollary}[Linear entropy with parity-minimal fixed points]
\label{cor:taxicab-minimal-fixed}
Let $J_0(n)$ count legal fixed-point-free
cube-sum involutions on $[1,n]$ for
even $n$, and let $J_1(n)$ count
legal cube-sum involutions with
exactly one fixed point for odd $n$.
For every \emph{odd} positive integer
$u$, with $T=u^3$,
\[
\boxed{\begin{aligned}
J_0(54406T)&\ge2^{(4913T-1)/2},\\
J_1(4118149T)&\ge2^{(970299T-1)/2}.
\end{aligned}}
\]
In particular, both parity classes
admit explicit infinite subsequences
on which the number of legal
\emph{minimally fixed} involutions
grows exponentially in their lengths:
\[
\boxed{\begin{aligned}
\limsup_{\substack{n\to\infty\\n\ {\rm even}}}
 \frac{\log J_0(n)}n
 &\ge\frac{4913\log2}{108812}>0,\\
\limsup_{\substack{n\to\infty\\n\ {\rm odd}}}
 \frac{\log J_1(n)}n
 &\ge\frac{970299\log2}{8236298}>0.
\end{aligned}}
\]
\end{corollary}
\begin{proof}
Use the two exact identities
from Corollary~\ref{cor:taxicab-cubes}.
In Theorem~\ref{thm:taxicab-switching},
the initial two-block permutation
has one fixed point in its first
block if and only if $AT$ is even,
and one in its second block if
and only if $BT$ is even.
No taxicab switch can create or
remove a fixed point: its four
distinct endpoints occur in
two transpositions before and
after switching, while all
unselected positions are unchanged.

For $(a,b)=(17,39)$ and odd $u$,
both $AT=(au)^3$ and
$BT=(bu)^3$ are odd.
Thus the baseline and every
switched involution are
fixed-point-free, proving
the even-length bound.

For $(a,b)=(99,172)$ and odd $u$,
the first cube sum $AT$ is odd,
while $BT$ is even, so the
second reflection block has
exactly one fixed point
and the first has none.
All switches preserve this
single fixed point.
The lengths are odd in this
case, so one is the minimum
possible fixed-point count.
The entropy assertions follow
on letting odd $u\to\infty$.
\end{proof}

\section{A computation-free exponential construction}
\label{sec:three-cycle}
A particularly simple weighted identity of four cubes
produces more permutations than all the previous
asymptotic constructions, with \emph{no} finite
matching certificate or eventual-existence theorem.

\begin{proposition}[Weighted-cube odd-cycle principle]
\label{prop:odd-cycle}
Let $q\ge1$, and let $A,B,X,Y$ be positive
cubes such that
\[
(q+1)X+qY=A+2qB.
\]
For $0\le j\le2q$ put
\[
D_j=\left\lceil j/2\right\rceil(B-X)
+\left\lfloor j/2\right\rfloor(B-Y).
\]
Assume $A\le D_j\le B-2A$ for $j\ge1$
and $|D_j-D_k|\ge A$ whenever $j\ne k$.
Then for every integer $u\ge1$,
\[
\boxed{P((B-A)u^3)\ge2^{Au^3-1}.}
\]
In fact the family encodes $Au^3-1$ independent
binary choices, each of which changes precisely
$2q+1$ permutation positions.
\end{proposition}
\begin{proof}
Put $T=u^3$ and define the two-interval
reflection $p(i)=AT-i$ for $i<AT$ and
$p(i)=BT-i$ for $AT\le i\le(B-A)T$.
The interval hypotheses make this a valid
permutation with cube sums $AT$ and $BT$.
For each $1\le s<AT$, let
$v_j=s+D_jT$ for $0\le j\le2q$.
The gap conditions place $v_0$ in the first
reflection block and the other vertices in
the second, while ensuring that the translated
cycles for different $s$ are disjoint.
The edge sums $v_j+p(v_{j+1})$ are
alternately $XT$ and $YT$ for $j<2q$.
The final edge is a cube as well, since
\[
v_{2q}+p(v_0)
=(A+2qB-qX-qY)T=XT.
\]
Thus the $AT-1$ cycles can be switched
independently, proving the count and
Hamming-distance statement.
\end{proof}

For $q=1$ this recovers the three-cycle
identity $2X+Y=A+2B$. For $q=2$ it
becomes $3X+2Y=A+4B$. In particular,
\[
3\cdot5^3+2\cdot8^3=3^3+4\cdot7^3
\]
gives $(A,B,X,Y)=(27,343,125,512)$
and offsets $(D_0,\ldots,D_4)
=(0,218,49,267,98)$, satisfying every
hypothesis. Consequently
\[
\boxed{P(316u^3)\ge2^{27u^3-1}}
\]
for every positive integer $u$.
This unifies the three- and five-cycle
constructions through one algebraic identity.

\begin{lemma}[Three-cycle identity gadget]
\label{lem:three-cycle}
Let $A,B,X,Y$ be positive perfect cubes satisfying
\[
A+2B=2X+Y,\qquad
A+\tfrac12B\le X\le B-A.
\]
For every positive integer $u$, put $T=u^3$
and $N=(B-A)T$. There are at least
\[
\boxed{2^{AT-1}}
\]
different permutations $p$ of $[1,N]$ for which
$i+p(i)$ is a perfect cube.
\end{lemma}
\begin{proof}
The displayed inequalities give $B\ge4A$.
Define a reference permutation by two reflections:
\[
p(i)=
\begin{cases}
AT-i,&1\le i<AT,\\
BT-i,&AT\le i\le (B-A)T.
\end{cases}
\]
The first interval maps to itself under the
first reflection, and the second also maps to
itself under the second reflection.
Every original sum is $AT$ or $BT$,
both perfect cubes.

Set $d_1=B-X$ and $d_2=X-A$.
The hypotheses imply
\[
d_1\ge A,\qquad
d_2-d_1\ge A,\qquad
d_2+A=X\le B-A.
\]
For every $s=1,\ldots,AT-1$, the triple
\[
C_s=(s,\ s+d_1T,\ s+d_2T)
\]
lies inside $[1,N]$, with the second
and third vertices in the second reflection
interval. The three translated ranges are
pairwise disjoint by the displayed gap
inequalities. Thus all $AT-1$ triples are
pairwise vertex-disjoint.

Along the directed three-cycle
$s\to s+d_1T\to s+d_2T\to s$,
the three alternative matching sums are
\[
\begin{aligned}
s+p(s+d_1T)&=XT,\\
(s+d_1T)+p(s+d_2T)&=(A+2B-2X)T=YT,\\
(s+d_2T)+p(s)&=XT.
\end{aligned}
\]
Because $T$ and $A,B,X,Y$ are cubes,
every new sum is a cube. Switching
any subset of the $AT-1$ disjoint
cycles yields $2^{AT-1}$ pairwise
different legal permutations.
\end{proof}

\begin{theorem}[Explicit high-entropy cube-sum permutations]
\label{thm:three-cycle}
For every integer $u\ge1$,
\begin{equation}\label{eq:three-cycle}
\boxed{\displaystyle
P(681707u^3)\ge2^{148877u^3-1}.}
\end{equation}
Consequently,
\begin{equation}\label{eq:three-cycle-rate}
\boxed{\displaystyle
\limsup_{n\to\infty}\frac{\log P(n)}n
\ge \frac{148877\log2}{681707}>0.}
\end{equation}
This theorem is entirely computation-free,
apart from elementary multiplication of
four displayed integers.
\end{theorem}
\begin{proof}
Apply Lemma~\ref{lem:three-cycle} with
\[
\begin{array}{ll}
A=53^3=148877,&B=94^3=830584,\\
X=87^3=658503,&Y=79^3=493039.
\end{array}
\]
Indeed,
\[
A+2B=1810045=2X+Y,
\]
and
\[
A+\tfrac12B=564169
\le X=658503\le681707=B-A.
\]
Hence the $AT-1=148877u^3-1$
explicit cycles are disjoint. More precisely,
each has the form
\[
s\longrightarrow s+172081u^3
\longrightarrow s+509626u^3
\longrightarrow s,
\]
with alternative sums $(87u)^3,(79u)^3,(87u)^3$.
No matching search, graph decomposition,
or auxiliary finite base case is involved.
Taking logarithms along the sequence
$n=681707u^3$ proves the second claim.
\end{proof}

\begin{corollary}[Linear-distance binary permutation codes]
\label{cor:three-cycle-codes}
For $N=681707u^3$ there is an injective
map from $\{0,1\}^{148877u^3-1}$
into legal cube-sum permutations of $[1,N]$
that multiplies all Hamming distances by
exactly three. Its distance enumerator is
\[
\bigl(1+z^3\bigr)^{148877u^3-1}.
\]
In particular, a greedy binary code
subfamily has exponentially many legal
permutations with pairwise Hamming
distance proportional to $N$.
\end{corollary}
\begin{proof}
Each bit chooses whether to switch one
of the pairwise disjoint three-cycles.
A changed bit changes precisely three
mapped positions. The distance identity
and polynomial follow immediately.
The greedy Hamming-ball argument from
Corollary~\ref{cor:goodcode} applies
unchanged.
\end{proof}

\section{Power-block lifting and exponential growth on a subsequence}
\label{sec:lifting}

The nine disjoint alternating cycles at $236$
can be amplified by a purely algebraic interval
construction, yielding exponentially many solutions
at infinitely many \emph{different} lengths.
The lifting mechanism works for every perfect-power
exponent, not only cubes.

\begin{lemma}[Independent-fiber perfect-power tensorization]
\label{lem:block}
Fix $d\ge2$ and $M\ge1$, and suppose $P_d(M)>0$.
For any odd positive integer $t$ put
$L=t^d$ and $A=(L-1)/2$.
For each $r\in[1,L]$ choose independently a
$d$th-power-sum permutation $p_r$ of $[1,M]$.
Define a permutation $Q$ on the interval
$I=[A+1,A+ML]$ by
\begin{equation}\label{eq:lift}
Q\bigl(A+(i-1)L+r\bigr)
=A+(p_r(i)-1)L+(L+1-r).
\end{equation}
Then every $x+Q(x)$ is a perfect $d$th power,
and the $L$ independent choices give exactly
\[
\boxed{P_d(M)^L}
\]
different permutations of $I$.

If instead all fibers use the same $p$ possessing
$h$ pairwise vertex-disjoint directed alternating
cycles, then $Q$ has at least $hL$ pairwise
vertex-disjoint directed alternating cycles.
\end{lemma}
\begin{proof}
The input positions have unique coordinates $(i,r)$.
For fixed $r$, the map $i\mapsto p_r(i)$ is a
permutation of $[1,M]$, while the output
residue $L+1-r$ identifies the original fiber
$r$ uniquely. Thus \eqref{eq:lift} is a bijection.
For $x=A+(i-1)L+r$,
\[
\begin{aligned}
x+Q(x)
 &=2A+L(i+p_r(i)-1)+1\\
 &=L(i+p_r(i)).
\end{aligned}
\]
If $i+p_r(i)=k_{i,r}^d$, then this sum equals
$(t k_{i,r})^d$.
The map from the tuple $(p_1,\ldots,p_L)$ to $Q$
is injective, since all $p_r(i)$ are recovered
from the output block indices; hence there are
exactly $P_d(M)^L$ such interval permutations.

Finally, when every $p_r=p$, an alternating
edge $i\to j$ of $p$ gives
$x_{i,r}+Q(x_{j,r})=L(i+p(j))$ for each
$x_{i,r}=A+(i-1)L+r$, again a perfect power.
Every original alternating cycle consequently
lifts to one cycle in each fiber, and disjoint
fibers and original cycles remain disjoint.
\end{proof}

\begin{proposition}[Inverse-paired fiber involutions]
\label{prop:inverse-fibers}
In the setting of
Lemma~\ref{lem:block},
let $I_d(M)$ be the number of
valid power-sum involutions of length
$M$. The exact number of lifted
interval permutations of the form
\eqref{eq:lift} which themselves
are involutions is
\begin{equation}\label{eq:inverse-fibers}
\boxed{\displaystyle
I_d(M)\,P_d(M)^{(L-1)/2}.}
\end{equation}
\end{proposition}
\begin{proof}
Write $r'=L+1-r$. Since
\eqref{eq:lift} sends $(i,r)$
to $(p_r(i),r')$, its square sends
$(i,r)$ to
$(p_{r'}(p_r(i)),r)$.
Thus the lifted interval permutation
is an involution if and only if
\[
p_{L+1-r}=p_r^{-1}
\qquad(1\le r\le L).
\]
The inverse of any valid power-sum
permutation is valid, by symmetry
of the condition $i+p(i)$.
The $(L-1)/2$ pairs of distinct
fibers may therefore be chosen
independently in $P_d(M)$ ways
each; on the unique central fiber
$r=(L+1)/2$, the chosen
permutation must be an involution
and has $I_d(M)$ choices.
Injectivity of the fiber encoding
gives \eqref{eq:inverse-fibers}.
\end{proof}

\begin{corollary}[Transfer from general multiplicity to involution entropy]
\label{cor:inverse-transfer}
Fix $d\ge2$. Suppose valid
$d$th-power-sum involutions exist
for every sufficiently large length.
If some $M$ satisfies
$I_d(M)>0$ and $P_d(M)\ge2$,
then
\[
\boxed{\displaystyle
\limsup_{n\to\infty}
\frac{\log I_d(n)}n
\ge\frac{\log P_d(M)}{2M+1}>0.}
\]
\end{corollary}
\begin{proof}
For sufficiently large odd $t$,
put $L=t^d$, $A=(L-1)/2$
and $B=A+ML$. Choose a
valid involution on the prefix
$[1,A]$. Proposition~\ref{prop:inverse-fibers}
constructs at least
$P_d(M)^{(L-1)/2}$ valid involutions
on the middle interval.
For the suffix choose the least
perfect $d$th power
$C=k^d>(2M+1)L$, and reflect
$[B+1,C-B-1]$ through $C$.
This is a legal involution because
$B+(C-B-1)+1=C$.
The resulting full size is
\[
n=C-B-1
=\bigl(M+\tfrac12\bigr)L
 +O(t^{d-1}),
\]
since $k=(2M+1)^{1/d}t+O(1)$.
The prefix and suffix are fixed
across all middle choices.
Taking logarithms and dividing
by $n$ gives the result.
\end{proof}

This principle shows that
power-sum involution families
may have positive linear entropy
even when the independent local
choices on the middle blocks are
\emph{not themselves involutions}:
inverse-pairing complementary fibers
is enough, provided one central
involution exists.
In the cubic case the explicit
$316$-element two-reflection
construction supplies such a
central involution and several
alternative ordinary permutations.
The direct taxicab-switching theorem
in Section~\ref{sec:taxicab-involutions}
does not require this eventual
prefix-existence hypothesis.

\begin{theorem}[An exponential subsequence]
\label{thm:exponential-subsequence}
For every odd integer $t\ge9$, define
\[
\boxed{n_t=\frac{551t^3-1}{2}.}
\]
Then
\[
\boxed{P(n_t)\ge
P\!\left(\frac{t^3-1}{2}\right)512^{t^3}
\ge 3\cdot 2^{9t^3}.}
\]
In particular,
\[
\boxed{\displaystyle
\limsup_{n\to\infty}\frac{\log P(n)}n
 \ge\frac{18\log2}{551}>0.}
\]
\end{theorem}
\begin{proof}
Take the explicit $M=236$ permutation
from Section~\ref{sec:236}, possessing $h=9$
disjoint alternating $5$-cycles.
Let $L=t^3$, $A=(L-1)/2$, and
$B=A+236L=(473L-1)/2$.
Because $t\ge9$, $A\ge364\ge173$,
so Theorem~\ref{thm:exist} supplies a valid
permutation on the prefix $[1,A]$.
By Theorem~\ref{thm:multi}, $P(A)\ge3$.
The exact count $P(236)=512$ from
Theorem~\ref{thm:236} and
Lemma~\ref{lem:block} provide $512^L$
distinct valid middle-block permutations
on $[A+1,B]$, without needing to select
a preferred alternating-cycle orientation.

Now let $C=(8t)^3=512L$ and set
\[
n_t=C-B-1=\frac{551L-1}{2}.
\]
The final suffix $[B+1,n_t]$ has exactly
$39L>0$ elements. Define the permutation
there by $x\mapsto C-x$.
It permutes that suffix because
$C=B+n_t+1$, and every new sum is the cube
$C$. Combining the three disjoint intervals
produces a cube-sum permutation on $[1,n_t]$.
The prefix and middle-block choices are
independent, while the suffix is fixed.
Consequently at least $P(A)\,512^L$
distinct permutations are obtained, and
$P(A)\ge3$ completes the proof.
The older $9L$ independent-cycle construction
is a special subfamily of the new tensorized
construction.
\end{proof}

\begin{theorem}[Asymptotically optimal terminal reflection]
\label{thm:optimized-entropy}
Let $t\ge9$ be odd and $L=t^3$. Let $k_t$
be the least integer satisfying
\[
k_t^3>473L,
\]
and define
\[
\boxed{N_t=k_t^3-\frac{473L+1}{2}.}
\]
Then
\[
\boxed{P(N_t)\ge
P\!\left(\frac{L-1}{2}\right)512^L
\ge3\cdot2^{9L}.}
\]
Furthermore,
\[
N_t=\left(\frac{473}{2}+o(1)\right)L
\qquad(t\to\infty),
\]
and therefore
\begin{equation}\label{eq:optimized-entropy}
\boxed{\displaystyle
\limsup_{n\to\infty}\frac{\log P(n)}n
\ge\frac{18\log2}{473}.}
\end{equation}
\end{theorem}
\begin{proof}
Take the $236$-element tensorized gadget
of Lemma~\ref{lem:block}.
The valid prefix has length $A=(L-1)/2$
and the lifted middle block ends at
$B=(473L-1)/2$.
The strict inequality $k_t^3>473L=2B+1$
guarantees the suffix interval
$[B+1,N_t]$ is nonempty: indeed
\[
N_t-B=k_t^3-473L>0.
\]
Reflection through the cube $k_t^3$
permutes that suffix since
$B+N_t+1=k_t^3$.
The middle blocks have $512^L$ independent
choices, the prefix has $P(A)$ choices,
and the suffix is fixed.
The first bound follows, and $P(A)\ge3$
because $A\ge364$.
By minimality,
\[
(k_t-1)^3\le473t^3<k_t^3,
\]
so $k_t=473^{1/3}t+O(1)$.
Thus $k_t^3=473t^3+O(t^2)$ and
$N_t=(473/2)t^3+O(t^2)$.
Dividing $9L\log2$ by $N_t$
establishes \eqref{eq:optimized-entropy}.
\end{proof}

The fixed terminal cube $(8t)^3$ used
in Theorem~\ref{thm:exponential-subsequence}
is convenient for a simple closed formula
$n_t=(551t^3-1)/2$, but is not
asymptotically optimal. Selecting the
\emph{least} admissible terminal cube
as above shrinks the suffix from
$39t^3$ to $O(t^2)$, and improves the
linear entropy coefficient from
$18\log2/551$ to $18\log2/473$.

\begin{corollary}[Multiplicity forces exponential subsequences]
\label{cor:amplification}
Fix $d\ge2$. Suppose $P_d(n)>0$ for all
sufficiently large $n$ and that $P_d(M)\ge2$
for some $M\ge1$. Choose any integer $b$
with $b^d>2M+1$. For all sufficiently large
odd $t$, define
\[
n_t=\Bigl(b^d-M-\tfrac12\Bigr)t^d-\tfrac12.
\]
Then
\[
P_d(n_t)\ge
P_d\!\left(\frac{t^d-1}{2}\right)
P_d(M)^{t^d},
\]
and in particular
\[
\boxed{\displaystyle
\limsup_{n\to\infty}\frac{\log P_d(n)}n
\ge
\frac{\log P_d(M)}
{b^d-M-\frac12}>0.}
\]
\end{corollary}
\begin{proof}
Let $L=t^d$, $A=(L-1)/2$ and $B=A+ML$.
Eventual existence supplies a valid prefix
on $[1,A]$ for all sufficiently large odd $t$,
with $P_d(A)$ possible choices.
Lemma~\ref{lem:block} gives
$P_d(M)^L$ independent middle-block choices.
Now put $C=(bt)^d=b^dL$ and
$n_t=C-B-1$. Since $b^d>2M+1$,
the suffix $[B+1,n_t]$ is nonempty;
reflection $x\mapsto C-x$ permutes it
and makes every sum equal to $C$.
The three intervals are invariant under
the constructed permutation, and their
choices are independent. Taking logarithms
and dividing by $n_t$ gives the bound.
\end{proof}

\begin{theorem}[An exact binary cube-sum permutation code]
\label{thm:code}
For each odd $t\ge9$, put $L=t^3$ and
$n_t=(551L-1)/2$.
There is an injective map
\[
\Phi_t:\{0,1\}^{9L}\longrightarrow
\{\text{cube-sum permutations of }[1,n_t]\}
\]
such that for all bit strings $u,v$,
\begin{equation}\label{eq:isometric}
\boxed{d_{\rm H}(\Phi_t(u),\Phi_t(v))
       =5\,d_{\rm H}(u,v).}
\end{equation}
In particular, relative to $\Phi_t(0)$, the
exact Hamming-distance enumerator of this
explicit family is
\[
\boxed{\sum_{u\in\{0,1\}^{9L}}
 z^{d_{\rm H}(\Phi_t(u),\Phi_t(0))}
 =(1+z^5)^{9L}.}
\]
\end{theorem}
\begin{proof}
The proof of Theorem~\ref{thm:236}
identifies every $236$-element cube-sum
permutation with a unique nine-bit vector
recording which of the nine disjoint directed
five-cycles have been switched.
Two different choices on one five-cycle
change all five mapped values, while
all untouched components agree.
Thus their permutation Hamming distance
is exactly five times the bitwise Hamming
distance.

Take the $L$ independent base-permutation
fibers of Lemma~\ref{lem:block}, and concatenate
their nine-bit selectors into $u\in\{0,1\}^{9L}$.
Different fibers occupy disjoint input positions,
so distances add exactly.
Finally append a fixed valid prefix and the
fixed cube-reflection suffix of
Theorem~\ref{thm:exponential-subsequence}.
Neither changes the Hamming distances.
The generating polynomial follows because
each of the $9L$ bits independently contributes
either distance zero or distance five.
\end{proof}

\begin{corollary}[Positive-rate, linear-distance families]
\label{cor:goodcode}
Let $H_2(x)=-x\log_2x-(1-x)\log_2(1-x)$
be binary entropy. For every odd $t\ge9$,
there is a subfamily $\mathcal F_t$ of cube-sum
permutations on $[1,n_t]$ such that
\[
|\mathcal F_t|
\ge 2^{\,9L(1-H_2(1/4))},
\qquad
d_{\min}(\mathcal F_t)>
\frac{45L}{4}.
\]
Consequently both
$\log|\mathcal F_t|=\Omega(n_t)$ and
$d_{\min}(\mathcal F_t)=\Omega(n_t)$.
\end{corollary}
\begin{proof}
Greedily select a binary word of length
$N=9L$, and delete all words within
Hamming distance $\lfloor N/4\rfloor$
of the selected word.
A binary Hamming ball of this radius has
at most $2^{NH_2(1/4)}$ members, so the
greedy packing yields at least
$2^{N(1-H_2(1/4))}$ codewords, with
pairwise binary distance exceeding $N/4$.
Applying $\Phi_t$ multiplies all distances
by exactly five, while preserving cardinality.
Since $n_t=(551L-1)/2$, both claimed
linear asymptotics follow.
\end{proof}

\begin{corollary}[Optimal general terminal choice]
\label{cor:optimal-general}
Under the assumptions of
Corollary~\ref{cor:amplification},
the stronger universal estimate holds:
\[
\boxed{\displaystyle
\limsup_{n\to\infty}
\frac{\log P_d(n)}n
\ge\frac{\log P_d(M)}{M+\frac12}.}
\]
\end{corollary}
\begin{proof}
Put $L=t^d$ with odd $t$ and
$B=((2M+1)L-1)/2$. Take the least
integer $k$ for which
$k^d>(2M+1)L$. Since
$k=(2M+1)^{1/d}t+O(1)$, the
reflected suffix gives a terminal
size $n=k^d-B-1$ with
$n=(M+\tfrac12)L+O(t^{d-1})$.
The independent-fiber middle and
eventually valid prefix yield
$P_d(n)\ge P_d(M)^L$.
Taking logarithms proves the bound.
\end{proof}

\begin{theorem}[Power-divisor tensor recurrence]\label{thm:divisor-recurrence}
Fix an exponent $d\ge2$ and a nonnegative
integer $n$. For any odd positive integer $t$
such that $L=t^d$ divides $2n+1$, set
\[
A=\frac{L-1}{2},\qquad
q=\frac{2n+1}{L},
\qquad M_b=b^d-\frac{q+1}{2}.
\]
With $P_d(0)=1$, we have
\begin{equation}\label{eq:divisor-recurrence}
\boxed{\displaystyle
P_d(n)\ge P_d(A)
\!\!\!\sum_{\substack{b\ge2,\ M_b\ge0\\b^d<q}}
P_d(M_b)^{\,L}.}
\end{equation}
For $t=1$, this is exactly the earlier
one-step interval-reflection counting
recurrence (now also allowing the empty prefix).
\end{theorem}
\begin{proof}
For each summand choose one permutation on
$[1,A]$ and, independently, $L$ permutations
of $[1,M_b]$. Apply
Lemma~\ref{lem:block} on
$[A+1,B]$, where $B=A+M_bL$.
The equation defining $M_b$ gives
\[
B+n+1=b^dL=(bt)^d,\qquad
n-B=(q-b^d)L>0.
\]
Thus the final interval $[B+1,n]$ is
permuted by the reflection
$x\mapsto(bt)^d-x$. All three pieces
are valid and disjoint, producing
$P_d(A)P_d(M_b)^L$ distinct permutations.
Different $b$ have different last
entries $p(n)=(bt)^d-n$, so these
families are disjoint and can be added.
When $t=1$, $L=1$, $A=0$ and
$M_b=b^d-n-1$, recovering the reflection
recurrence.
\end{proof}

This formula gives a precise arithmetic
description of the advantage and limitation
of the tensor construction: a nontrivial
block scale $t>1$ requires the cubic
divisibility condition $t^d\mid(2n+1)$.
It does not, by itself, prove exponential
growth at \emph{every} length.

\begin{proposition}[Arithmetic sparsity of nontrivial block lifting]
\label{prop:cubefree}
The set of positive integers $n$ for which
some odd integer $t\ge3$ satisfies
$t^3\mid(2n+1)$ has natural density
\[
\boxed{\displaystyle
1-\frac{8}{7\zeta(3)}
=0.0492487\ldots.}
\]
Thus more than $95\%$ of lengths have no
nontrivial integral cube scale for the
particular centered tensor construction
used in Theorem~\ref{thm:divisor-recurrence}.
\end{proposition}
\begin{proof}
An odd integer is cube-free precisely
when it is divisible by no prime cube
$p^3$, with $p$ odd. By
Möbius inclusion--exclusion, the number
of positive $n\le X$ for which
$2n+1$ is cube-free equals
\[
\begin{aligned}
\sum_{\substack{n\le X\\ 2n+1\ {\rm cube\mbox{-}free}}}1
&=\sum_{\substack{d\le(2X+1)^{1/3}\\d\ {\rm odd}}}
 \mu(d)\#\{n\le X:d^3\mid2n+1\}\\
&=X\!\!\sum_{\substack{d\ge1\\d\ {\rm odd}}}
  \frac{\mu(d)}{d^3}
  +O(X^{1/3})\\
&=X\prod_{p\ {\rm odd}}(1-p^{-3})
  +O(X^{1/3})\\
&=\frac{8X}{7\zeta(3)}+O(X^{1/3}).
\end{aligned}
\]
Every odd composite $t\ge3$ has an odd
prime divisor, so $t^3\mid2n+1$
is possible for some $t\ge3$ exactly
when $2n+1$ is not cube-free.
Taking complements proves the claim.
\end{proof}

This is a limitation of the \emph{specific
centered tensor construction}, not an
upper bound on $P(n)$ or a no-go
theorem for other embeddings.
In particular, the exponential-in-$n$
subsequence result cannot by itself
be promoted to a uniform exponential
lower bound using the same integral
scales for every $n$.

\begin{lemma}[A denser finite alternating-cycle gadget]
\label{lem:362}
There exists a cube-sum permutation of
$[1,362]$ admitting at least $40$ pairwise
vertex-disjoint alternating cycles.
Consequently,
\[
\boxed{P(362)\ge 2^{40}.}
\]
\end{lemma}
\begin{proof}
An explicit certificate in
\texttt{certificates/forty\_cycles\_362.json}
describes a valid permutation using
$43$ integer interval-reflection rows
$p(i)=c-i$ and lists $40$ alternating
directed cycles on $182$ distinct vertices.
Each listed edge is checked by the
single integer identity $i+p(j)=k^3$.
These cycles are vertex-disjoint and can
be switched independently. The standalone
\texttt{verify\_362.py} checker directly
verifies every row and edge without any
matching solver or permanent algorithm.
\end{proof}

\begin{theorem}[Improved explicit exponential subsequence]
\label{thm:362subsequence}
For every odd integer $t\ge7$, let
\[
N_t=\frac{733t^3-1}{2}.
\]
Then
\[
\boxed{\displaystyle
P(N_t)\ge
P\!\left(\frac{t^3-1}{2}\right)
2^{40t^3}\ge2^{40t^3}.}
\]
In particular,
\[
\boxed{\displaystyle
\limsup_{n\to\infty}\frac{\log P(n)}n
\ge\frac{80\log2}{733}
>0.0756.}
\]
\end{theorem}
\begin{proof}
In Theorem~\ref{thm:divisor-recurrence}
take $d=3$, $M_b=362$, $b=9$.
Then $b^3=729$ and $q=2\cdot729-2\cdot362-1=733$,
so $b^3<q$ and $n=N_t$.
At $t=7$ the prefix length
$A=(7^3-1)/2=171$ is a known positive
index from the finite classification.
For every larger odd $t$, $A\ge364$,
so the existence theorem applies.
Hence $P(A)>0$ and
\[
P(N_t)\ge P(A)P(362)^{t^3}
\ge2^{40t^3}.
\]
Since $N_t\sim(733/2)t^3$,
division by $N_t$ gives the stated
exponential coefficient.
\end{proof}

The constant here is stronger than
the earlier $236$-seed coefficient
$18\log2/551$. It is a rigorously
certified \emph{lower bound}, not a
claim of optimal asymptotic entropy.

This gives an \emph{exponential limsup},
not an exponential lower bound for all $n$.
In particular, it is logically distinct from
the uniform, near-critical reflection
estimates in Section~\ref{sec:entropy}.

\section{Many permutations: an exponential cube-root bound}
Besides qualitative multiplicity, the reflection family
yields a strong quantitative statement.
Let $K(n)=\lfloor\sqrt[3]{2n}\rfloor$, with the floor
defined using integer comparisons (not floating point).

\begin{lemma}[Three disjoint embeddings]\label{lem:triple}
If $K(n)=k\ge27$, each of
$k^3,(k-1)^3,(k-2)^3$ gives a reflection extension
from a predecessor $m\ge173$ with $m<n$,
and all three predecessors obey $K(m)\ge k-6$.
\end{lemma}
\begin{proof}
By definition, $k^3\le2n<(k+1)^3$.
For the smallest of these cubes, set
$m=(k-2)^3-n-1$. Then
\[
2m>2(k-2)^3-(k+1)^3-2
=k^3-15k^2+21k-19.
\]
The right side exceeds $(k-6)^3$ for $k\ge27$, since
the difference is
\[
3k^2-87k+197>0
\quad (k\ge27).
\]
Thus $2m>(k-6)^3\ge21^3$, which implies
$m\ge173$ and $K(m)\ge k-6$.
Also $(k-2)^3<k^3\le2n$, so $m<n$.
The other two cubes give larger predecessors
and satisfy the same bounds.
Distinct cubes yield disjoint matching families
by Lemma~\ref{lem:reflect}.
\end{proof}

\begin{theorem}[A096904 quantitative bound]\label{thm:many}
For all integers $n\ge173$, with
$q(n)=\max\{0,\lfloor(K(n)-27)/6\rfloor\}$,
\[
\boxed{P(n)\ge3^{q(n)}}.
\]
In particular
\[
\log P(n)\ \ge\
\frac{\log3}{6}(2n)^{1/3}-O(1),
\]
so $P(n)$ grows at least exponentially in $n^{1/3}$.
\end{theorem}
\begin{proof}
If $K(n)<27$ then $q(n)=0$, and
Theorem~\ref{thm:exist} supplies $P(n)\ge1$.
For $K(n)\ge27$, Lemma~\ref{lem:triple}
provides three distinct extensions from smaller
predecessors $m_i\ge173$ with
$K(m_i)\ge K(n)-6$. Thus
$q(m_i)\ge q(n)-1$.
By strong induction and the disjointness of the
three extension images,
\[
P(n)\ge\sum_{i=1}^{3}P(m_i)
\ge3\cdot3^{\max(0,q(n)-1)}
\ge3^{q(n)}.
\]
The displayed logarithmic asymptotic follows immediately.
\end{proof}


\section{A universal reflection-entropy bound}\label{sec:entropy}
The fixed three-branch construction above is only
one instance of a general branching mechanism. In fact,
one may use \emph{any fixed number} of the largest
available cubes and obtain a sharp limiting coefficient
for this family of lower bounds.

For each integer $s\ge2$, let $\rho_s\in(1,2)$
be the unique solution of
\begin{equation}\label{eq:rho}
\sum_{h=1}^{s}\rho_s^{-h}=1,
\qquad
B_s=2s^2+4s-1.
\end{equation}
Let $K(n)=\lfloor\sqrt[3]{2n}\rfloor$ as before.

\begin{theorem}[Reflection entropy]\label{thm:entropy}
For all $s\ge2$ and $n\ge173$,
\begin{equation}\label{eq:entropy}
\boxed{\displaystyle P(n)\ge
\rho_s^{\left(K(n)-B_s\right)/2}.}
\end{equation}
In particular,
\begin{equation}\label{eq:liminf}
\boxed{\displaystyle
\liminf_{n\to\infty}
\frac{\log P(n)}{(2n)^{1/3}}
\ \ge\ \frac{\log 2}{2}.}
\end{equation}
Equivalently, for every $\varepsilon>0$ there is
a constant $C_\varepsilon$ such that
\[
P(n)\ge
\exp\left(\left(\tfrac12\log2-\varepsilon\right)
(2n)^{1/3}-C_\varepsilon\right)
\quad(n\ge173).
\]
\end{theorem}
\begin{proof}
Fix $s$ and write $B=B_s$, $\rho=\rho_s$.
We prove \eqref{eq:entropy} by strong induction.
For $K(n)<B$, the right side is at most $1$,
and $P(n)\ge1$ by Theorem~\ref{thm:exist}.

Now suppose $k=K(n)\ge B$. For each
$j=0,\dots,s-1$, take the cube
$c_j=(k-j)^3$ and put $m_j=c_j-n-1$.
Since $k^3\le2n<(k+1)^3$, the following
exact polynomial inequality holds:
\[
\begin{aligned}
2m_j-(k-2j-2)^3
&\ge 2(k-j)^3-(k+1)^3-1
 -(k-2j-2)^3\\
&=3k^2-(6j^2+24j+15)k\\
&\qquad+6j^3+24j^2+24j+6\\
&\ge0.
\end{aligned}
\]
Dividing the final polynomial by $3$, its
coefficient of $-k$ is $2j^2+8j+5$ and its
constant term is positive. Hence the last inequality
follows from $k\ge B\ge2j^2+8j+5$.
Also $k-2j-2\ge k-2s\ge8$, so
$m_j\ge256\ge173$ and
$K(m_j)\ge k-2(j+1)$.
Each $c_j\le k^3\le2n$, so $m_j<n$.

The $s$ reflection extensions are pairwise disjoint
by their distinct last images, as in
Lemma~\ref{lem:reflect}.
Using the strong induction hypothesis, we obtain
\[
\begin{aligned}
P(n)
&\ge\sum_{j=0}^{s-1}P(m_j)\\
&\ge\sum_{j=0}^{s-1}
\rho^{(K(m_j)-B)/2}\\
&\ge\rho^{(k-B)/2}
\sum_{j=0}^{s-1}\rho^{-(j+1)}
=\rho^{(k-B)/2}.
\end{aligned}
\]
This proves the first assertion. The $\rho_s$ are
strictly increasing and converge to $2$, since
their defining finite geometric sums approach
$\sum_{h\ge1}2^{-h}=1$. For each fixed $s$,
$B_s$ is a constant independent of $n$ and
$K(n)=(2n)^{1/3}+O(1)$.
Therefore the ratio on the left side of
\eqref{eq:liminf} is at least
$\tfrac12\log\rho_s$ in the limit.
Letting $s\to\infty$ proves \eqref{eq:liminf}.
\end{proof}


\begin{corollary}[Other power sums, conditional on existence]
For an integer $d\ge2$, let $P_d(n)$ count
permutations of $[1,n]$ in which every $i+p_i$
is a perfect $d$th power. If $P_d(n)>0$
for all sufficiently large $n$, then
\[
\liminf_{n\to\infty}
\frac{\log P_d(n)}{(2n)^{1/d}}
\ge\frac{\log2}{2}.
\]
\end{corollary}
\begin{proof}
For any fixed $s$ and $j=0,\ldots,s-1$, write
$k=\lfloor(2n)^{1/d}\rfloor$ and use
$c_j=(k-j)^d$, $m_j=c_j-n-1$.
The binomial theorem gives, uniformly over
these finitely many $j$,
\[
2(k-j)^d-(k+1)^d-(k-2j-2)^d
= d k^{d-1}+O_{d,s}(k^{d-2})>1
\]
for all sufficiently large $k$.
It follows that $m_j$ is a smaller positive
length in the eventual-existence range and
$\lfloor(2m_j)^{1/d}\rfloor\ge k-2j-2$.
The distinct reflection families obey the
same $s$-branch induction as
Theorem~\ref{thm:entropy}, with a suitable
constant initial offset in $k$.
The $\rho_s$ tend to $2$ as before.
\end{proof}


\begin{theorem}[A near-critical lower bound with explicit loss]
\label{thm:critical}
For $K(n)=\lfloor\sqrt[3]{2n}\rfloor$ and every $n\ge173$,
\begin{equation}\label{eq:critical}
\boxed{\displaystyle
P(n)\ge\frac{127}{K(n)}\,2^{(K(n)-127)/2}.}
\end{equation}
More generally, for every real $\varepsilon>0$
there exists $C_\varepsilon>0$ such that
\begin{equation}\label{eq:eps}
\boxed{\displaystyle
P(n)\ge C_\varepsilon\,
\frac{2^{K(n)/2}}{K(n)^\varepsilon}
\qquad(n\ge173).}
\end{equation}
\end{theorem}
\begin{proof}
Write $k=K(n)$. The exact cubic polynomial
inequality in the proof of
Theorem~\ref{thm:entropy} works whenever
\[
k\ge B_s=2s^2+4s-1,\qquad
s\ge2;
\]
it supplies $s$ \emph{disjoint} reflection
embeddings with predecessor indices
$m_h\in[173,n-1]$ and
$K(m_h)\ge k-2h$ for $h=1,\dots,s$.

First set $F(k)=2^{k/2}/k$, which is increasing
for all integer $k\ge3$. For $k\ge128$ choose
$s=\lfloor\log_2 k\rfloor$.
Then $s\ge7$ and $k\ge 2^s\ge B_s$:
the last inequality holds at $s=7$, since
$128\ge125$, and is preserved upon doubling,
because $2B_s\ge B_{s+1}$ for all $s\ge7$.
Moreover $k-2s\ge8$, so every predecessor
is in the increasing range of $F$.
By the reflection counting lemma,
\[
\begin{aligned}
\frac{\sum_{h=1}^{s}F(k-2h)}{F(k)}
&=\sum_{h=1}^{s}\frac{2^{-h}k}{k-2h}\\
&\ge\sum_{h=1}^{s}2^{-h}
  \left(1+\frac{2h}{k}\right)\\
&\ge1-2^{-s}+\frac2k>1,
\end{aligned}
\]
using $\sum_{h=1}^s h2^{-h}\ge1$ for
$s\ge2$ and $k<2^{s+1}$.
For $k<128$, Theorem~\ref{thm:exist}
ensures $P(n)\ge1$.
Because $F$ is increasing for $k\ge7$,
this implies
$P(n)\ge F(k)/F(127)$.
Strong induction using the above inequality
extends the same bound to all $k\ge128$.
Since $F(127)=2^{127/2}/127$,
this is exactly \eqref{eq:critical}.

For the more general result, fix
$\varepsilon>0$ and define
$F_\varepsilon(k)=2^{k/2}/k^\varepsilon$.
It is increasing for all sufficiently
large $k$. For sufficiently large $k$,
set $s=\lceil\log_2(k/\varepsilon)\rceil$,
so $s=O_\varepsilon(\log k)$,
$k\ge B_s$ and
$2^{-s}\le\varepsilon/k$.
Each predecessor is in the increasing
range of $F_\varepsilon$. By the
elementary Bernoulli inequality,
\[
\left(\frac{k}{k-2h}\right)^\varepsilon
\ge1+\frac{2\varepsilon h}{k}.
\]
Consequently
\[
\begin{aligned}
\frac{\sum_{h=1}^s F_\varepsilon(k-2h)}
{F_\varepsilon(k)}
&\ge 1-2^{-s}
+\frac{2\varepsilon}{k}
 \sum_{h=1}^s h2^{-h}\\
&\ge1-2^{-s}+\frac{2\varepsilon}{k}>1.
\end{aligned}
\]
Choose an initial $K_\varepsilon$
large enough for all these inequalities
to hold at $k\ge K_\varepsilon$.
By positivity, the bound
$P(n)\ge C_\varepsilon F_\varepsilon(k)$
holds for $k<K_\varepsilon$ on taking
$C_\varepsilon=2^{-K_\varepsilon/2}$.
The same strong induction proves it
for every larger $k$, yielding
\eqref{eq:eps}.
\end{proof}

These bounds approach the \emph{critical}
reflection coefficient $\frac12\log2$
with a logarithmic (indeed arbitrarily small
polynomial) loss. Their proofs use no new
computer-assisted certificate beyond the
existence theorem.

The coefficient $(\log2)/2$ is a \emph{proved lower
bound}, not an asserted exact exponential growth rate.
The argument needs only the qualitative existence
Theorem~\ref{thm:exist}; it makes no further use of
numerical matching data.

\section{An efficient constructive algorithm}
The proof is also algorithmic. Given $n>364$, choose the
\emph{least} cube $c=k^3\ge n+174$ and put $m=c-n-1$.
The cube-gap lemma ensures $173\le m<n$. Minimality gives
\[
(k-1)^3<n+174\le k^3
\quad\Longrightarrow\quad
173\le m\le3k^2-3k+173=O(n^{2/3}).
\]
Repeat until the remaining size is in the certified interval
$173\le m\le364$, then fill each recorded interval by
the explicit reflection $p_i=c-i$.

\begin{proposition}[Linear-output construction]
A valid cube-sum permutation of any $n\ge173$ can be
constructed in $O(n)$ unit-cost arithmetic operations,
with $O(\log\log n)$ reflection layers.
\end{proposition}
\begin{proof}
The predecessor bound $m=O(n^{2/3})$ implies
$\log m\le(2/3)\log n+O(1)$, so the number of
iterations before the fixed-size base is
$O(\log\log n)$. The returned intervals are disjoint
and their union, together with the base interval,
is $\{1,\dots,n\}$. Therefore exactly $n$ entries
are written. Each exact integer cube root can
be computed using $O(\log n)$ comparisons,
which contributes $O(\log n\log\log n)=o(n)$
additional unit-cost operations.
\end{proof}
The accompanying \texttt{construct\_linear.py}
implements this procedure with exact integer
cube roots, and can also display the recursion
plan without allocating the full output.

\section{Reproduction and limitations}
All programs require only Python 3 standard-library modules:
\begin{verbatim}
python papers/A006063/verify.py
python papers/A006063/independent_match.py
python papers/A006063/verify_extended.py
python papers/A006063/construct_linear.py
python papers/A006063/verify_word_paths.py
python papers/A006063/count_236.py
\end{verbatim}
The archived data consist of $109$ core witnesses,
the legacy $192$-row all-base cross-check, $48$ small
positive permutations, $124$ negative Hall sets,
and compact pairs of alternating-cycle certificates.
The infinite inductive and counting lemmas are
mathematical arguments rather than extrapolations
from numerical experiments.

The status is a \emph{computer-assisted proof candidate}
pending independent scrutiny of the finite certificate,
its interpretation, and literature priority.
The present manuscript does not compute the exact values
of $P(n)$ in general, nor does it claim external peer review.


\clearpage
\appendix

\clearpage
\section{The 57 algebraic cubic-order shapes}
\label{app:algebraic-shapes}
For each target $n$, the line gives a cube-root
word $k_1,\ldots,k_t$ and an output-block order
$\pi$. Its block lengths are not assumed or
listed: they are the integer determinants
in Theorem~\ref{thm:unimodular-block-factory}.
The target length is a checksum for their sum.
All listed lengths are strictly positive.
\begin{multicols}{2}
\scriptsize
\setlength{\parindent}{0pt}
\setlength{\parskip}{2pt}
% 57 finite algebraic shapes: each line specifies cubic roots and output order, no endpoint table.
% n is only a checksum: the determinant length formulas derive n as their sum.
% Read alongside UNIMODULAR_BLOCK_FACTORY.md and verify_unimodular_factory.py.
\noindent \textbf{82}: roots \texttt{4,\allowbreak 2,\allowbreak 3,\allowbreak 4,\allowbreak 5,\allowbreak 4,\allowbreak 5}; order \texttt{6,\allowbreak 2,\allowbreak 3,\allowbreak 4,\allowbreak 7,\allowbreak 1,\allowbreak 5}\par
\noindent \textbf{116}: roots \texttt{4,\allowbreak 3,\allowbreak 2,\allowbreak 3,\allowbreak 5,\allowbreak 3,\allowbreak 5,\allowbreak 3,\allowbreak 5,\allowbreak 3,\allowbreak 5,\allowbreak 3,\allowbreak 5,\allowbreak 3,\allowbreak 5,\allowbreak 4,\allowbreak 5,\allowbreak 6,\allowbreak 5,\allowbreak 6,\allowbreak 5,\allowbreak 6,\allowbreak 5,\allowbreak 6,\allowbreak 5,\allowbreak 6,\allowbreak 5}; order \texttt{16,\allowbreak 14,\allowbreak 3,\allowbreak 12,\allowbreak 27,\allowbreak 10,\allowbreak 25,\allowbreak 8,\allowbreak 23,\allowbreak 6,\allowbreak 21,\allowbreak 4,\allowbreak 19,\allowbreak 2,\allowbreak 17,\allowbreak 1,\allowbreak 15,\allowbreak 26,\allowbreak 13,\allowbreak 24,\allowbreak 11,\allowbreak 22,\allowbreak 9,\allowbreak 20,\allowbreak 7,\allowbreak 18,\allowbreak 5}\par
\noindent \textbf{118}: roots \texttt{3,\allowbreak 2,\allowbreak 3,\allowbreak 2,\allowbreak 3,\allowbreak 2,\allowbreak 5,\allowbreak 3,\allowbreak 5,\allowbreak 3,\allowbreak 5,\allowbreak 3,\allowbreak 5,\allowbreak 3,\allowbreak 5,\allowbreak 3,\allowbreak 5,\allowbreak 3,\allowbreak 5,\allowbreak 3,\allowbreak 5,\allowbreak 3,\allowbreak 5,\allowbreak 3,\allowbreak 5,\allowbreak 6,\allowbreak 5,\allowbreak 6,\allowbreak 5,\allowbreak 6,\allowbreak 5,\allowbreak 6,\allowbreak 5,\allowbreak 6,\allowbreak 5,\allowbreak 6,\allowbreak 5,\allowbreak 6,\allowbreak 5,\allowbreak 6,\allowbreak 5,\allowbreak 6,\allowbreak 5}; order \texttt{24,\allowbreak 6,\allowbreak 22,\allowbreak 4,\allowbreak 20,\allowbreak 2,\allowbreak 43,\allowbreak 18,\allowbreak 41,\allowbreak 16,\allowbreak 39,\allowbreak 14,\allowbreak 37,\allowbreak 12,\allowbreak 35,\allowbreak 10,\allowbreak 33,\allowbreak 8,\allowbreak 31,\allowbreak 5,\allowbreak 29,\allowbreak 3,\allowbreak 27,\allowbreak 1,\allowbreak 25,\allowbreak 42,\allowbreak 23,\allowbreak 40,\allowbreak 21,\allowbreak 38,\allowbreak 19,\allowbreak 36,\allowbreak 17,\allowbreak 34,\allowbreak 15,\allowbreak 32,\allowbreak 13,\allowbreak 30,\allowbreak 11,\allowbreak 28,\allowbreak 9,\allowbreak 26,\allowbreak 7}\par
\noindent \textbf{120}: roots \texttt{4,\allowbreak 3,\allowbreak 2,\allowbreak 5,\allowbreak 3,\allowbreak 5,\allowbreak 3,\allowbreak 5,\allowbreak 3,\allowbreak 5,\allowbreak 4,\allowbreak 5,\allowbreak 6,\allowbreak 5,\allowbreak 6,\allowbreak 5,\allowbreak 6,\allowbreak 5}; order \texttt{11,\allowbreak 9,\allowbreak 3,\allowbreak 18,\allowbreak 7,\allowbreak 16,\allowbreak 5,\allowbreak 14,\allowbreak 2,\allowbreak 12,\allowbreak 1,\allowbreak 10,\allowbreak 17,\allowbreak 8,\allowbreak 15,\allowbreak 6,\allowbreak 13,\allowbreak 4}\par
\noindent \textbf{121}: roots \texttt{4,\allowbreak 3,\allowbreak 5,\allowbreak 3,\allowbreak 5,\allowbreak 3,\allowbreak 5,\allowbreak 3,\allowbreak 5,\allowbreak 4,\allowbreak 5,\allowbreak 6,\allowbreak 5,\allowbreak 6,\allowbreak 5,\allowbreak 6,\allowbreak 5}; order \texttt{10,\allowbreak 8,\allowbreak 17,\allowbreak 6,\allowbreak 15,\allowbreak 4,\allowbreak 13,\allowbreak 2,\allowbreak 11,\allowbreak 1,\allowbreak 9,\allowbreak 16,\allowbreak 7,\allowbreak 14,\allowbreak 5,\allowbreak 12,\allowbreak 3}\par
\noindent \textbf{123}: roots \texttt{4,\allowbreak 5,\allowbreak 4,\allowbreak 5,\allowbreak 4,\allowbreak 5,\allowbreak 6,\allowbreak 5,\allowbreak 6}; order \texttt{5,\allowbreak 8,\allowbreak 3,\allowbreak 6,\allowbreak 1,\allowbreak 4,\allowbreak 9,\allowbreak 2,\allowbreak 7}\par
\noindent \textbf{125}: roots \texttt{5,\allowbreak 2,\allowbreak 5,\allowbreak 4,\allowbreak 5,\allowbreak 4,\allowbreak 5,\allowbreak 6,\allowbreak 5,\allowbreak 6,\allowbreak 5,\allowbreak 6,\allowbreak 5,\allowbreak 6}; order \texttt{13,\allowbreak 2,\allowbreak 11,\allowbreak 6,\allowbreak 9,\allowbreak 4,\allowbreak 7,\allowbreak 14,\allowbreak 5,\allowbreak 12,\allowbreak 3,\allowbreak 10,\allowbreak 1,\allowbreak 8}\par
\noindent \textbf{126}: roots \texttt{4,\allowbreak 2,\allowbreak 5,\allowbreak 4,\allowbreak 5,\allowbreak 4,\allowbreak 5,\allowbreak 6,\allowbreak 5,\allowbreak 6}; order \texttt{6,\allowbreak 2,\allowbreak 9,\allowbreak 4,\allowbreak 7,\allowbreak 1,\allowbreak 5,\allowbreak 10,\allowbreak 3,\allowbreak 8}\par
\noindent \textbf{127}: roots \texttt{5,\allowbreak 2,\allowbreak 3,\allowbreak 2,\allowbreak 5,\allowbreak 3,\allowbreak 5,\allowbreak 3,\allowbreak 5,\allowbreak 3,\allowbreak 5,\allowbreak 4,\allowbreak 5,\allowbreak 4,\allowbreak 5,\allowbreak 6,\allowbreak 5,\allowbreak 6,\allowbreak 5,\allowbreak 6,\allowbreak 5,\allowbreak 6,\allowbreak 5,\allowbreak 6,\allowbreak 5,\allowbreak 6,\allowbreak 5,\allowbreak 6}; order \texttt{10,\allowbreak 4,\allowbreak 27,\allowbreak 2,\allowbreak 25,\allowbreak 8,\allowbreak 23,\allowbreak 6,\allowbreak 21,\allowbreak 3,\allowbreak 19,\allowbreak 14,\allowbreak 17,\allowbreak 12,\allowbreak 15,\allowbreak 28,\allowbreak 13,\allowbreak 26,\allowbreak 11,\allowbreak 24,\allowbreak 9,\allowbreak 22,\allowbreak 7,\allowbreak 20,\allowbreak 5,\allowbreak 18,\allowbreak 1,\allowbreak 16}\par
\noindent \textbf{129}: roots \texttt{5,\allowbreak 2,\allowbreak 3,\allowbreak 5,\allowbreak 3,\allowbreak 5,\allowbreak 3,\allowbreak 5,\allowbreak 4,\allowbreak 5,\allowbreak 4,\allowbreak 5,\allowbreak 6,\allowbreak 5,\allowbreak 6,\allowbreak 5,\allowbreak 6,\allowbreak 5,\allowbreak 6,\allowbreak 5,\allowbreak 6,\allowbreak 5,\allowbreak 6}; order \texttt{22,\allowbreak 2,\allowbreak 7,\allowbreak 20,\allowbreak 5,\allowbreak 18,\allowbreak 3,\allowbreak 16,\allowbreak 11,\allowbreak 14,\allowbreak 9,\allowbreak 12,\allowbreak 23,\allowbreak 10,\allowbreak 21,\allowbreak 8,\allowbreak 19,\allowbreak 6,\allowbreak 17,\allowbreak 4,\allowbreak 15,\allowbreak 1,\allowbreak 13}\par
\noindent \textbf{173}: roots \texttt{2,\allowbreak 3,\allowbreak 5,\allowbreak 3,\allowbreak 5,\allowbreak 4,\allowbreak 6,\allowbreak 5,\allowbreak 6,\allowbreak 5,\allowbreak 6}; order \texttt{1,\allowbreak 4,\allowbreak 10,\allowbreak 2,\allowbreak 8,\allowbreak 6,\allowbreak 11,\allowbreak 5,\allowbreak 9,\allowbreak 3,\allowbreak 7}\par
\noindent \textbf{176}: roots \texttt{2,\allowbreak 3,\allowbreak 5,\allowbreak 4,\allowbreak 3,\allowbreak 5,\allowbreak 4,\allowbreak 6,\allowbreak 5,\allowbreak 6,\allowbreak 5,\allowbreak 6,\allowbreak 7,\allowbreak 6}; order \texttt{1,\allowbreak 5,\allowbreak 11,\allowbreak 2,\allowbreak 9,\allowbreak 7,\allowbreak 14,\allowbreak 4,\allowbreak 12,\allowbreak 6,\allowbreak 10,\allowbreak 3,\allowbreak 8,\allowbreak 13}\par
\noindent \textbf{177}: roots \texttt{5,\allowbreak 2,\allowbreak 3,\allowbreak 5,\allowbreak 3,\allowbreak 5,\allowbreak 3,\allowbreak 5,\allowbreak 4,\allowbreak 5,\allowbreak 4,\allowbreak 6,\allowbreak 5,\allowbreak 6,\allowbreak 5,\allowbreak 6,\allowbreak 5,\allowbreak 6,\allowbreak 5,\allowbreak 6,\allowbreak 5,\allowbreak 6}; order \texttt{21,\allowbreak 2,\allowbreak 7,\allowbreak 19,\allowbreak 5,\allowbreak 17,\allowbreak 3,\allowbreak 15,\allowbreak 11,\allowbreak 13,\allowbreak 9,\allowbreak 22,\allowbreak 10,\allowbreak 20,\allowbreak 8,\allowbreak 18,\allowbreak 6,\allowbreak 16,\allowbreak 4,\allowbreak 14,\allowbreak 1,\allowbreak 12}\par
\noindent \textbf{179}: roots \texttt{5,\allowbreak 3,\allowbreak 4,\allowbreak 3,\allowbreak 4,\allowbreak 5,\allowbreak 4,\allowbreak 6,\allowbreak 4,\allowbreak 6,\allowbreak 5,\allowbreak 6,\allowbreak 5,\allowbreak 6,\allowbreak 7,\allowbreak 6,\allowbreak 7}; order \texttt{13,\allowbreak 4,\allowbreak 9,\allowbreak 2,\allowbreak 7,\allowbreak 11,\allowbreak 5,\allowbreak 16,\allowbreak 3,\allowbreak 14,\allowbreak 6,\allowbreak 12,\allowbreak 1,\allowbreak 10,\allowbreak 17,\allowbreak 8,\allowbreak 15}\par
\noindent \textbf{180}: roots \texttt{3,\allowbreak 5,\allowbreak 3,\allowbreak 5,\allowbreak 3,\allowbreak 5,\allowbreak 3,\allowbreak 5,\allowbreak 4,\allowbreak 6,\allowbreak 5,\allowbreak 6,\allowbreak 5,\allowbreak 6,\allowbreak 5,\allowbreak 6,\allowbreak 5,\allowbreak 6}; order \texttt{7,\allowbreak 17,\allowbreak 5,\allowbreak 15,\allowbreak 3,\allowbreak 13,\allowbreak 1,\allowbreak 11,\allowbreak 9,\allowbreak 18,\allowbreak 8,\allowbreak 16,\allowbreak 6,\allowbreak 14,\allowbreak 4,\allowbreak 12,\allowbreak 2,\allowbreak 10}\par
\noindent \textbf{181}: roots \texttt{5,\allowbreak 2,\allowbreak 5,\allowbreak 4,\allowbreak 5,\allowbreak 4,\allowbreak 6,\allowbreak 5,\allowbreak 6,\allowbreak 5,\allowbreak 6,\allowbreak 5,\allowbreak 6}; order \texttt{12,\allowbreak 2,\allowbreak 10,\allowbreak 6,\allowbreak 8,\allowbreak 4,\allowbreak 13,\allowbreak 5,\allowbreak 11,\allowbreak 3,\allowbreak 9,\allowbreak 1,\allowbreak 7}\par
\noindent \textbf{184}: roots \texttt{2,\allowbreak 4,\allowbreak 6,\allowbreak 4,\allowbreak 6,\allowbreak 7,\allowbreak 6}; order \texttt{1,\allowbreak 4,\allowbreak 7,\allowbreak 2,\allowbreak 5,\allowbreak 6,\allowbreak 3}\par
\noindent \textbf{185}: roots \texttt{5,\allowbreak 3,\allowbreak 4,\allowbreak 6,\allowbreak 4,\allowbreak 6,\allowbreak 5,\allowbreak 6,\allowbreak 7}; order \texttt{2,\allowbreak 5,\allowbreak 7,\allowbreak 3,\allowbreak 8,\allowbreak 1,\allowbreak 6,\allowbreak 9,\allowbreak 4}\par
\noindent \textbf{188}: roots \texttt{5,\allowbreak 4,\allowbreak 2,\allowbreak 4,\allowbreak 6,\allowbreak 4,\allowbreak 6,\allowbreak 5,\allowbreak 6,\allowbreak 7,\allowbreak 6,\allowbreak 7}; order \texttt{3,\allowbreak 6,\allowbreak 8,\allowbreak 4,\allowbreak 11,\allowbreak 2,\allowbreak 9,\allowbreak 1,\allowbreak 7,\allowbreak 12,\allowbreak 5,\allowbreak 10}\par
\noindent \textbf{194}: roots \texttt{2,\allowbreak 3,\allowbreak 5,\allowbreak 3,\allowbreak 5,\allowbreak 6,\allowbreak 5,\allowbreak 6,\allowbreak 5,\allowbreak 6}; order \texttt{1,\allowbreak 4,\allowbreak 9,\allowbreak 2,\allowbreak 7,\allowbreak 10,\allowbreak 5,\allowbreak 8,\allowbreak 3,\allowbreak 6}\par
\noindent \textbf{197}: roots \texttt{5,\allowbreak 3,\allowbreak 6,\allowbreak 4,\allowbreak 6,\allowbreak 5,\allowbreak 6,\allowbreak 7,\allowbreak 6}; order \texttt{4,\allowbreak 2,\allowbreak 9,\allowbreak 6,\allowbreak 7,\allowbreak 1,\allowbreak 5,\allowbreak 8,\allowbreak 3}\par
\noindent \textbf{198}: roots \texttt{4,\allowbreak 2,\allowbreak 5,\allowbreak 4,\allowbreak 6,\allowbreak 5,\allowbreak 4,\allowbreak 5,\allowbreak 6,\allowbreak 5,\allowbreak 6,\allowbreak 7,\allowbreak 6,\allowbreak 7}; order \texttt{7,\allowbreak 2,\allowbreak 10,\allowbreak 4,\allowbreak 13,\allowbreak 8,\allowbreak 1,\allowbreak 6,\allowbreak 11,\allowbreak 3,\allowbreak 9,\allowbreak 14,\allowbreak 5,\allowbreak 12}\par
\noindent \textbf{199}: roots \texttt{4,\allowbreak 2,\allowbreak 4,\allowbreak 6,\allowbreak 4,\allowbreak 6,\allowbreak 4,\allowbreak 5,\allowbreak 4,\allowbreak 5,\allowbreak 6,\allowbreak 7,\allowbreak 6,\allowbreak 7,\allowbreak 6}; order \texttt{9,\allowbreak 2,\allowbreak 7,\allowbreak 15,\allowbreak 5,\allowbreak 13,\allowbreak 3,\allowbreak 10,\allowbreak 1,\allowbreak 8,\allowbreak 11,\allowbreak 14,\allowbreak 6,\allowbreak 12,\allowbreak 4}\par
\noindent \textbf{200}: roots \texttt{3,\allowbreak 4,\allowbreak 6,\allowbreak 5,\allowbreak 6,\allowbreak 7,\allowbreak 6,\allowbreak 7}; order \texttt{1,\allowbreak 2,\allowbreak 7,\allowbreak 4,\allowbreak 5,\allowbreak 8,\allowbreak 3,\allowbreak 6}\par
\noindent \textbf{201}: roots \texttt{3,\allowbreak 4,\allowbreak 3,\allowbreak 4,\allowbreak 6,\allowbreak 5,\allowbreak 4,\allowbreak 6,\allowbreak 7,\allowbreak 6,\allowbreak 7,\allowbreak 6}; order \texttt{7,\allowbreak 3,\allowbreak 12,\allowbreak 1,\allowbreak 4,\allowbreak 10,\allowbreak 2,\allowbreak 6,\allowbreak 8,\allowbreak 11,\allowbreak 5,\allowbreak 9}\par
\noindent \textbf{202}: roots \texttt{4,\allowbreak 6,\allowbreak 4,\allowbreak 6,\allowbreak 4,\allowbreak 6,\allowbreak 7,\allowbreak 6,\allowbreak 7,\allowbreak 6}; order \texttt{5,\allowbreak 10,\allowbreak 3,\allowbreak 8,\allowbreak 1,\allowbreak 6,\allowbreak 9,\allowbreak 4,\allowbreak 7,\allowbreak 2}\par
\noindent \textbf{203}: roots \texttt{4,\allowbreak 6,\allowbreak 3,\allowbreak 4,\allowbreak 6,\allowbreak 4,\allowbreak 6,\allowbreak 7,\allowbreak 6,\allowbreak 7,\allowbreak 6}; order \texttt{3,\allowbreak 6,\allowbreak 11,\allowbreak 4,\allowbreak 9,\allowbreak 1,\allowbreak 7,\allowbreak 10,\allowbreak 5,\allowbreak 8,\allowbreak 2}\par
\noindent \textbf{204}: roots \texttt{4,\allowbreak 6,\allowbreak 3,\allowbreak 4,\allowbreak 6,\allowbreak 5,\allowbreak 6,\allowbreak 7,\allowbreak 6,\allowbreak 7,\allowbreak 6}; order \texttt{3,\allowbreak 11,\allowbreak 4,\allowbreak 9,\allowbreak 1,\allowbreak 6,\allowbreak 7,\allowbreak 10,\allowbreak 5,\allowbreak 8,\allowbreak 2}\par
\noindent \textbf{205}: roots \texttt{4,\allowbreak 2,\allowbreak 3,\allowbreak 4,\allowbreak 6,\allowbreak 5,\allowbreak 4,\allowbreak 5,\allowbreak 6,\allowbreak 7,\allowbreak 6,\allowbreak 7}; order \texttt{7,\allowbreak 2,\allowbreak 3,\allowbreak 4,\allowbreak 11,\allowbreak 8,\allowbreak 1,\allowbreak 6,\allowbreak 9,\allowbreak 12,\allowbreak 5,\allowbreak 10}\par
\noindent \textbf{206}: roots \texttt{2,\allowbreak 3,\allowbreak 4,\allowbreak 6,\allowbreak 4,\allowbreak 6,\allowbreak 5,\allowbreak 4,\allowbreak 5,\allowbreak 6,\allowbreak 7,\allowbreak 6,\allowbreak 7,\allowbreak 6,\allowbreak 7}; order \texttt{1,\allowbreak 8,\allowbreak 2,\allowbreak 5,\allowbreak 14,\allowbreak 3,\allowbreak 9,\allowbreak 12,\allowbreak 7,\allowbreak 10,\allowbreak 15,\allowbreak 6,\allowbreak 13,\allowbreak 4,\allowbreak 11}\par
\noindent \textbf{233}: roots \texttt{2,\allowbreak 6,\allowbreak 4,\allowbreak 5,\allowbreak 6,\allowbreak 5,\allowbreak 6,\allowbreak 5,\allowbreak 6,\allowbreak 7,\allowbreak 6,\allowbreak 7,\allowbreak 6,\allowbreak 7,\allowbreak 6,\allowbreak 7}; order \texttt{1,\allowbreak 15,\allowbreak 8,\allowbreak 13,\allowbreak 3,\allowbreak 6,\allowbreak 11,\allowbreak 4,\allowbreak 9,\allowbreak 16,\allowbreak 7,\allowbreak 14,\allowbreak 5,\allowbreak 12,\allowbreak 2,\allowbreak 10}\par
\noindent \textbf{234}: roots \texttt{6,\allowbreak 5,\allowbreak 6,\allowbreak 5,\allowbreak 6,\allowbreak 5,\allowbreak 7,\allowbreak 6,\allowbreak 7,\allowbreak 6,\allowbreak 7,\allowbreak 6,\allowbreak 7}; order \texttt{12,\allowbreak 6,\allowbreak 10,\allowbreak 4,\allowbreak 8,\allowbreak 2,\allowbreak 13,\allowbreak 5,\allowbreak 11,\allowbreak 3,\allowbreak 9,\allowbreak 1,\allowbreak 7}\par
\noindent \textbf{235}: roots \texttt{2,\allowbreak 6,\allowbreak 3,\allowbreak 5,\allowbreak 6,\allowbreak 5,\allowbreak 6,\allowbreak 5,\allowbreak 7,\allowbreak 6,\allowbreak 7,\allowbreak 6,\allowbreak 7,\allowbreak 6,\allowbreak 7}; order \texttt{1,\allowbreak 14,\allowbreak 3,\allowbreak 8,\allowbreak 12,\allowbreak 6,\allowbreak 10,\allowbreak 4,\allowbreak 15,\allowbreak 7,\allowbreak 13,\allowbreak 5,\allowbreak 11,\allowbreak 2,\allowbreak 9}\par
\noindent \textbf{236}: roots \texttt{2,\allowbreak 6,\allowbreak 5,\allowbreak 6,\allowbreak 4,\allowbreak 5,\allowbreak 6,\allowbreak 5,\allowbreak 7,\allowbreak 6,\allowbreak 7,\allowbreak 6,\allowbreak 7,\allowbreak 6,\allowbreak 7}; order \texttt{1,\allowbreak 14,\allowbreak 5,\allowbreak 8,\allowbreak 12,\allowbreak 6,\allowbreak 10,\allowbreak 3,\allowbreak 15,\allowbreak 7,\allowbreak 13,\allowbreak 4,\allowbreak 11,\allowbreak 2,\allowbreak 9}\par
\noindent \textbf{243}: roots \texttt{4,\allowbreak 5,\allowbreak 4,\allowbreak 6,\allowbreak 5,\allowbreak 7,\allowbreak 6,\allowbreak 7}; order \texttt{3,\allowbreak 7,\allowbreak 5,\allowbreak 1,\allowbreak 2,\allowbreak 8,\allowbreak 4,\allowbreak 6}\par
\noindent \textbf{245}: roots \texttt{4,\allowbreak 5,\allowbreak 6,\allowbreak 5,\allowbreak 7,\allowbreak 6,\allowbreak 7}; order \texttt{6,\allowbreak 4,\allowbreak 1,\allowbreak 2,\allowbreak 7,\allowbreak 3,\allowbreak 5}\par
\noindent \textbf{246}: roots \texttt{6,\allowbreak 3,\allowbreak 6,\allowbreak 5,\allowbreak 4,\allowbreak 5,\allowbreak 6,\allowbreak 4,\allowbreak 6,\allowbreak 5,\allowbreak 7,\allowbreak 6,\allowbreak 7,\allowbreak 6,\allowbreak 7,\allowbreak 6,\allowbreak 7,\allowbreak 6,\allowbreak 7}; order \texttt{8,\allowbreak 2,\allowbreak 18,\allowbreak 5,\allowbreak 16,\allowbreak 10,\allowbreak 14,\allowbreak 6,\allowbreak 12,\allowbreak 4,\allowbreak 19,\allowbreak 9,\allowbreak 17,\allowbreak 7,\allowbreak 15,\allowbreak 3,\allowbreak 13,\allowbreak 1,\allowbreak 11}\par
\noindent \textbf{247}: roots \texttt{6,\allowbreak 2,\allowbreak 3,\allowbreak 2,\allowbreak 6,\allowbreak 3,\allowbreak 6,\allowbreak 5,\allowbreak 4,\allowbreak 5,\allowbreak 4,\allowbreak 5,\allowbreak 6,\allowbreak 4,\allowbreak 6,\allowbreak 5,\allowbreak 7,\allowbreak 6,\allowbreak 7,\allowbreak 6,\allowbreak 7,\allowbreak 6,\allowbreak 7,\allowbreak 6,\allowbreak 7,\allowbreak 6,\allowbreak 7}; order \texttt{14,\allowbreak 4,\allowbreak 6,\allowbreak 2,\allowbreak 11,\allowbreak 3,\allowbreak 26,\allowbreak 9,\allowbreak 24,\allowbreak 16,\allowbreak 22,\allowbreak 12,\allowbreak 20,\allowbreak 10,\allowbreak 18,\allowbreak 8,\allowbreak 27,\allowbreak 15,\allowbreak 25,\allowbreak 13,\allowbreak 23,\allowbreak 7,\allowbreak 21,\allowbreak 5,\allowbreak 19,\allowbreak 1,\allowbreak 17}\par
\noindent \textbf{248}: roots \texttt{6,\allowbreak 2,\allowbreak 6,\allowbreak 5,\allowbreak 4,\allowbreak 5,\allowbreak 4,\allowbreak 6,\allowbreak 4,\allowbreak 6,\allowbreak 5,\allowbreak 7,\allowbreak 6,\allowbreak 7,\allowbreak 6,\allowbreak 7,\allowbreak 6,\allowbreak 7,\allowbreak 6,\allowbreak 7}; order \texttt{9,\allowbreak 2,\allowbreak 7,\allowbreak 19,\allowbreak 5,\allowbreak 17,\allowbreak 11,\allowbreak 15,\allowbreak 6,\allowbreak 13,\allowbreak 4,\allowbreak 20,\allowbreak 10,\allowbreak 18,\allowbreak 8,\allowbreak 16,\allowbreak 3,\allowbreak 14,\allowbreak 1,\allowbreak 12}\par
\noindent \textbf{249}: roots \texttt{3,\allowbreak 2,\allowbreak 6,\allowbreak 5,\allowbreak 4,\allowbreak 6,\allowbreak 4,\allowbreak 6,\allowbreak 5,\allowbreak 7,\allowbreak 6,\allowbreak 7,\allowbreak 6,\allowbreak 7,\allowbreak 6,\allowbreak 7}; order \texttt{7,\allowbreak 2,\allowbreak 5,\allowbreak 15,\allowbreak 1,\allowbreak 13,\allowbreak 9,\allowbreak 11,\allowbreak 4,\allowbreak 16,\allowbreak 8,\allowbreak 14,\allowbreak 6,\allowbreak 12,\allowbreak 3,\allowbreak 10}\par
\noindent \textbf{250}: roots \texttt{2,\allowbreak 6,\allowbreak 5,\allowbreak 4,\allowbreak 6,\allowbreak 5,\allowbreak 7,\allowbreak 6,\allowbreak 7,\allowbreak 6,\allowbreak 7}; order \texttt{1,\allowbreak 4,\allowbreak 10,\allowbreak 6,\allowbreak 8,\allowbreak 3,\allowbreak 11,\allowbreak 5,\allowbreak 9,\allowbreak 2,\allowbreak 7}\par
\noindent \textbf{251}: roots \texttt{2,\allowbreak 3,\allowbreak 6,\allowbreak 5,\allowbreak 6,\allowbreak 5,\allowbreak 7,\allowbreak 6,\allowbreak 7,\allowbreak 6,\allowbreak 7}; order \texttt{1,\allowbreak 2,\allowbreak 10,\allowbreak 6,\allowbreak 8,\allowbreak 4,\allowbreak 11,\allowbreak 5,\allowbreak 9,\allowbreak 3,\allowbreak 7}\par
\noindent \textbf{252}: roots \texttt{2,\allowbreak 6,\allowbreak 4,\allowbreak 6,\allowbreak 5,\allowbreak 6,\allowbreak 5,\allowbreak 6,\allowbreak 5,\allowbreak 7,\allowbreak 6,\allowbreak 7,\allowbreak 6,\allowbreak 7,\allowbreak 6,\allowbreak 7,\allowbreak 6,\allowbreak 7}; order \texttt{1,\allowbreak 17,\allowbreak 9,\allowbreak 15,\allowbreak 3,\allowbreak 13,\allowbreak 7,\allowbreak 11,\allowbreak 5,\allowbreak 18,\allowbreak 8,\allowbreak 16,\allowbreak 6,\allowbreak 14,\allowbreak 4,\allowbreak 12,\allowbreak 2,\allowbreak 10}\par
\noindent \textbf{253}: roots \texttt{4,\allowbreak 3,\allowbreak 6,\allowbreak 4,\allowbreak 6,\allowbreak 4,\allowbreak 5,\allowbreak 6,\allowbreak 5,\allowbreak 6,\allowbreak 5,\allowbreak 7,\allowbreak 6,\allowbreak 7,\allowbreak 6,\allowbreak 7,\allowbreak 6,\allowbreak 7,\allowbreak 6,\allowbreak 7}; order \texttt{19,\allowbreak 2,\allowbreak 17,\allowbreak 6,\allowbreak 11,\allowbreak 4,\allowbreak 9,\allowbreak 15,\allowbreak 1,\allowbreak 13,\allowbreak 7,\allowbreak 20,\allowbreak 10,\allowbreak 18,\allowbreak 8,\allowbreak 16,\allowbreak 5,\allowbreak 14,\allowbreak 3,\allowbreak 12}\par
\noindent \textbf{339}: roots \texttt{3,\allowbreak 6,\allowbreak 4,\allowbreak 3,\allowbreak 6,\allowbreak 8,\allowbreak 7,\allowbreak 8,\allowbreak 7}; order \texttt{4,\allowbreak 9,\allowbreak 1,\allowbreak 7,\allowbreak 3,\allowbreak 5,\allowbreak 8,\allowbreak 2,\allowbreak 6}\par
\noindent \textbf{341}: roots \texttt{3,\allowbreak 7,\allowbreak 6,\allowbreak 4,\allowbreak 6,\allowbreak 7,\allowbreak 8,\allowbreak 7}; order \texttt{4,\allowbreak 8,\allowbreak 1,\allowbreak 5,\allowbreak 6,\allowbreak 3,\allowbreak 7,\allowbreak 2}\par
\noindent \textbf{343}: roots \texttt{5,\allowbreak 4,\allowbreak 6,\allowbreak 4,\allowbreak 6,\allowbreak 5,\allowbreak 6,\allowbreak 8,\allowbreak 7,\allowbreak 8}; order \texttt{9,\allowbreak 4,\allowbreak 6,\allowbreak 2,\allowbreak 7,\allowbreak 1,\allowbreak 5,\allowbreak 10,\allowbreak 3,\allowbreak 8}\par
\noindent \textbf{344}: roots \texttt{7,\allowbreak 6,\allowbreak 5,\allowbreak 6,\allowbreak 7,\allowbreak 8,\allowbreak 7,\allowbreak 6,\allowbreak 7,\allowbreak 8,\allowbreak 7,\allowbreak 8}; order \texttt{11,\allowbreak 8,\allowbreak 3,\allowbreak 4,\allowbreak 9,\allowbreak 12,\allowbreak 7,\allowbreak 2,\allowbreak 5,\allowbreak 10,\allowbreak 1,\allowbreak 6}\par
\noindent \textbf{346}: roots \texttt{5,\allowbreak 7,\allowbreak 4,\allowbreak 6,\allowbreak 5,\allowbreak 6,\allowbreak 8,\allowbreak 7,\allowbreak 6,\allowbreak 7,\allowbreak 8}; order \texttt{3,\allowbreak 9,\allowbreak 5,\allowbreak 6,\allowbreak 1,\allowbreak 10,\allowbreak 4,\allowbreak 8,\allowbreak 11,\allowbreak 2,\allowbreak 7}\par
\noindent \textbf{347}: roots \texttt{3,\allowbreak 4,\allowbreak 6,\allowbreak 4,\allowbreak 6,\allowbreak 7,\allowbreak 8,\allowbreak 7,\allowbreak 8}; order \texttt{8,\allowbreak 4,\allowbreak 1,\allowbreak 2,\allowbreak 5,\allowbreak 9,\allowbreak 3,\allowbreak 6,\allowbreak 7}\par
\noindent \textbf{348}: roots \texttt{6,\allowbreak 5,\allowbreak 6,\allowbreak 5,\allowbreak 6,\allowbreak 7,\allowbreak 6,\allowbreak 7,\allowbreak 8,\allowbreak 7,\allowbreak 8}; order \texttt{10,\allowbreak 4,\allowbreak 7,\allowbreak 2,\allowbreak 5,\allowbreak 11,\allowbreak 3,\allowbreak 8,\allowbreak 1,\allowbreak 6,\allowbreak 9}\par
\noindent \textbf{349}: roots \texttt{7,\allowbreak 4,\allowbreak 6,\allowbreak 5,\allowbreak 6,\allowbreak 8,\allowbreak 6,\allowbreak 7,\allowbreak 6,\allowbreak 8}; order \texttt{9,\allowbreak 2,\allowbreak 7,\allowbreak 4,\allowbreak 5,\allowbreak 8,\allowbreak 3,\allowbreak 10,\allowbreak 1,\allowbreak 6}\par
\noindent \textbf{350}: roots \texttt{7,\allowbreak 4,\allowbreak 3,\allowbreak 4,\allowbreak 6,\allowbreak 8,\allowbreak 6,\allowbreak 8}; order \texttt{3,\allowbreak 4,\allowbreak 7,\allowbreak 2,\allowbreak 5,\allowbreak 8,\allowbreak 1,\allowbreak 6}\par
\noindent \textbf{353}: roots \texttt{3,\allowbreak 2,\allowbreak 6,\allowbreak 3,\allowbreak 5,\allowbreak 4,\allowbreak 6,\allowbreak 7,\allowbreak 6,\allowbreak 8,\allowbreak 7,\allowbreak 8}; order \texttt{2,\allowbreak 4,\allowbreak 6,\allowbreak 1,\allowbreak 11,\allowbreak 9,\allowbreak 5,\allowbreak 7,\allowbreak 12,\allowbreak 3,\allowbreak 8,\allowbreak 10}\par
\noindent \textbf{354}: roots \texttt{7,\allowbreak 4,\allowbreak 3,\allowbreak 4,\allowbreak 6,\allowbreak 8,\allowbreak 7,\allowbreak 6,\allowbreak 8,\allowbreak 7,\allowbreak 8}; order \texttt{10,\allowbreak 3,\allowbreak 4,\allowbreak 8,\allowbreak 2,\allowbreak 5,\allowbreak 11,\allowbreak 7,\allowbreak 9,\allowbreak 1,\allowbreak 6}\par
\noindent \textbf{358}: roots \texttt{3,\allowbreak 2,\allowbreak 7,\allowbreak 4,\allowbreak 5,\allowbreak 6,\allowbreak 8,\allowbreak 6,\allowbreak 7,\allowbreak 8}; order \texttt{2,\allowbreak 4,\allowbreak 1,\allowbreak 8,\allowbreak 5,\allowbreak 6,\allowbreak 9,\allowbreak 10,\allowbreak 3,\allowbreak 7}\par
\noindent \textbf{363}: roots \texttt{7,\allowbreak 6,\allowbreak 4,\allowbreak 6,\allowbreak 5,\allowbreak 6,\allowbreak 7,\allowbreak 8,\allowbreak 7,\allowbreak 8,\allowbreak 7,\allowbreak 8}; order \texttt{11,\allowbreak 3,\allowbreak 5,\allowbreak 6,\allowbreak 12,\allowbreak 9,\allowbreak 4,\allowbreak 7,\allowbreak 2,\allowbreak 10,\allowbreak 1,\allowbreak 8}\par

\end{multicols}

\section{Redundant 57-row endpoint atlas}
\label{app:interval-atlas}

\noindent Each row $n$ lists pairs $b:k$:
the consecutive input interval ending at $b$
is mapped by $p(i)=k^3-i$.
The first interval starts at $1$.
There are $57$ exceptional rows and $795$
segments. Six parameterized block families
prove twenty-five additional base lengths. For every row,
Lemma~\ref{lem:endpoint-cancellation}
reduces checking the permutation
to equality of two lists of integers.

\begin{multicols}{2}
\scriptsize\ttfamily
\setlength{\parindent}{0pt}
\setlength{\parskip}{2pt}
% 57 exceptional cube-sum interval identities, totaling 795 reflected blocks.
% The other 25 seeds follow from six algebraic permutation block-rearrangement families.
\noindent \textbf{82}: 2:4\allowbreak\ 5:2\allowbreak\ 21:3\allowbreak\ 42:4\allowbreak\ 61:5\allowbreak\ 63:4\allowbreak\ 82:5\par
\noindent \textbf{116}: 2:4\allowbreak\ 3:3\allowbreak\ 4:2\allowbreak\ 8:3\allowbreak\ 9:5\allowbreak\ 10:3\allowbreak\ 11:5\allowbreak\ 15:3\allowbreak\ 16:5\allowbreak\ 17:3\allowbreak\ 18:5\allowbreak\ 22:3\allowbreak\ 23:5\allowbreak\ 24:3\allowbreak\ 61:5\allowbreak\ 63:4\allowbreak\ 100:5\allowbreak\ 101:6\allowbreak\ 102:5\allowbreak\ 106:6\allowbreak\ 107:5\allowbreak\ 108:6\allowbreak\ 109:5\allowbreak\ 113:6\allowbreak\ 114:5\allowbreak\ 115:6\allowbreak\ 116:5\par
\noindent \textbf{118}: 1:3\allowbreak\ 2:2\allowbreak\ 3:3\allowbreak\ 4:2\allowbreak\ 5:3\allowbreak\ 6:2\allowbreak\ 7:5\allowbreak\ 8:3\allowbreak\ 9:5\allowbreak\ 10:3\allowbreak\ 11:5\allowbreak\ 12:3\allowbreak\ 14:5\allowbreak\ 15:3\allowbreak\ 16:5\allowbreak\ 17:3\allowbreak\ 18:5\allowbreak\ 19:3\allowbreak\ 21:5\allowbreak\ 22:3\allowbreak\ 23:5\allowbreak\ 24:3\allowbreak\ 25:5\allowbreak\ 26:3\allowbreak\ 98:5\allowbreak\ 99:6\allowbreak\ 100:5\allowbreak\ 101:6\allowbreak\ 102:5\allowbreak\ 103:6\allowbreak\ 105:5\allowbreak\ 106:6\allowbreak\ 107:5\allowbreak\ 108:6\allowbreak\ 109:5\allowbreak\ 110:6\allowbreak\ 112:5\allowbreak\ 113:6\allowbreak\ 114:5\allowbreak\ 115:6\allowbreak\ 116:5\allowbreak\ 117:6\allowbreak\ 118:5\par
\noindent \textbf{120}: 2:4\allowbreak\ 3:3\allowbreak\ 4:2\allowbreak\ 9:5\allowbreak\ 10:3\allowbreak\ 16:5\allowbreak\ 17:3\allowbreak\ 23:5\allowbreak\ 24:3\allowbreak\ 61:5\allowbreak\ 63:4\allowbreak\ 100:5\allowbreak\ 101:6\allowbreak\ 107:5\allowbreak\ 108:6\allowbreak\ 114:5\allowbreak\ 115:6\allowbreak\ 120:5\par
\noindent \textbf{121}: 2:4\allowbreak\ 3:3\allowbreak\ 9:5\allowbreak\ 10:3\allowbreak\ 16:5\allowbreak\ 17:3\allowbreak\ 23:5\allowbreak\ 24:3\allowbreak\ 61:5\allowbreak\ 63:4\allowbreak\ 100:5\allowbreak\ 101:6\allowbreak\ 107:5\allowbreak\ 108:6\allowbreak\ 114:5\allowbreak\ 115:6\allowbreak\ 121:5\par
\noindent \textbf{123}: 2:4\allowbreak\ 31:5\allowbreak\ 32:4\allowbreak\ 61:5\allowbreak\ 63:4\allowbreak\ 92:5\allowbreak\ 93:6\allowbreak\ 122:5\allowbreak\ 123:6\par
\noindent \textbf{125}: 3:5\allowbreak\ 4:2\allowbreak\ 29:5\allowbreak\ 30:4\allowbreak\ 33:5\allowbreak\ 34:4\allowbreak\ 90:5\allowbreak\ 91:6\allowbreak\ 94:5\allowbreak\ 95:6\allowbreak\ 120:5\allowbreak\ 121:6\allowbreak\ 124:5\allowbreak\ 125:6\par
\noindent \textbf{126}: 2:4\allowbreak\ 5:2\allowbreak\ 28:5\allowbreak\ 35:4\allowbreak\ 61:5\allowbreak\ 63:4\allowbreak\ 89:5\allowbreak\ 96:6\allowbreak\ 119:5\allowbreak\ 126:6\par
\noindent \textbf{127}: 1:5\allowbreak\ 4:2\allowbreak\ 5:3\allowbreak\ 6:2\allowbreak\ 11:5\allowbreak\ 12:3\allowbreak\ 18:5\allowbreak\ 19:3\allowbreak\ 25:5\allowbreak\ 26:3\allowbreak\ 27:5\allowbreak\ 30:4\allowbreak\ 31:5\allowbreak\ 36:4\allowbreak\ 88:5\allowbreak\ 91:6\allowbreak\ 92:5\allowbreak\ 97:6\allowbreak\ 102:5\allowbreak\ 103:6\allowbreak\ 109:5\allowbreak\ 110:6\allowbreak\ 116:5\allowbreak\ 117:6\allowbreak\ 118:5\allowbreak\ 121:6\allowbreak\ 122:5\allowbreak\ 127:6\par
\noindent \textbf{129}: 3:5\allowbreak\ 4:2\allowbreak\ 8:3\allowbreak\ 11:5\allowbreak\ 15:3\allowbreak\ 18:5\allowbreak\ 22:3\allowbreak\ 25:5\allowbreak\ 30:4\allowbreak\ 33:5\allowbreak\ 38:4\allowbreak\ 86:5\allowbreak\ 91:6\allowbreak\ 94:5\allowbreak\ 99:6\allowbreak\ 102:5\allowbreak\ 106:6\allowbreak\ 109:5\allowbreak\ 113:6\allowbreak\ 116:5\allowbreak\ 121:6\allowbreak\ 124:5\allowbreak\ 129:6\par
\noindent \textbf{173}: 7:2\allowbreak\ 12:3\allowbreak\ 14:5\allowbreak\ 19:3\allowbreak\ 21:5\allowbreak\ 42:4\allowbreak\ 103:6\allowbreak\ 105:5\allowbreak\ 110:6\allowbreak\ 112:5\allowbreak\ 173:6\par
\noindent \textbf{176}: 7:2\allowbreak\ 12:3\allowbreak\ 14:5\allowbreak\ 17:4\allowbreak\ 19:3\allowbreak\ 24:5\allowbreak\ 42:4\allowbreak\ 103:6\allowbreak\ 105:5\allowbreak\ 110:6\allowbreak\ 115:5\allowbreak\ 166:6\allowbreak\ 169:7\allowbreak\ 176:6\par
\noindent \textbf{177}: 3:5\allowbreak\ 4:2\allowbreak\ 8:3\allowbreak\ 11:5\allowbreak\ 15:3\allowbreak\ 18:5\allowbreak\ 22:3\allowbreak\ 25:5\allowbreak\ 30:4\allowbreak\ 33:5\allowbreak\ 38:4\allowbreak\ 91:6\allowbreak\ 94:5\allowbreak\ 99:6\allowbreak\ 102:5\allowbreak\ 106:6\allowbreak\ 109:5\allowbreak\ 113:6\allowbreak\ 116:5\allowbreak\ 121:6\allowbreak\ 124:5\allowbreak\ 177:6\par
\noindent \textbf{179}: 3:5\allowbreak\ 11:3\allowbreak\ 15:4\allowbreak\ 23:3\allowbreak\ 30:4\allowbreak\ 33:5\allowbreak\ 40:4\allowbreak\ 48:6\allowbreak\ 52:4\allowbreak\ 91:6\allowbreak\ 94:5\allowbreak\ 121:6\allowbreak\ 124:5\allowbreak\ 163:6\allowbreak\ 167:7\allowbreak\ 175:6\allowbreak\ 179:7\par
\noindent \textbf{180}: 5:3\allowbreak\ 7:5\allowbreak\ 12:3\allowbreak\ 14:5\allowbreak\ 19:3\allowbreak\ 21:5\allowbreak\ 26:3\allowbreak\ 28:5\allowbreak\ 35:4\allowbreak\ 96:6\allowbreak\ 98:5\allowbreak\ 103:6\allowbreak\ 105:5\allowbreak\ 110:6\allowbreak\ 112:5\allowbreak\ 117:6\allowbreak\ 119:5\allowbreak\ 180:6\par
\noindent \textbf{181}: 3:5\allowbreak\ 4:2\allowbreak\ 29:5\allowbreak\ 30:4\allowbreak\ 33:5\allowbreak\ 34:4\allowbreak\ 91:6\allowbreak\ 94:5\allowbreak\ 95:6\allowbreak\ 120:5\allowbreak\ 121:6\allowbreak\ 124:5\allowbreak\ 181:6\par
\noindent \textbf{184}: 7:2\allowbreak\ 31:4\allowbreak\ 32:6\allowbreak\ 56:4\allowbreak\ 159:6\allowbreak\ 183:7\allowbreak\ 184:6\par
\noindent \textbf{185}: 21:5\allowbreak\ 26:3\allowbreak\ 30:4\allowbreak\ 51:6\allowbreak\ 58:4\allowbreak\ 91:6\allowbreak\ 112:5\allowbreak\ 178:6\allowbreak\ 185:7\par
\noindent \textbf{188}: 2:5\allowbreak\ 5:4\allowbreak\ 7:2\allowbreak\ 30:4\allowbreak\ 32:6\allowbreak\ 61:4\allowbreak\ 91:6\allowbreak\ 93:5\allowbreak\ 154:6\allowbreak\ 157:7\allowbreak\ 159:6\allowbreak\ 188:7\par
\noindent \textbf{194}: 7:2\allowbreak\ 12:3\allowbreak\ 14:5\allowbreak\ 19:3\allowbreak\ 21:5\allowbreak\ 103:6\allowbreak\ 105:5\allowbreak\ 110:6\allowbreak\ 112:5\allowbreak\ 194:6\par
\noindent \textbf{197}: 8:5\allowbreak\ 18:3\allowbreak\ 55:6\allowbreak\ 63:4\allowbreak\ 91:6\allowbreak\ 99:5\allowbreak\ 182:6\allowbreak\ 190:7\allowbreak\ 197:6\par
\noindent \textbf{198}: 2:4\allowbreak\ 5:2\allowbreak\ 28:5\allowbreak\ 35:4\allowbreak\ 53:6\allowbreak\ 61:5\allowbreak\ 63:4\allowbreak\ 71:5\allowbreak\ 96:6\allowbreak\ 119:5\allowbreak\ 144:6\allowbreak\ 162:7\allowbreak\ 180:6\allowbreak\ 198:7\par
\noindent \textbf{199}: 2:4\allowbreak\ 5:2\allowbreak\ 16:4\allowbreak\ 22:6\allowbreak\ 41:4\allowbreak\ 47:6\allowbreak\ 58:4\allowbreak\ 61:5\allowbreak\ 63:4\allowbreak\ 66:5\allowbreak\ 149:6\allowbreak\ 168:7\allowbreak\ 174:6\allowbreak\ 193:7\allowbreak\ 199:6\par
\noindent \textbf{200}: 26:3\allowbreak\ 37:4\allowbreak\ 51:6\allowbreak\ 73:5\allowbreak\ 142:6\allowbreak\ 164:7\allowbreak\ 178:6\allowbreak\ 200:7\par
\noindent \textbf{201}: 1:3\allowbreak\ 12:4\allowbreak\ 25:3\allowbreak\ 37:4\allowbreak\ 50:6\allowbreak\ 62:5\allowbreak\ 63:4\allowbreak\ 141:6\allowbreak\ 164:7\allowbreak\ 177:6\allowbreak\ 190:7\allowbreak\ 201:6\par
\noindent \textbf{202}: 13:4\allowbreak\ 25:6\allowbreak\ 38:4\allowbreak\ 50:6\allowbreak\ 63:4\allowbreak\ 152:6\allowbreak\ 165:7\allowbreak\ 177:6\allowbreak\ 190:7\allowbreak\ 202:6\par
\noindent \textbf{203}: 12:4\allowbreak\ 25:6\allowbreak\ 26:3\allowbreak\ 37:4\allowbreak\ 51:6\allowbreak\ 62:4\allowbreak\ 152:6\allowbreak\ 164:7\allowbreak\ 178:6\allowbreak\ 189:7\allowbreak\ 203:6\par
\noindent \textbf{204}: 11:4\allowbreak\ 15:6\allowbreak\ 26:3\allowbreak\ 36:4\allowbreak\ 51:6\allowbreak\ 61:5\allowbreak\ 142:6\allowbreak\ 163:7\allowbreak\ 178:6\allowbreak\ 188:7\allowbreak\ 204:6\par
\noindent \textbf{205}: 2:4\allowbreak\ 5:2\allowbreak\ 21:3\allowbreak\ 42:4\allowbreak\ 46:6\allowbreak\ 61:5\allowbreak\ 63:4\allowbreak\ 78:5\allowbreak\ 137:6\allowbreak\ 169:7\allowbreak\ 173:6\allowbreak\ 205:7\par
\noindent \textbf{206}: 7:2\allowbreak\ 18:3\allowbreak\ 19:4\allowbreak\ 20:6\allowbreak\ 44:4\allowbreak\ 45:6\allowbreak\ 55:5\allowbreak\ 56:4\allowbreak\ 79:5\allowbreak\ 136:6\allowbreak\ 146:7\allowbreak\ 147:6\allowbreak\ 171:7\allowbreak\ 172:6\allowbreak\ 206:7\par
\noindent \textbf{233}: 7:2\allowbreak\ 9:6\allowbreak\ 18:4\allowbreak\ 43:5\allowbreak\ 45:6\allowbreak\ 70:5\allowbreak\ 81:6\allowbreak\ 106:5\allowbreak\ 109:6\allowbreak\ 134:7\allowbreak\ 136:6\allowbreak\ 170:7\allowbreak\ 172:6\allowbreak\ 197:7\allowbreak\ 208:6\allowbreak\ 233:7\par
\noindent \textbf{234}: 16:6\allowbreak\ 36:5\allowbreak\ 52:6\allowbreak\ 72:5\allowbreak\ 88:6\allowbreak\ 108:5\allowbreak\ 127:7\allowbreak\ 143:6\allowbreak\ 163:7\allowbreak\ 179:6\allowbreak\ 199:7\allowbreak\ 215:6\allowbreak\ 234:7\par
\noindent \textbf{235}: 7:2\allowbreak\ 9:6\allowbreak\ 17:3\allowbreak\ 43:5\allowbreak\ 45:6\allowbreak\ 79:5\allowbreak\ 81:6\allowbreak\ 107:5\allowbreak\ 134:7\allowbreak\ 136:6\allowbreak\ 170:7\allowbreak\ 172:6\allowbreak\ 206:7\allowbreak\ 208:6\allowbreak\ 235:7\par
\noindent \textbf{236}: 7:2\allowbreak\ 18:6\allowbreak\ 43:5\allowbreak\ 45:6\allowbreak\ 54:4\allowbreak\ 79:5\allowbreak\ 81:6\allowbreak\ 106:5\allowbreak\ 134:7\allowbreak\ 145:6\allowbreak\ 170:7\allowbreak\ 172:6\allowbreak\ 206:7\allowbreak\ 208:6\allowbreak\ 236:7\par
\noindent \textbf{243}: 25:4\allowbreak\ 61:5\allowbreak\ 63:4\allowbreak\ 86:6\allowbreak\ 99:5\allowbreak\ 190:7\allowbreak\ 213:6\allowbreak\ 243:7\par
\noindent \textbf{245}: 27:4\allowbreak\ 61:5\allowbreak\ 88:6\allowbreak\ 97:5\allowbreak\ 188:7\allowbreak\ 215:6\allowbreak\ 245:7\par
\noindent \textbf{246}: 2:6\allowbreak\ 24:3\allowbreak\ 28:6\allowbreak\ 36:5\allowbreak\ 38:4\allowbreak\ 60:5\allowbreak\ 61:6\allowbreak\ 63:4\allowbreak\ 64:6\allowbreak\ 96:5\allowbreak\ 127:7\allowbreak\ 129:6\allowbreak\ 151:7\allowbreak\ 155:6\allowbreak\ 187:7\allowbreak\ 188:6\allowbreak\ 190:7\allowbreak\ 191:6\allowbreak\ 246:7\par
\noindent \textbf{247}: 2:6\allowbreak\ 3:2\allowbreak\ 4:3\allowbreak\ 5:2\allowbreak\ 22:6\allowbreak\ 23:3\allowbreak\ 29:6\allowbreak\ 36:5\allowbreak\ 38:4\allowbreak\ 41:5\allowbreak\ 58:4\allowbreak\ 59:5\allowbreak\ 61:6\allowbreak\ 63:4\allowbreak\ 65:6\allowbreak\ 95:5\allowbreak\ 127:7\allowbreak\ 129:6\allowbreak\ 132:7\allowbreak\ 149:6\allowbreak\ 150:7\allowbreak\ 156:6\allowbreak\ 186:7\allowbreak\ 188:6\allowbreak\ 190:7\allowbreak\ 192:6\allowbreak\ 247:7\par
\noindent \textbf{248}: 2:6\allowbreak\ 5:2\allowbreak\ 30:6\allowbreak\ 36:5\allowbreak\ 38:4\allowbreak\ 41:5\allowbreak\ 58:4\allowbreak\ 61:6\allowbreak\ 63:4\allowbreak\ 66:6\allowbreak\ 94:5\allowbreak\ 127:7\allowbreak\ 129:6\allowbreak\ 132:7\allowbreak\ 157:6\allowbreak\ 185:7\allowbreak\ 188:6\allowbreak\ 190:7\allowbreak\ 193:6\allowbreak\ 248:7\par
\noindent \textbf{249}: 1:3\allowbreak\ 6:2\allowbreak\ 31:6\allowbreak\ 42:5\allowbreak\ 57:4\allowbreak\ 62:6\allowbreak\ 63:4\allowbreak\ 67:6\allowbreak\ 93:5\allowbreak\ 133:7\allowbreak\ 158:6\allowbreak\ 184:7\allowbreak\ 189:6\allowbreak\ 190:7\allowbreak\ 194:6\allowbreak\ 249:7\par
\noindent \textbf{250}: 7:2\allowbreak\ 32:6\allowbreak\ 43:5\allowbreak\ 56:4\allowbreak\ 68:6\allowbreak\ 92:5\allowbreak\ 134:7\allowbreak\ 159:6\allowbreak\ 183:7\allowbreak\ 195:6\allowbreak\ 250:7\par
\noindent \textbf{251}: 7:2\allowbreak\ 19:3\allowbreak\ 33:6\allowbreak\ 55:5\allowbreak\ 69:6\allowbreak\ 91:5\allowbreak\ 146:7\allowbreak\ 160:6\allowbreak\ 182:7\allowbreak\ 196:6\allowbreak\ 251:7\par
\noindent \textbf{252}: 7:2\allowbreak\ 9:6\allowbreak\ 18:4\allowbreak\ 34:6\allowbreak\ 43:5\allowbreak\ 45:6\allowbreak\ 54:5\allowbreak\ 81:6\allowbreak\ 90:5\allowbreak\ 134:7\allowbreak\ 136:6\allowbreak\ 145:7\allowbreak\ 161:6\allowbreak\ 170:7\allowbreak\ 172:6\allowbreak\ 181:7\allowbreak\ 208:6\allowbreak\ 252:7\par
\noindent \textbf{253}: 2:4\allowbreak\ 24:3\allowbreak\ 25:6\allowbreak\ 27:4\allowbreak\ 28:6\allowbreak\ 35:4\allowbreak\ 60:5\allowbreak\ 64:6\allowbreak\ 86:5\allowbreak\ 88:6\allowbreak\ 89:5\allowbreak\ 151:7\allowbreak\ 152:6\allowbreak\ 154:7\allowbreak\ 155:6\allowbreak\ 187:7\allowbreak\ 191:6\allowbreak\ 213:7\allowbreak\ 215:6\allowbreak\ 253:7\par
\noindent \textbf{339}: 6:3\allowbreak\ 20:6\allowbreak\ 23:4\allowbreak\ 26:3\allowbreak\ 172:6\allowbreak\ 302:8\allowbreak\ 316:7\allowbreak\ 322:8\allowbreak\ 339:7\par
\noindent \textbf{341}: 1:3\allowbreak\ 21:7\allowbreak\ 62:6\allowbreak\ 63:4\allowbreak\ 189:6\allowbreak\ 190:7\allowbreak\ 317:8\allowbreak\ 341:7\par
\noindent \textbf{343}: 16:5\allowbreak\ 30:4\allowbreak\ 46:6\allowbreak\ 47:4\allowbreak\ 91:6\allowbreak\ 107:5\allowbreak\ 168:6\allowbreak\ 326:8\allowbreak\ 342:7\allowbreak\ 343:8\par
\noindent \textbf{344}: 40:7\allowbreak\ 42:6\allowbreak\ 82:5\allowbreak\ 133:6\allowbreak\ 167:7\allowbreak\ 169:8\allowbreak\ 173:7\allowbreak\ 175:6\allowbreak\ 209:7\allowbreak\ 302:8\allowbreak\ 342:7\allowbreak\ 344:8\par
\noindent \textbf{346}: 8:5\allowbreak\ 49:7\allowbreak\ 63:4\allowbreak\ 74:6\allowbreak\ 99:5\allowbreak\ 165:6\allowbreak\ 177:8\allowbreak\ 190:7\allowbreak\ 201:6\allowbreak\ 218:7\allowbreak\ 346:8\par
\noindent \textbf{347}: 9:3\allowbreak\ 37:4\allowbreak\ 46:6\allowbreak\ 51:4\allowbreak\ 161:6\allowbreak\ 164:7\allowbreak\ 330:8\allowbreak\ 342:7\allowbreak\ 347:8\par
\noindent \textbf{348}: 10:6\allowbreak\ 36:5\allowbreak\ 46:6\allowbreak\ 52:5\allowbreak\ 101:6\allowbreak\ 127:7\allowbreak\ 137:6\allowbreak\ 163:7\allowbreak\ 270:8\allowbreak\ 342:7\allowbreak\ 348:8\par
\noindent \textbf{349}: 46:7\allowbreak\ 53:4\allowbreak\ 71:6\allowbreak\ 78:5\allowbreak\ 162:6\allowbreak\ 169:8\allowbreak\ 198:6\allowbreak\ 205:7\allowbreak\ 215:6\allowbreak\ 349:8\par
\noindent \textbf{350}: 9:7\allowbreak\ 17:4\allowbreak\ 26:3\allowbreak\ 54:4\allowbreak\ 161:6\allowbreak\ 169:8\allowbreak\ 178:6\allowbreak\ 350:8\par
\noindent \textbf{353}: 6:3\allowbreak\ 7:2\allowbreak\ 20:6\allowbreak\ 25:3\allowbreak\ 43:5\allowbreak\ 57:4\allowbreak\ 116:6\allowbreak\ 134:7\allowbreak\ 158:6\allowbreak\ 285:8\allowbreak\ 316:7\allowbreak\ 353:8\par
\noindent \textbf{354}: 5:7\allowbreak\ 21:4\allowbreak\ 22:3\allowbreak\ 58:4\allowbreak\ 157:6\allowbreak\ 169:8\allowbreak\ 173:7\allowbreak\ 174:6\allowbreak\ 338:8\allowbreak\ 342:7\allowbreak\ 354:8\par
\noindent \textbf{358}: 6:3\allowbreak\ 7:2\allowbreak\ 43:7\allowbreak\ 62:4\allowbreak\ 85:5\allowbreak\ 153:6\allowbreak\ 176:8\allowbreak\ 189:6\allowbreak\ 212:7\allowbreak\ 358:8\par
\noindent \textbf{363}: 4:7\allowbreak\ 6:6\allowbreak\ 21:4\allowbreak\ 42:6\allowbreak\ 67:5\allowbreak\ 133:6\allowbreak\ 148:7\allowbreak\ 169:8\allowbreak\ 173:7\allowbreak\ 300:8\allowbreak\ 342:7\allowbreak\ 363:8\par

\end{multicols}

\clearpage
\section{All 191 short alternating-path words}
\noindent The integer at the left is the new size $n$;
the following one-digit symbols are the successive
cube roots $r_j$ in Lemma~\ref{lem:words}.
Together with the explicit seed in
Lemma~\ref{lem:seed173}, these words form an
independently and locally checkable existence
certificate for $173\le n\le364$.

\begin{multicols}{3}
\scriptsize\ttfamily
\setlength{\parindent}{0pt}
\setlength{\parskip}{0pt}
% Generated by finite exact cube-sum augmentation path checks.
% Each line lists n: successive integer cube roots on the alternating path.
\noindent 174: 63555536363555536363555536\par
\noindent 175: 6336\par
\noindent 176: 65555556\par
\noindent 177: 6364636\par
\noindent 178: 6364625555556\par
\noindent 179: 74635555555547\par
\noindent 180: 74474247447\par
\noindent 181: 6554746474556\par
\noindent 182: 655556\par
\noindent 183: 663647\par
\noindent 184: 66\par
\noindent 185: 66264746366\par
\noindent 186: 66266\par
\noindent 187: 7447\par
\noindent 188: 7455665547\par
\noindent 189: 666363647447447\par
\noindent 190: 647\par
\noindent 191: 6665547556\par
\noindent 192: 75257\par
\noindent 193: 7575555757\par
\noindent 194: 66655666\par
\noindent 195: 6264757\par
\noindent 196: 663666626\par
\noindent 197: 75746264757\par
\noindent 198: 65575264757\par
\noindent 199: 666\par
\noindent 200: 66633666\par
\noindent 201: 64747574757\par
\noindent 202: 6474752556\par
\noindent 203: 663466666\par
\noindent 204: 757424757\par
\noindent 205: 65755756\par
\noindent 206: 6666446666\par
\noindent 207: 666665546666\par
\noindent 208: 64757\par
\noindent 209: 6443756\par
\noindent 210: 757464757\par
\noindent 211: 6666357\par
\noindent 212: 66664666\par
\noindent 213: 757475757\par
\noindent 214: 653756\par
\noindent 215: 6666556\par
\noindent 216: 7555666666665557\par
\noindent 217: 75757375375757375757\par
\noindent 218: 7555757573573757575557\par
\noindent 219: 757455757464757\par
\noindent 220: 75757554665575757\par
\noindent 221: 757573753666375757\par
\noindent 222: 7375757\par
\noindent 223: 757455755557554757\par
\noindent 224: 7363753637\par
\noindent 225: 7557357447557\par
\noindent 226: 75447557\par
\noindent 227: 757537525735757\par
\noindent 228: 7524665566666657\par
\noindent 229: 746664667\par
\noindent 230: 752757573637557\par
\noindent 231: 75666637575725257\par
\noindent 232: 757536357556666665757\par
\noindent 233: 75666666557537\par
\noindent 234: 7375757256666667\par
\noindent 235: 77\par
\noindent 236: 74664664664664664664757275747\par
\noindent 237: 7375747\par
\noindent 238: 74664747\par
\noindent 239: 73444667\par
\noindent 240: 7466466444667\par
\noindent 241: 7337\par
\noindent 242: 734664667\par
\noindent 243: 73527537534667\par
\noindent 244: 73667\par
\noindent 245: 746666667\par
\noindent 246: 766375375375266667\par
\noindent 247: 76666666666666666667\par
\noindent 248: 76666266667\par
\noindent 249: 766662757475757266667\par
\noindent 250: 747\par
\noindent 251: 766665736666667\par
\noindent 252: 747573375747\par
\noindent 253: 7474466375747\par
\noindent 254: 74744664664747\par
\noindent 255: 74757377\par
\noindent 256: 8\par
\noindent 257: 76666667\par
\noindent 258: 747525747\par
\noindent 259: 85666658566277\par
\noindent 260: 85662666666667\par
\noindent 261: 7475375758\par
\noindent 262: 8558\par
\noindent 263: 7736666444758\par
\noindent 264: 7737574744758\par
\noindent 265: 8558534758\par
\noindent 266: 766664758\par
\noindent 267: 85585757377\par
\noindent 268: 7477\par
\noindent 269: 766666677477\par
\noindent 270: 8575247477\par
\noindent 271: 8575266667\par
\noindent 272: 7475375375758575758\par
\noindent 273: 855857573757447\par
\noindent 274: 85266258\par
\noindent 275: 7437447\par
\noindent 276: 74586666666658\par
\noindent 277: 868\par
\noindent 278: 7447\par
\noindent 279: 868525377\par
\noindent 280: 773368\par
\noindent 281: 76666477\par
\noindent 282: 72755857477\par
\noindent 283: 727\par
\noindent 284: 777558\par
\noindent 285: 727575868\par
\noindent 286: 85475758\par
\noindent 287: 76658\par
\noindent 288: 7577666658\par
\noindent 289: 8666857537\par
\noindent 290: 8558\par
\noindent 291: 856666667\par
\noindent 292: 86468\par
\noindent 293: 8558667\par
\noindent 294: 737575727737\par
\noindent 295: 86477737\par
\noindent 296: 7775486477\par
\noindent 297: 83747775868\par
\noindent 298: 7577738\par
\noindent 299: 7775468\par
\noindent 300: 8646368585486868\par
\noindent 301: 7368666468\par
\noindent 302: 76644867\par
\noindent 303: 734775557\par
\noindent 304: 736868\par
\noindent 305: 8677777\par
\noindent 306: 7586666856768\par
\noindent 307: 7684777558\par
\noindent 308: 768458667\par
\noindent 309: 77776667\par
\noindent 310: 76848\par
\noindent 311: 77\par
\noindent 312: 836366625868\par
\noindent 313: 775868577\par
\noindent 314: 77457777\par
\noindent 315: 7764777\par
\noindent 316: 867738\par
\noindent 317: 76664468\par
\noindent 318: 7668667\par
\noindent 319: 864866367\par
\noindent 320: 764777\par
\noindent 321: 8685485677\par
\noindent 322: 7684847777\par
\noindent 323: 75737668\par
\noindent 324: 7667677367\par
\noindent 325: 7586268\par
\noindent 326: 7586847\par
\noindent 327: 8685758268\par
\noindent 328: 8626846668\par
\noindent 329: 7577\par
\noindent 330: 86856768667\par
\noindent 331: 7585268368\par
\noindent 332: 867\par
\noindent 333: 777777\par
\noindent 334: 7686855758\par
\noindent 335: 76358668668\par
\noindent 336: 784686847586868\par
\noindent 337: 8638486777\par
\noindent 338: 7468487\par
\noindent 339: 77\par
\noindent 340: 77583867\par
\noindent 341: 7777577\par
\noindent 342: 88288\par
\noindent 343: 86388\par
\noindent 344: 86867777778\par
\noindent 345: 858557288\par
\noindent 346: 85777488\par
\noindent 347: 8586677768\par
\noindent 348: 8855676668\par
\noindent 349: 8848\par
\noindent 350: 8485738686388\par
\noindent 351: 86864485882778\par
\noindent 352: 8686488\par
\noindent 353: 8577664648\par
\noindent 354: 85868684848\par
\noindent 355: 877388\par
\noindent 356: 84577488\par
\noindent 357: 86777755778\par
\noindent 358: 858557785588\par
\noindent 359: 8468\par
\noindent 360: 8586482775858\par
\noindent 361: 8586855588\par
\noindent 362: 883477778\par
\noindent 363: 8528686488\par
\noindent 364: 855557388\par

\end{multicols}

\begin{thebibliography}{9}
\bibitem{oeisA006063}
N.~J.~A. Sloane et al., \emph{A006063: A card-arranging problem},
\url{https://oeis.org/A006063}; existence conjecture
attributed there to Robert Israel (2018).
\bibitem{oeisA096680}
Ray Chandler, \emph{A096680: Values with more than one cube-sum
permutation}, \url{https://oeis.org/A096680}.
\bibitem{oeisA096904}
Ray Chandler, \emph{A096904: Number of cube-sum permutations},
\url{https://oeis.org/A096904}.
\bibitem{hall}
P. Hall, On representatives of subsets,
\emph{J. London Math. Soc.} 10 (1935), 26--30.
\end{thebibliography}

\clearpage
\section{Reduced set of 56 algebraic cubic-order shapes}
\label{app:reduced-shapes}
The critical length $n=173$ is proved directly
by forced-reflection peeling rather than from
the determinant parameter table. For each of
the remaining $56$ shapes, the block lengths
are positive integer determinants by
Theorem~\ref{thm:unimodular-block-factory}.
The excluded original $57$-shape table
remains an optional independent audit.
\begin{multicols}{2}
\scriptsize
\setlength{\parindent}{0pt}
\setlength{\parskip}{2pt}
% 56 residual cubic-order determinant shapes, 784 blocks.
% The excluded seed n=173 has a direct forced-leaf uniqueness proof in FORCED_REFLECTION_PEELING.md.
% Outputs are not given; all block lengths are computed symbolically from Cramer determinants.
\noindent \textbf{82}: roots \texttt{4,\allowbreak 2,\allowbreak 3,\allowbreak 4,\allowbreak 5,\allowbreak 4,\allowbreak 5}; order \texttt{6,\allowbreak 2,\allowbreak 3,\allowbreak 4,\allowbreak 7,\allowbreak 1,\allowbreak 5}\par
\noindent \textbf{116}: roots \texttt{4,\allowbreak 3,\allowbreak 2,\allowbreak 3,\allowbreak 5,\allowbreak 3,\allowbreak 5,\allowbreak 3,\allowbreak 5,\allowbreak 3,\allowbreak 5,\allowbreak 3,\allowbreak 5,\allowbreak 3,\allowbreak 5,\allowbreak 4,\allowbreak 5,\allowbreak 6,\allowbreak 5,\allowbreak 6,\allowbreak 5,\allowbreak 6,\allowbreak 5,\allowbreak 6,\allowbreak 5,\allowbreak 6,\allowbreak 5}; order \texttt{16,\allowbreak 14,\allowbreak 3,\allowbreak 12,\allowbreak 27,\allowbreak 10,\allowbreak 25,\allowbreak 8,\allowbreak 23,\allowbreak 6,\allowbreak 21,\allowbreak 4,\allowbreak 19,\allowbreak 2,\allowbreak 17,\allowbreak 1,\allowbreak 15,\allowbreak 26,\allowbreak 13,\allowbreak 24,\allowbreak 11,\allowbreak 22,\allowbreak 9,\allowbreak 20,\allowbreak 7,\allowbreak 18,\allowbreak 5}\par
\noindent \textbf{118}: roots \texttt{3,\allowbreak 2,\allowbreak 3,\allowbreak 2,\allowbreak 3,\allowbreak 2,\allowbreak 5,\allowbreak 3,\allowbreak 5,\allowbreak 3,\allowbreak 5,\allowbreak 3,\allowbreak 5,\allowbreak 3,\allowbreak 5,\allowbreak 3,\allowbreak 5,\allowbreak 3,\allowbreak 5,\allowbreak 3,\allowbreak 5,\allowbreak 3,\allowbreak 5,\allowbreak 3,\allowbreak 5,\allowbreak 6,\allowbreak 5,\allowbreak 6,\allowbreak 5,\allowbreak 6,\allowbreak 5,\allowbreak 6,\allowbreak 5,\allowbreak 6,\allowbreak 5,\allowbreak 6,\allowbreak 5,\allowbreak 6,\allowbreak 5,\allowbreak 6,\allowbreak 5,\allowbreak 6,\allowbreak 5}; order \texttt{24,\allowbreak 6,\allowbreak 22,\allowbreak 4,\allowbreak 20,\allowbreak 2,\allowbreak 43,\allowbreak 18,\allowbreak 41,\allowbreak 16,\allowbreak 39,\allowbreak 14,\allowbreak 37,\allowbreak 12,\allowbreak 35,\allowbreak 10,\allowbreak 33,\allowbreak 8,\allowbreak 31,\allowbreak 5,\allowbreak 29,\allowbreak 3,\allowbreak 27,\allowbreak 1,\allowbreak 25,\allowbreak 42,\allowbreak 23,\allowbreak 40,\allowbreak 21,\allowbreak 38,\allowbreak 19,\allowbreak 36,\allowbreak 17,\allowbreak 34,\allowbreak 15,\allowbreak 32,\allowbreak 13,\allowbreak 30,\allowbreak 11,\allowbreak 28,\allowbreak 9,\allowbreak 26,\allowbreak 7}\par
\noindent \textbf{120}: roots \texttt{4,\allowbreak 3,\allowbreak 2,\allowbreak 5,\allowbreak 3,\allowbreak 5,\allowbreak 3,\allowbreak 5,\allowbreak 3,\allowbreak 5,\allowbreak 4,\allowbreak 5,\allowbreak 6,\allowbreak 5,\allowbreak 6,\allowbreak 5,\allowbreak 6,\allowbreak 5}; order \texttt{11,\allowbreak 9,\allowbreak 3,\allowbreak 18,\allowbreak 7,\allowbreak 16,\allowbreak 5,\allowbreak 14,\allowbreak 2,\allowbreak 12,\allowbreak 1,\allowbreak 10,\allowbreak 17,\allowbreak 8,\allowbreak 15,\allowbreak 6,\allowbreak 13,\allowbreak 4}\par
\noindent \textbf{121}: roots \texttt{4,\allowbreak 3,\allowbreak 5,\allowbreak 3,\allowbreak 5,\allowbreak 3,\allowbreak 5,\allowbreak 3,\allowbreak 5,\allowbreak 4,\allowbreak 5,\allowbreak 6,\allowbreak 5,\allowbreak 6,\allowbreak 5,\allowbreak 6,\allowbreak 5}; order \texttt{10,\allowbreak 8,\allowbreak 17,\allowbreak 6,\allowbreak 15,\allowbreak 4,\allowbreak 13,\allowbreak 2,\allowbreak 11,\allowbreak 1,\allowbreak 9,\allowbreak 16,\allowbreak 7,\allowbreak 14,\allowbreak 5,\allowbreak 12,\allowbreak 3}\par
\noindent \textbf{123}: roots \texttt{4,\allowbreak 5,\allowbreak 4,\allowbreak 5,\allowbreak 4,\allowbreak 5,\allowbreak 6,\allowbreak 5,\allowbreak 6}; order \texttt{5,\allowbreak 8,\allowbreak 3,\allowbreak 6,\allowbreak 1,\allowbreak 4,\allowbreak 9,\allowbreak 2,\allowbreak 7}\par
\noindent \textbf{125}: roots \texttt{5,\allowbreak 2,\allowbreak 5,\allowbreak 4,\allowbreak 5,\allowbreak 4,\allowbreak 5,\allowbreak 6,\allowbreak 5,\allowbreak 6,\allowbreak 5,\allowbreak 6,\allowbreak 5,\allowbreak 6}; order \texttt{13,\allowbreak 2,\allowbreak 11,\allowbreak 6,\allowbreak 9,\allowbreak 4,\allowbreak 7,\allowbreak 14,\allowbreak 5,\allowbreak 12,\allowbreak 3,\allowbreak 10,\allowbreak 1,\allowbreak 8}\par
\noindent \textbf{126}: roots \texttt{4,\allowbreak 2,\allowbreak 5,\allowbreak 4,\allowbreak 5,\allowbreak 4,\allowbreak 5,\allowbreak 6,\allowbreak 5,\allowbreak 6}; order \texttt{6,\allowbreak 2,\allowbreak 9,\allowbreak 4,\allowbreak 7,\allowbreak 1,\allowbreak 5,\allowbreak 10,\allowbreak 3,\allowbreak 8}\par
\noindent \textbf{127}: roots \texttt{5,\allowbreak 2,\allowbreak 3,\allowbreak 2,\allowbreak 5,\allowbreak 3,\allowbreak 5,\allowbreak 3,\allowbreak 5,\allowbreak 3,\allowbreak 5,\allowbreak 4,\allowbreak 5,\allowbreak 4,\allowbreak 5,\allowbreak 6,\allowbreak 5,\allowbreak 6,\allowbreak 5,\allowbreak 6,\allowbreak 5,\allowbreak 6,\allowbreak 5,\allowbreak 6,\allowbreak 5,\allowbreak 6,\allowbreak 5,\allowbreak 6}; order \texttt{10,\allowbreak 4,\allowbreak 27,\allowbreak 2,\allowbreak 25,\allowbreak 8,\allowbreak 23,\allowbreak 6,\allowbreak 21,\allowbreak 3,\allowbreak 19,\allowbreak 14,\allowbreak 17,\allowbreak 12,\allowbreak 15,\allowbreak 28,\allowbreak 13,\allowbreak 26,\allowbreak 11,\allowbreak 24,\allowbreak 9,\allowbreak 22,\allowbreak 7,\allowbreak 20,\allowbreak 5,\allowbreak 18,\allowbreak 1,\allowbreak 16}\par
\noindent \textbf{129}: roots \texttt{5,\allowbreak 2,\allowbreak 3,\allowbreak 5,\allowbreak 3,\allowbreak 5,\allowbreak 3,\allowbreak 5,\allowbreak 4,\allowbreak 5,\allowbreak 4,\allowbreak 5,\allowbreak 6,\allowbreak 5,\allowbreak 6,\allowbreak 5,\allowbreak 6,\allowbreak 5,\allowbreak 6,\allowbreak 5,\allowbreak 6,\allowbreak 5,\allowbreak 6}; order \texttt{22,\allowbreak 2,\allowbreak 7,\allowbreak 20,\allowbreak 5,\allowbreak 18,\allowbreak 3,\allowbreak 16,\allowbreak 11,\allowbreak 14,\allowbreak 9,\allowbreak 12,\allowbreak 23,\allowbreak 10,\allowbreak 21,\allowbreak 8,\allowbreak 19,\allowbreak 6,\allowbreak 17,\allowbreak 4,\allowbreak 15,\allowbreak 1,\allowbreak 13}\par
\noindent \textbf{176}: roots \texttt{2,\allowbreak 3,\allowbreak 5,\allowbreak 4,\allowbreak 3,\allowbreak 5,\allowbreak 4,\allowbreak 6,\allowbreak 5,\allowbreak 6,\allowbreak 5,\allowbreak 6,\allowbreak 7,\allowbreak 6}; order \texttt{1,\allowbreak 5,\allowbreak 11,\allowbreak 2,\allowbreak 9,\allowbreak 7,\allowbreak 14,\allowbreak 4,\allowbreak 12,\allowbreak 6,\allowbreak 10,\allowbreak 3,\allowbreak 8,\allowbreak 13}\par
\noindent \textbf{177}: roots \texttt{5,\allowbreak 2,\allowbreak 3,\allowbreak 5,\allowbreak 3,\allowbreak 5,\allowbreak 3,\allowbreak 5,\allowbreak 4,\allowbreak 5,\allowbreak 4,\allowbreak 6,\allowbreak 5,\allowbreak 6,\allowbreak 5,\allowbreak 6,\allowbreak 5,\allowbreak 6,\allowbreak 5,\allowbreak 6,\allowbreak 5,\allowbreak 6}; order \texttt{21,\allowbreak 2,\allowbreak 7,\allowbreak 19,\allowbreak 5,\allowbreak 17,\allowbreak 3,\allowbreak 15,\allowbreak 11,\allowbreak 13,\allowbreak 9,\allowbreak 22,\allowbreak 10,\allowbreak 20,\allowbreak 8,\allowbreak 18,\allowbreak 6,\allowbreak 16,\allowbreak 4,\allowbreak 14,\allowbreak 1,\allowbreak 12}\par
\noindent \textbf{179}: roots \texttt{5,\allowbreak 3,\allowbreak 4,\allowbreak 3,\allowbreak 4,\allowbreak 5,\allowbreak 4,\allowbreak 6,\allowbreak 4,\allowbreak 6,\allowbreak 5,\allowbreak 6,\allowbreak 5,\allowbreak 6,\allowbreak 7,\allowbreak 6,\allowbreak 7}; order \texttt{13,\allowbreak 4,\allowbreak 9,\allowbreak 2,\allowbreak 7,\allowbreak 11,\allowbreak 5,\allowbreak 16,\allowbreak 3,\allowbreak 14,\allowbreak 6,\allowbreak 12,\allowbreak 1,\allowbreak 10,\allowbreak 17,\allowbreak 8,\allowbreak 15}\par
\noindent \textbf{180}: roots \texttt{3,\allowbreak 5,\allowbreak 3,\allowbreak 5,\allowbreak 3,\allowbreak 5,\allowbreak 3,\allowbreak 5,\allowbreak 4,\allowbreak 6,\allowbreak 5,\allowbreak 6,\allowbreak 5,\allowbreak 6,\allowbreak 5,\allowbreak 6,\allowbreak 5,\allowbreak 6}; order \texttt{7,\allowbreak 17,\allowbreak 5,\allowbreak 15,\allowbreak 3,\allowbreak 13,\allowbreak 1,\allowbreak 11,\allowbreak 9,\allowbreak 18,\allowbreak 8,\allowbreak 16,\allowbreak 6,\allowbreak 14,\allowbreak 4,\allowbreak 12,\allowbreak 2,\allowbreak 10}\par
\noindent \textbf{181}: roots \texttt{5,\allowbreak 2,\allowbreak 5,\allowbreak 4,\allowbreak 5,\allowbreak 4,\allowbreak 6,\allowbreak 5,\allowbreak 6,\allowbreak 5,\allowbreak 6,\allowbreak 5,\allowbreak 6}; order \texttt{12,\allowbreak 2,\allowbreak 10,\allowbreak 6,\allowbreak 8,\allowbreak 4,\allowbreak 13,\allowbreak 5,\allowbreak 11,\allowbreak 3,\allowbreak 9,\allowbreak 1,\allowbreak 7}\par
\noindent \textbf{184}: roots \texttt{2,\allowbreak 4,\allowbreak 6,\allowbreak 4,\allowbreak 6,\allowbreak 7,\allowbreak 6}; order \texttt{1,\allowbreak 4,\allowbreak 7,\allowbreak 2,\allowbreak 5,\allowbreak 6,\allowbreak 3}\par
\noindent \textbf{185}: roots \texttt{5,\allowbreak 3,\allowbreak 4,\allowbreak 6,\allowbreak 4,\allowbreak 6,\allowbreak 5,\allowbreak 6,\allowbreak 7}; order \texttt{2,\allowbreak 5,\allowbreak 7,\allowbreak 3,\allowbreak 8,\allowbreak 1,\allowbreak 6,\allowbreak 9,\allowbreak 4}\par
\noindent \textbf{188}: roots \texttt{5,\allowbreak 4,\allowbreak 2,\allowbreak 4,\allowbreak 6,\allowbreak 4,\allowbreak 6,\allowbreak 5,\allowbreak 6,\allowbreak 7,\allowbreak 6,\allowbreak 7}; order \texttt{3,\allowbreak 6,\allowbreak 8,\allowbreak 4,\allowbreak 11,\allowbreak 2,\allowbreak 9,\allowbreak 1,\allowbreak 7,\allowbreak 12,\allowbreak 5,\allowbreak 10}\par
\noindent \textbf{194}: roots \texttt{2,\allowbreak 3,\allowbreak 5,\allowbreak 3,\allowbreak 5,\allowbreak 6,\allowbreak 5,\allowbreak 6,\allowbreak 5,\allowbreak 6}; order \texttt{1,\allowbreak 4,\allowbreak 9,\allowbreak 2,\allowbreak 7,\allowbreak 10,\allowbreak 5,\allowbreak 8,\allowbreak 3,\allowbreak 6}\par
\noindent \textbf{197}: roots \texttt{5,\allowbreak 3,\allowbreak 6,\allowbreak 4,\allowbreak 6,\allowbreak 5,\allowbreak 6,\allowbreak 7,\allowbreak 6}; order \texttt{4,\allowbreak 2,\allowbreak 9,\allowbreak 6,\allowbreak 7,\allowbreak 1,\allowbreak 5,\allowbreak 8,\allowbreak 3}\par
\noindent \textbf{198}: roots \texttt{4,\allowbreak 2,\allowbreak 5,\allowbreak 4,\allowbreak 6,\allowbreak 5,\allowbreak 4,\allowbreak 5,\allowbreak 6,\allowbreak 5,\allowbreak 6,\allowbreak 7,\allowbreak 6,\allowbreak 7}; order \texttt{7,\allowbreak 2,\allowbreak 10,\allowbreak 4,\allowbreak 13,\allowbreak 8,\allowbreak 1,\allowbreak 6,\allowbreak 11,\allowbreak 3,\allowbreak 9,\allowbreak 14,\allowbreak 5,\allowbreak 12}\par
\noindent \textbf{199}: roots \texttt{4,\allowbreak 2,\allowbreak 4,\allowbreak 6,\allowbreak 4,\allowbreak 6,\allowbreak 4,\allowbreak 5,\allowbreak 4,\allowbreak 5,\allowbreak 6,\allowbreak 7,\allowbreak 6,\allowbreak 7,\allowbreak 6}; order \texttt{9,\allowbreak 2,\allowbreak 7,\allowbreak 15,\allowbreak 5,\allowbreak 13,\allowbreak 3,\allowbreak 10,\allowbreak 1,\allowbreak 8,\allowbreak 11,\allowbreak 14,\allowbreak 6,\allowbreak 12,\allowbreak 4}\par
\noindent \textbf{200}: roots \texttt{3,\allowbreak 4,\allowbreak 6,\allowbreak 5,\allowbreak 6,\allowbreak 7,\allowbreak 6,\allowbreak 7}; order \texttt{1,\allowbreak 2,\allowbreak 7,\allowbreak 4,\allowbreak 5,\allowbreak 8,\allowbreak 3,\allowbreak 6}\par
\noindent \textbf{201}: roots \texttt{3,\allowbreak 4,\allowbreak 3,\allowbreak 4,\allowbreak 6,\allowbreak 5,\allowbreak 4,\allowbreak 6,\allowbreak 7,\allowbreak 6,\allowbreak 7,\allowbreak 6}; order \texttt{7,\allowbreak 3,\allowbreak 12,\allowbreak 1,\allowbreak 4,\allowbreak 10,\allowbreak 2,\allowbreak 6,\allowbreak 8,\allowbreak 11,\allowbreak 5,\allowbreak 9}\par
\noindent \textbf{202}: roots \texttt{4,\allowbreak 6,\allowbreak 4,\allowbreak 6,\allowbreak 4,\allowbreak 6,\allowbreak 7,\allowbreak 6,\allowbreak 7,\allowbreak 6}; order \texttt{5,\allowbreak 10,\allowbreak 3,\allowbreak 8,\allowbreak 1,\allowbreak 6,\allowbreak 9,\allowbreak 4,\allowbreak 7,\allowbreak 2}\par
\noindent \textbf{203}: roots \texttt{4,\allowbreak 6,\allowbreak 3,\allowbreak 4,\allowbreak 6,\allowbreak 4,\allowbreak 6,\allowbreak 7,\allowbreak 6,\allowbreak 7,\allowbreak 6}; order \texttt{3,\allowbreak 6,\allowbreak 11,\allowbreak 4,\allowbreak 9,\allowbreak 1,\allowbreak 7,\allowbreak 10,\allowbreak 5,\allowbreak 8,\allowbreak 2}\par
\noindent \textbf{204}: roots \texttt{4,\allowbreak 6,\allowbreak 3,\allowbreak 4,\allowbreak 6,\allowbreak 5,\allowbreak 6,\allowbreak 7,\allowbreak 6,\allowbreak 7,\allowbreak 6}; order \texttt{3,\allowbreak 11,\allowbreak 4,\allowbreak 9,\allowbreak 1,\allowbreak 6,\allowbreak 7,\allowbreak 10,\allowbreak 5,\allowbreak 8,\allowbreak 2}\par
\noindent \textbf{205}: roots \texttt{4,\allowbreak 2,\allowbreak 3,\allowbreak 4,\allowbreak 6,\allowbreak 5,\allowbreak 4,\allowbreak 5,\allowbreak 6,\allowbreak 7,\allowbreak 6,\allowbreak 7}; order \texttt{7,\allowbreak 2,\allowbreak 3,\allowbreak 4,\allowbreak 11,\allowbreak 8,\allowbreak 1,\allowbreak 6,\allowbreak 9,\allowbreak 12,\allowbreak 5,\allowbreak 10}\par
\noindent \textbf{206}: roots \texttt{2,\allowbreak 3,\allowbreak 4,\allowbreak 6,\allowbreak 4,\allowbreak 6,\allowbreak 5,\allowbreak 4,\allowbreak 5,\allowbreak 6,\allowbreak 7,\allowbreak 6,\allowbreak 7,\allowbreak 6,\allowbreak 7}; order \texttt{1,\allowbreak 8,\allowbreak 2,\allowbreak 5,\allowbreak 14,\allowbreak 3,\allowbreak 9,\allowbreak 12,\allowbreak 7,\allowbreak 10,\allowbreak 15,\allowbreak 6,\allowbreak 13,\allowbreak 4,\allowbreak 11}\par
\noindent \textbf{233}: roots \texttt{2,\allowbreak 6,\allowbreak 4,\allowbreak 5,\allowbreak 6,\allowbreak 5,\allowbreak 6,\allowbreak 5,\allowbreak 6,\allowbreak 7,\allowbreak 6,\allowbreak 7,\allowbreak 6,\allowbreak 7,\allowbreak 6,\allowbreak 7}; order \texttt{1,\allowbreak 15,\allowbreak 8,\allowbreak 13,\allowbreak 3,\allowbreak 6,\allowbreak 11,\allowbreak 4,\allowbreak 9,\allowbreak 16,\allowbreak 7,\allowbreak 14,\allowbreak 5,\allowbreak 12,\allowbreak 2,\allowbreak 10}\par
\noindent \textbf{234}: roots \texttt{6,\allowbreak 5,\allowbreak 6,\allowbreak 5,\allowbreak 6,\allowbreak 5,\allowbreak 7,\allowbreak 6,\allowbreak 7,\allowbreak 6,\allowbreak 7,\allowbreak 6,\allowbreak 7}; order \texttt{12,\allowbreak 6,\allowbreak 10,\allowbreak 4,\allowbreak 8,\allowbreak 2,\allowbreak 13,\allowbreak 5,\allowbreak 11,\allowbreak 3,\allowbreak 9,\allowbreak 1,\allowbreak 7}\par
\noindent \textbf{235}: roots \texttt{2,\allowbreak 6,\allowbreak 3,\allowbreak 5,\allowbreak 6,\allowbreak 5,\allowbreak 6,\allowbreak 5,\allowbreak 7,\allowbreak 6,\allowbreak 7,\allowbreak 6,\allowbreak 7,\allowbreak 6,\allowbreak 7}; order \texttt{1,\allowbreak 14,\allowbreak 3,\allowbreak 8,\allowbreak 12,\allowbreak 6,\allowbreak 10,\allowbreak 4,\allowbreak 15,\allowbreak 7,\allowbreak 13,\allowbreak 5,\allowbreak 11,\allowbreak 2,\allowbreak 9}\par
\noindent \textbf{236}: roots \texttt{2,\allowbreak 6,\allowbreak 5,\allowbreak 6,\allowbreak 4,\allowbreak 5,\allowbreak 6,\allowbreak 5,\allowbreak 7,\allowbreak 6,\allowbreak 7,\allowbreak 6,\allowbreak 7,\allowbreak 6,\allowbreak 7}; order \texttt{1,\allowbreak 14,\allowbreak 5,\allowbreak 8,\allowbreak 12,\allowbreak 6,\allowbreak 10,\allowbreak 3,\allowbreak 15,\allowbreak 7,\allowbreak 13,\allowbreak 4,\allowbreak 11,\allowbreak 2,\allowbreak 9}\par
\noindent \textbf{243}: roots \texttt{4,\allowbreak 5,\allowbreak 4,\allowbreak 6,\allowbreak 5,\allowbreak 7,\allowbreak 6,\allowbreak 7}; order \texttt{3,\allowbreak 7,\allowbreak 5,\allowbreak 1,\allowbreak 2,\allowbreak 8,\allowbreak 4,\allowbreak 6}\par
\noindent \textbf{245}: roots \texttt{4,\allowbreak 5,\allowbreak 6,\allowbreak 5,\allowbreak 7,\allowbreak 6,\allowbreak 7}; order \texttt{6,\allowbreak 4,\allowbreak 1,\allowbreak 2,\allowbreak 7,\allowbreak 3,\allowbreak 5}\par
\noindent \textbf{246}: roots \texttt{6,\allowbreak 3,\allowbreak 6,\allowbreak 5,\allowbreak 4,\allowbreak 5,\allowbreak 6,\allowbreak 4,\allowbreak 6,\allowbreak 5,\allowbreak 7,\allowbreak 6,\allowbreak 7,\allowbreak 6,\allowbreak 7,\allowbreak 6,\allowbreak 7,\allowbreak 6,\allowbreak 7}; order \texttt{8,\allowbreak 2,\allowbreak 18,\allowbreak 5,\allowbreak 16,\allowbreak 10,\allowbreak 14,\allowbreak 6,\allowbreak 12,\allowbreak 4,\allowbreak 19,\allowbreak 9,\allowbreak 17,\allowbreak 7,\allowbreak 15,\allowbreak 3,\allowbreak 13,\allowbreak 1,\allowbreak 11}\par
\noindent \textbf{247}: roots \texttt{6,\allowbreak 2,\allowbreak 3,\allowbreak 2,\allowbreak 6,\allowbreak 3,\allowbreak 6,\allowbreak 5,\allowbreak 4,\allowbreak 5,\allowbreak 4,\allowbreak 5,\allowbreak 6,\allowbreak 4,\allowbreak 6,\allowbreak 5,\allowbreak 7,\allowbreak 6,\allowbreak 7,\allowbreak 6,\allowbreak 7,\allowbreak 6,\allowbreak 7,\allowbreak 6,\allowbreak 7,\allowbreak 6,\allowbreak 7}; order \texttt{14,\allowbreak 4,\allowbreak 6,\allowbreak 2,\allowbreak 11,\allowbreak 3,\allowbreak 26,\allowbreak 9,\allowbreak 24,\allowbreak 16,\allowbreak 22,\allowbreak 12,\allowbreak 20,\allowbreak 10,\allowbreak 18,\allowbreak 8,\allowbreak 27,\allowbreak 15,\allowbreak 25,\allowbreak 13,\allowbreak 23,\allowbreak 7,\allowbreak 21,\allowbreak 5,\allowbreak 19,\allowbreak 1,\allowbreak 17}\par
\noindent \textbf{248}: roots \texttt{6,\allowbreak 2,\allowbreak 6,\allowbreak 5,\allowbreak 4,\allowbreak 5,\allowbreak 4,\allowbreak 6,\allowbreak 4,\allowbreak 6,\allowbreak 5,\allowbreak 7,\allowbreak 6,\allowbreak 7,\allowbreak 6,\allowbreak 7,\allowbreak 6,\allowbreak 7,\allowbreak 6,\allowbreak 7}; order \texttt{9,\allowbreak 2,\allowbreak 7,\allowbreak 19,\allowbreak 5,\allowbreak 17,\allowbreak 11,\allowbreak 15,\allowbreak 6,\allowbreak 13,\allowbreak 4,\allowbreak 20,\allowbreak 10,\allowbreak 18,\allowbreak 8,\allowbreak 16,\allowbreak 3,\allowbreak 14,\allowbreak 1,\allowbreak 12}\par
\noindent \textbf{249}: roots \texttt{3,\allowbreak 2,\allowbreak 6,\allowbreak 5,\allowbreak 4,\allowbreak 6,\allowbreak 4,\allowbreak 6,\allowbreak 5,\allowbreak 7,\allowbreak 6,\allowbreak 7,\allowbreak 6,\allowbreak 7,\allowbreak 6,\allowbreak 7}; order \texttt{7,\allowbreak 2,\allowbreak 5,\allowbreak 15,\allowbreak 1,\allowbreak 13,\allowbreak 9,\allowbreak 11,\allowbreak 4,\allowbreak 16,\allowbreak 8,\allowbreak 14,\allowbreak 6,\allowbreak 12,\allowbreak 3,\allowbreak 10}\par
\noindent \textbf{250}: roots \texttt{2,\allowbreak 6,\allowbreak 5,\allowbreak 4,\allowbreak 6,\allowbreak 5,\allowbreak 7,\allowbreak 6,\allowbreak 7,\allowbreak 6,\allowbreak 7}; order \texttt{1,\allowbreak 4,\allowbreak 10,\allowbreak 6,\allowbreak 8,\allowbreak 3,\allowbreak 11,\allowbreak 5,\allowbreak 9,\allowbreak 2,\allowbreak 7}\par
\noindent \textbf{251}: roots \texttt{2,\allowbreak 3,\allowbreak 6,\allowbreak 5,\allowbreak 6,\allowbreak 5,\allowbreak 7,\allowbreak 6,\allowbreak 7,\allowbreak 6,\allowbreak 7}; order \texttt{1,\allowbreak 2,\allowbreak 10,\allowbreak 6,\allowbreak 8,\allowbreak 4,\allowbreak 11,\allowbreak 5,\allowbreak 9,\allowbreak 3,\allowbreak 7}\par
\noindent \textbf{252}: roots \texttt{2,\allowbreak 6,\allowbreak 4,\allowbreak 6,\allowbreak 5,\allowbreak 6,\allowbreak 5,\allowbreak 6,\allowbreak 5,\allowbreak 7,\allowbreak 6,\allowbreak 7,\allowbreak 6,\allowbreak 7,\allowbreak 6,\allowbreak 7,\allowbreak 6,\allowbreak 7}; order \texttt{1,\allowbreak 17,\allowbreak 9,\allowbreak 15,\allowbreak 3,\allowbreak 13,\allowbreak 7,\allowbreak 11,\allowbreak 5,\allowbreak 18,\allowbreak 8,\allowbreak 16,\allowbreak 6,\allowbreak 14,\allowbreak 4,\allowbreak 12,\allowbreak 2,\allowbreak 10}\par
\noindent \textbf{253}: roots \texttt{4,\allowbreak 3,\allowbreak 6,\allowbreak 4,\allowbreak 6,\allowbreak 4,\allowbreak 5,\allowbreak 6,\allowbreak 5,\allowbreak 6,\allowbreak 5,\allowbreak 7,\allowbreak 6,\allowbreak 7,\allowbreak 6,\allowbreak 7,\allowbreak 6,\allowbreak 7,\allowbreak 6,\allowbreak 7}; order \texttt{19,\allowbreak 2,\allowbreak 17,\allowbreak 6,\allowbreak 11,\allowbreak 4,\allowbreak 9,\allowbreak 15,\allowbreak 1,\allowbreak 13,\allowbreak 7,\allowbreak 20,\allowbreak 10,\allowbreak 18,\allowbreak 8,\allowbreak 16,\allowbreak 5,\allowbreak 14,\allowbreak 3,\allowbreak 12}\par
\noindent \textbf{339}: roots \texttt{3,\allowbreak 6,\allowbreak 4,\allowbreak 3,\allowbreak 6,\allowbreak 8,\allowbreak 7,\allowbreak 8,\allowbreak 7}; order \texttt{4,\allowbreak 9,\allowbreak 1,\allowbreak 7,\allowbreak 3,\allowbreak 5,\allowbreak 8,\allowbreak 2,\allowbreak 6}\par
\noindent \textbf{341}: roots \texttt{3,\allowbreak 7,\allowbreak 6,\allowbreak 4,\allowbreak 6,\allowbreak 7,\allowbreak 8,\allowbreak 7}; order \texttt{4,\allowbreak 8,\allowbreak 1,\allowbreak 5,\allowbreak 6,\allowbreak 3,\allowbreak 7,\allowbreak 2}\par
\noindent \textbf{343}: roots \texttt{5,\allowbreak 4,\allowbreak 6,\allowbreak 4,\allowbreak 6,\allowbreak 5,\allowbreak 6,\allowbreak 8,\allowbreak 7,\allowbreak 8}; order \texttt{9,\allowbreak 4,\allowbreak 6,\allowbreak 2,\allowbreak 7,\allowbreak 1,\allowbreak 5,\allowbreak 10,\allowbreak 3,\allowbreak 8}\par
\noindent \textbf{344}: roots \texttt{7,\allowbreak 6,\allowbreak 5,\allowbreak 6,\allowbreak 7,\allowbreak 8,\allowbreak 7,\allowbreak 6,\allowbreak 7,\allowbreak 8,\allowbreak 7,\allowbreak 8}; order \texttt{11,\allowbreak 8,\allowbreak 3,\allowbreak 4,\allowbreak 9,\allowbreak 12,\allowbreak 7,\allowbreak 2,\allowbreak 5,\allowbreak 10,\allowbreak 1,\allowbreak 6}\par
\noindent \textbf{346}: roots \texttt{5,\allowbreak 7,\allowbreak 4,\allowbreak 6,\allowbreak 5,\allowbreak 6,\allowbreak 8,\allowbreak 7,\allowbreak 6,\allowbreak 7,\allowbreak 8}; order \texttt{3,\allowbreak 9,\allowbreak 5,\allowbreak 6,\allowbreak 1,\allowbreak 10,\allowbreak 4,\allowbreak 8,\allowbreak 11,\allowbreak 2,\allowbreak 7}\par
\noindent \textbf{347}: roots \texttt{3,\allowbreak 4,\allowbreak 6,\allowbreak 4,\allowbreak 6,\allowbreak 7,\allowbreak 8,\allowbreak 7,\allowbreak 8}; order \texttt{8,\allowbreak 4,\allowbreak 1,\allowbreak 2,\allowbreak 5,\allowbreak 9,\allowbreak 3,\allowbreak 6,\allowbreak 7}\par
\noindent \textbf{348}: roots \texttt{6,\allowbreak 5,\allowbreak 6,\allowbreak 5,\allowbreak 6,\allowbreak 7,\allowbreak 6,\allowbreak 7,\allowbreak 8,\allowbreak 7,\allowbreak 8}; order \texttt{10,\allowbreak 4,\allowbreak 7,\allowbreak 2,\allowbreak 5,\allowbreak 11,\allowbreak 3,\allowbreak 8,\allowbreak 1,\allowbreak 6,\allowbreak 9}\par
\noindent \textbf{349}: roots \texttt{7,\allowbreak 4,\allowbreak 6,\allowbreak 5,\allowbreak 6,\allowbreak 8,\allowbreak 6,\allowbreak 7,\allowbreak 6,\allowbreak 8}; order \texttt{9,\allowbreak 2,\allowbreak 7,\allowbreak 4,\allowbreak 5,\allowbreak 8,\allowbreak 3,\allowbreak 10,\allowbreak 1,\allowbreak 6}\par
\noindent \textbf{350}: roots \texttt{7,\allowbreak 4,\allowbreak 3,\allowbreak 4,\allowbreak 6,\allowbreak 8,\allowbreak 6,\allowbreak 8}; order \texttt{3,\allowbreak 4,\allowbreak 7,\allowbreak 2,\allowbreak 5,\allowbreak 8,\allowbreak 1,\allowbreak 6}\par
\noindent \textbf{353}: roots \texttt{3,\allowbreak 2,\allowbreak 6,\allowbreak 3,\allowbreak 5,\allowbreak 4,\allowbreak 6,\allowbreak 7,\allowbreak 6,\allowbreak 8,\allowbreak 7,\allowbreak 8}; order \texttt{2,\allowbreak 4,\allowbreak 6,\allowbreak 1,\allowbreak 11,\allowbreak 9,\allowbreak 5,\allowbreak 7,\allowbreak 12,\allowbreak 3,\allowbreak 8,\allowbreak 10}\par
\noindent \textbf{354}: roots \texttt{7,\allowbreak 4,\allowbreak 3,\allowbreak 4,\allowbreak 6,\allowbreak 8,\allowbreak 7,\allowbreak 6,\allowbreak 8,\allowbreak 7,\allowbreak 8}; order \texttt{10,\allowbreak 3,\allowbreak 4,\allowbreak 8,\allowbreak 2,\allowbreak 5,\allowbreak 11,\allowbreak 7,\allowbreak 9,\allowbreak 1,\allowbreak 6}\par
\noindent \textbf{358}: roots \texttt{3,\allowbreak 2,\allowbreak 7,\allowbreak 4,\allowbreak 5,\allowbreak 6,\allowbreak 8,\allowbreak 6,\allowbreak 7,\allowbreak 8}; order \texttt{2,\allowbreak 4,\allowbreak 1,\allowbreak 8,\allowbreak 5,\allowbreak 6,\allowbreak 9,\allowbreak 10,\allowbreak 3,\allowbreak 7}\par
\noindent \textbf{363}: roots \texttt{7,\allowbreak 6,\allowbreak 4,\allowbreak 6,\allowbreak 5,\allowbreak 6,\allowbreak 7,\allowbreak 8,\allowbreak 7,\allowbreak 8,\allowbreak 7,\allowbreak 8}; order \texttt{11,\allowbreak 3,\allowbreak 5,\allowbreak 6,\allowbreak 12,\allowbreak 9,\allowbreak 4,\allowbreak 7,\allowbreak 2,\allowbreak 10,\allowbreak 1,\allowbreak 8}\par

\end{multicols}

\end{document}
